% ============================================================
%  F. C. Sibbern, On Love, or Love between Man and Woman
%  (Om Elskov eller Kjærlighed imellem Mand og Qvinde)
%  Kjøbenhavn: Paa eget Forlag (Trykt hos J. H. Schultz), 1858.
%  "Tredie uforandrede Udgave" — text identical to the 1st ed. (1819).
%  TRANSLATION (English)
%  Translated from the Danish transcription (transcription.tex).
%
%  TRANSLATION CONVENTIONS (see RESUME-NOTES.md for fuller notes):
%   * The book turns on a distinction English blurs: Elskov = love
%     between the sexes (eros, romantic-erotic love); Kjærlighed =
%     love in the wider sense (affection, charity, caritas). Both are
%     rendered "love"; where the two are contrasted or the choice of
%     word matters, the Danish is given in brackets, e.g. love [Elskov].
%   * 1819/1858 register kept but made readable; Sibbern's long periods
%     are preserved, only lightly broken where English demands.
%   * German (Goethe), Greek, and Latin left in the original with an
%     English gloss in a translator's footnote.
%   * Section breaks (the printed centered rules) reproduced as rules.
%   * Paragraphing follows the transcription.
% ============================================================
\documentclass[12pt,a4paper]{book}
\usepackage[utf8]{inputenc}
\usepackage[T1]{fontenc}
\usepackage{libertinus}
\usepackage{libertinust1math}
\usepackage{textalpha}   % Greek words quoted from the original
\usepackage[english]{babel}
\usepackage[protrusion=true,expansion=true]{microtype}
\usepackage{setspace}
\usepackage[margin=1.2in]{geometry}
\usepackage{fancyhdr}
\setlength{\headheight}{14.5pt}
\addtolength{\topmargin}{-2.5pt}
\usepackage{hyperref}

\hypersetup{
  pdftitle={On Love, or Love between Man and Woman},
  pdfauthor={F. C. Sibbern},
  hidelinks
}

\pagestyle{fancy}
\fancyhf{}
\fancyhead[LE,RO]{\thepage}
\fancyhead[RE]{\textit{On Love}}
\fancyhead[LO]{\textit{\leftmark}}

\setstretch{1.3}

\title{\textbf{On}\\[6pt]
  \Huge\textbf{Love}\\[6pt]
  \large or\\[6pt]
  \Large\textbf{Love between Man and Woman}\\[10pt]
  \normalsize\textit{Om Elskov eller Kjærlighed imellem Mand og Qvinde}\\[10pt]
  \normalsize by\\[6pt]
  \large Frederik Christian Sibbern.\\[16pt]
  \normalsize Third unaltered edition.}
\author{\small Translated from the Danish}
\date{Copenhagen 1858.\\
  Published by the author.\\
  Printed by J.\ H.\ Schultz.}

\begin{document}
\maketitle
\thispagestyle{empty}
\cleardoublepage

%% ------------------------------------------------------------
%%  TRANSLATOR'S NOTE
%% ------------------------------------------------------------
\thispagestyle{empty}
\chapter*{Translator's note}

Sibbern's title, and his whole argument, rest on a distinction that
English cannot carry in a single word. \textit{Elskov} is love between
the sexes---the ardour the Greeks called ἔρως; \textit{Kjærlighed}
is love in the broader sense, the affection and charity that need not
be erotic at all. The title names the first and then, with
\textit{eller Kjærlighed imellem Mand og Qvinde} (``or Love between Man
and Woman''), quietly widens it toward the second. Throughout, both
words are translated ``love''; where Sibbern plays the two against each
other, or where his choice of the one rather than the other carries the
thought, the Danish word is supplied in brackets---love [Elskov], love
[Kjærlighed].

The 1858 printing is a ``third unaltered edition'': its text is that of
the first edition of 1819. Passages of German, Greek, and Latin are left
as Sibbern gives them, with a plain English rendering in a footnote.

%% ============================================================
%%  PREFACE TO THE SECOND EDITION  (Fortale til anden Udgave)
%% ============================================================
\chapter*{Preface to the Second Edition.}
\markboth{Preface to the Second Edition}{Preface to the Second Edition}

For several years the first printing of the present book has been out of
print, and for some years there has been talk of my seeing to a new
edition of it. But each time I took the book out to that end, in order to
go through it, I found it repugnant to read, and I could not bring myself
to do so. At last a young man of learning offered to go through it for
me, to correct whatever might need correcting in the way of expression
and orthography---chiefly so that all those words might come to stand
with capital letters which I now, for the sake of greater clarity and
definiteness, write so, but formerly did not; for the rest, there was no
thought of making any change in style and presentation, on which I would
not enter.\footnote{That in many a place the word ``he'' might have had
``or she'' added to it, the lady readers will doubtless soon note.} Thus,
after the lapse of many years, I have first come to read this book again
in the proof-sheets, and I must confess that a good deal now seemed to me
not so very amiss, though a good deal was not at all to my liking, even
if it had to be left standing when the book was to appear again; for it
was to be as it was. Changes---trifles apart---I did not feel myself
capable of; they would probably also have clashed with the rest. So I
must repeat what I see I have already said at the close of the Preface to
the first edition, and now feel doubly called upon to say. A collection
of psychological observations and religious and moral intimations one
will indeed find here, and a few passages of greater worth. At some
places, especially in the third Book, I have imagined that, just as
novels have often been accompanied by copperplates which, when they came
from a Chodowiecki's hand, might be called truly valuable illustrations,
or heightenings, whereby what was set forth received, through another
art, a clarifying elucidation and effect, so likewise much in this book
might be suited to receive the illustrations of the novelist's art, or
might furnish a kind of theme for treatment in the novella. But there is
yet something that may perhaps be able to give much, especially in the
first Book, a peculiar significance. It is this: that the author drew it
directly from what he had, within himself, not merely experienced but
also lived, while in this respect he went to the fairest school. There
has sometimes been talk that he too should attempt to give depictions of
his life. Much as, in former times, he has often rejoiced at the thought
of one day being able to do so---especially in order to set up a grateful
memorial to many splendid men and women of his circle---he yet dare not
hope that he will ever attain to begin upon it. But in the present book,
as in his psychological Pathology, and perhaps in many another place in
his writings (besides in his two Gabrielises, of which nothing shall here
be said), he has given many a confession of his life; and on the whole he
may well, recasting some familiar lines of Goethe, say:

\begin{verse}
Was ich liebte, was ich lebte,\\
Was ich schaute, was ich strebte,\\
Wollte ausgesprochen seyn;\\
So vertraut' ich oft sub rosa.\\
Beichtend in ganz schlichter Prosa,\\
In der Forschung erstem Hain.\footnote{``What I loved, what I lived,
/ what I beheld, what I strove for, / wished to be uttered; / so I often
confided, \textit{sub rosa}, / confessing in quite plain prose, / in
research's first grove.''}
\end{verse}

\noindent Copenhagen, the 1st of November 1853.
\begin{flushright}
Fr.\ Chr.\ Sibbern.
\end{flushright}

%% ============================================================
%%  PREFACE TO THE FIRST EDITION  (Fortale til første Udgave)
%% ============================================================
\chapter*{Preface to the First Edition.}
\markboth{Preface to the First Edition}{Preface to the First Edition}

The immediate occasion for working out the present writing was given me
by the lectures I have for several years delivered at the University
under the name of Lectures on the Doctrine of Happiness. In these
lectures I linked together, following a certain definite train of
thought, a series of considerations directed at the enjoyment of a higher
life under the forms and relations of earthly life---without, however,
letting any strict scientific distinction between what properly belonged
to the purpose and what did not hold me back from elucidating in detail,
by some even very extensive side-considerations, several of the most
interesting among the objects that presented themselves, from every side,
so far as time and the main purpose of the lecture would allow. Hereby it
came about that some objects were treated in a far more comprehensive
manner than others taken up in the same plan and essentially belonging to
it. In a scientific work handed over through print to calmer
contemplation and fixed once and for all, I saw well enough, on the other
hand, that such a deviation, occurring at certain principal parts, from a
definite underlying main thought would not look too well, and still less
the unevenness in the treatment of the various parts that would go with
it. Those objects treated with particular fullness therefore showed
themselves more and more to me as ones that could far better be presented
in separate, self-standing elaborations. Thus the present writing has
come into being as a smaller whole for itself; and thus I hope in time to
be able to deliver separately my philosophical considerations on academic
study, which formed one part of those academic lectures, and likewise,
under God's assistance---which must thereby be doubly invoked---my
philosophical considerations on the chief doctrines of Christianity. What
may afterwards remain, as fitting material for an ever onward-progressing
development, in one constant direction, of a doctrine of human
happiness---and indeed whether such a doctrine is to come about through me
at all---will doubtless show itself to me by degrees. This provisionally
by way of answer to those of my hearers and friends who have kindly at
times urged me to publish the aforesaid lectures, which, for the rest, I
have not yet worked out in writing.

The considerations here communicated on love [Elskov]---not, precisely,
in any properly scientific form, whether as regards outer or inner
form---have of themselves divided into three principal parts. These it
seemed to me most fitting to call Books, since the designation of
sections, chapters, or the like would more suggest a systematic procedure
not here employed; though they are, to be sure, of somewhat too slight an
extent to deserve that name. A fourth Book, already for a great part
worked out, on the natural drive that co-operates in love, I have for
various reasons held back. The greatest part of it did not, moreover,
belong under the book's title, whereas it will have its proper place in
the psychological Pathology which, as the second part of my Psychology, I
hope one day to deliver. The considerations on unhappy love in the third
Book have turned out shorter than I had at first thought. On the whole I
have often felt, with them, how much better they might be
elucidated---since individual circumstances come so much into
consideration in them---by a presentation of quite another kind. An
attempt at such a one I may perhaps lay before the public one
day.\footnote{The allusion is to the ``Gabrielis's Posthumous Letters,''
then already for the most part written, which appeared for the first time
in 1826.}

For the rest---although I must confess that in my day I set down, and
several times went through, most of the principal parts of this writing
not without true delight and joy---I have since, during the printing, so
well seen and felt how far it is from being what, according to the idea of
such a presentation, it ought to be, that I cannot but beg the reader not
to come to its reading with great demands.

\noindent Copenhagen, the 12th of December 1819.
\begin{flushright}
The Author.
\end{flushright}

%% ------------------------------------------------------------
%%  POSTSCRIPT TO THE THIRD EDITION  (Efterskrivt ved tredie Udgave)
%% ------------------------------------------------------------
\bigskip
\noindent\textit{Postscript to the third edition,} the 26th of April
1858: Some small changes, made during the proof-correction, have not been
able to keep me from calling the present edition an unaltered one. It
follows the preceding one not merely page for page but line for
line---that preceding edition having been of only 350 copies, and
therefore having taken four and a half years to sell out. The present
printing is twice as large, and the price has been reduced to 48 sk.\ from
80 sk.---For the rest I will only remark that from a book such as this
everyone should take for himself whatever may be of use and furtherance to
him or her, of refreshment and gladdening, and let the rest lie.

%% ============================================================
%%  FIRST BOOK
%% ============================================================
\chapter*{First Book.}
\markboth{First Book}{First Book}

With inexpressible feelings the awakening of love [Elskov] fills every
inwardly loving soul, and in a new, glorious and rich life the existence
reborn in happy love unfolds. But only too often has Nature bestowed upon
a human being the richest fullness of this earthly blessing, by a good
fortune she has appointed for each, without the fortunate one letting
what was allotted to him so enter into himself and become his own that it
might become for him a genuine treasure in life and a perpetually flowing
well of gladness. Only too often does love's fair blossoming-time pass by
like a dream, and fails to become what it might become: the entrance to an
unchangeable existence. And yet no earthly way to the true life is fairer
than this, and no earthly light shines more clearly for genuine striving
than love's powerfully awakening and inwardly refreshing light. With
penetrating voice Nature, in this feeling, proclaims the inner meaning of
life, and by it calls us out of the widely spread manifoldness of the
visible into the deeper sanctuary of the invisible. For him who would
press into the inner truth and essentiality of life and Nature, the
fairest consecration is here prepared. But if it is to become for him such
a consecration, if from it a lasting joy in existence is to develop in
him, then he must grasp it in its inner essence and appropriate it from
the heart through a deep cognition [Erkjendelse]. Not that he, tracking
and sundering the shifting series of feelings---as they perhaps seem to
him to present themselves in a peculiar and remarkable way---should,
brooding over himself, fall out of the full sensation of life, and now
sunder apart for himself the life he formerly lived whole and as of one
piece. But grasp and cognize he must what has been bestowed on him,
precisely in order to be able to live and enjoy it inwardly and to the
full, as he feels its inner depth stirring in every shifting moment, and
lets this depth come to full clarity and to a penetrating diffusion
through his whole existence---which does not happen without a cognition.
By this, love must be raised up to show its power as a principle that
genuinely animates and forms life; and likewise it must thereby find its
firm ground, wherein it can take root, and its higher home, wherein it can
be transfigured, in the sphere of the Eternal. Without finding this
foundation and attaining this transfiguration, love will blossom itself
out, as it blossomed forth. The whole spiritual life of the human being
has its root in love [Kjærlighed] toward the Eternal One and in faith in
his infinite being; and this faith and this love must in every respect
constitute the heart---the inner warming and moving centre of the feelings
of life---in all the members and parts of the spiritual life. --- Therefore,
both in order to attain this cognition, and also in general to perceive
that voice of Nature which speaks to us in this fairest among the powers
of life that work toward a merely earthly blessing, we will enter into a
reflection upon love's nature and content. We only follow Nature's hint
when, in love---as it lives in happy, pure and strong souls, and unfolds
in their shared life---we seek an image of the true life. In this its
fairest and most fortunate shape we will, then, consider it first.

The two outermost, most mutually opposed points in the whole fullness of
love---the two fixed points, moreover, about which, as about two poles,
the whole consideration of it must turn---are: the wholly earthly in it
on the one side, and, on the other, that in it which is wholly turned
toward the Eternal. The former is the natural drive by which love's
longings are called forth. The latter is the joy in the beautiful and the
divine, which is likewise in motion wherever love comes forth. In
accordance with the first, the individual seeks its other half; in and
under the second, the nearness of its God. In genuine love, the one whose
striving is a happy one finds both in one: in his other half, the glorious
image of his God. What here, as inner natural prompting, seems to lead
only to an earthly satisfaction cannot, in the human being---because he is
a human being---find a human fulfilment without bringing him unspeakably
near to his eternal source and opening for him the temple of the Holiest
of Holies.

In a youthful mind, open to life's manifold impressions, love springs
almost always from a delight and joy of contemplation, and seeks in this
its first nourishment. Whether its beginning, in regard to an individual,
is a sudden movement of the mind and a rapture that shows how strongly the
heart is struck, or whether it advances slowly out of friendly intercourse
and fair companionship, in the course of which little by little two souls
unfold to each other: what the one beholds, and with delight ever anew
rejoices to behold, in the other is in both cases that which awakens
devotion and draws the one heart to the other. It is, to be sure, the
drive that lies in the background---it is the deeply stirring nature that
attunes the soul to inclination and to love's softening longing, it is a
need and a straitness rising from the individual's inmost depth, that
awakens and opens the eye for the beneficent element in another
individual, and makes it possible for the one to be so strongly attracted
by the other. But however much the stirring nature may set the human being
in restless motion, he yet cannot---as one in whom the sense for something
Higher is awake and everywhere ready to come forward---unless he wholly
denies his nature, attach himself to the object of his desire without
finding a pleasure in it, and that a spiritual pleasure. Thereby now, as
the object encountered is held up against an inner demand, the higher idea
of beauty is awakened to life. The longing becomes now a longing for the
beautiful and the comely, and a desire to be inwardly united with a
beautiful and captivating human being. And thus must then---to borrow an
image from the Nordic myth of Freia's carriage drawn by cats, so
beautifully interpreted by Oehlenschläger in his Saga of Vaulundur---the
appetites, which elsewhere everywhere in Nature are so often wild and
unruly, now in the soundly and purely feeling human being draw forth the
divinity of the highest love, while yet they are governed by it. Nature's
strong movements call forth a deep need and longing of the heart, a need
for the soul's fulfilment and for soul-joy. This now becomes predominant,
and what in Nature was, in the organic respect, the first ceases to be so
any longer, and steps back. Love's longing becomes captive to a definite
person only then, and rejoices at having found its object, when the
awakened sense is seized and filled by the inwardly attractive and
glorious thing that opens itself to the soul's eye in the contemplation of
the other being. Only then arises love's living feeling, which fulfils
sense, mind and heart. It shows its first power in the inward joy, the
inextinguishable delight, with which the one individual dwells and
rests in---nay, loses itself in---the contemplation of the other, and from
it, as from an inexhaustible spring, draws a perpetually fresh
refreshment, so that the heart is moved in a living gladness each time the
beloved form shows itself anew. And even if it is sometimes not the
delight of contemplation with which love begins---if it is the shelter
that one heart has found with another, and the trust and confidingness
springing therefrom, that have drawn forth the devotion and the
love---it is yet natural that this sympathy now gradually draws forth
contemplation, and the joy in it, and only then does it come to a proper
love. When a heart's inner worth is felt and cognized, it will also soon
be beheld with true delight in all its expressions; and then, when love's
drive and longing are thereby set in motion, soon the whole outer form in
which such a being reveals itself will also be beheld with inward
pleasure, and it is then fulfilled, in the most refreshing manner, that
love makes one see.\footnote{Sibbern inverts the proverb
\textit{Kjærlighed gjør blind}, ``love makes blind.''}

But however much this joy in the beloved attaches itself to the most
idiosyncratic and personal, and precisely thereby shows itself as a
proper love directed at a definite person and at no other than this one,
it must nevertheless have its ground in a cognition of something that, in
and for itself and without regard to all personal inclination, is good and
excellent in the beloved individual. Everyone who loves from the heart
loves because he beholds, or believes he beholds, something in itself
lovable in the other; but an individual is lovable only through the
Higher, the Divine, which it presents in itself, or whose stamp it bears
and of whose living power it is filled. Precisely because an individual
seems to us in itself to be fair, comely, lovely, of a pure or deep
nature, or in general of higher stamp, do we feel ourselves drawn toward
it. Love proceeds from an inward pleasure, which presupposes a hearty
approval. But this approval is possible only when we, referring what we
behold to an idea living within ourselves, feel and find true that it
deeply coincides with this. Indeed, only by this can we cognize the
excellent outside us at all, and grasp its worth. Did the divine fullness
not stir in our inmost depth, we should not be able to find it, and to
rejoice at finding it, outside us. Thus, then, that which even by mere
contemplation draws two beings to each other is something higher than
both. Just as, where two are gathered in Christ's name, Christ is in the
midst of them, so, where two are united by a genuine love, it is the
Divine that unites them, which is present in both, and which the one
beholds in the other when he looks into the other's soul as into the
heaven opened for him. Where it is not this mighty presence that holds
both united by a true joy in each other, there love is not built upon a
lasting ground either. A delusion has drawn the two together; for the one
has then beheld in the other what was not in that one. And only the
in-itself Good, as the genuine and abiding, can awaken and preserve the
deep trust without which love is nothing. The mere thought that the other
being perhaps is not what it seems---a pure vessel for the higher
Spirit---gnaws at once at the bond between the two, and the certainty of
it is like an evil demon standing between them. And even should the outer
allurements keep their power, yet the proper love is broken asunder, and
the thought that what still holds us captive is in its inner essence
spurious and untrue stifles even the delight of outer contemplation. Where
two believe themselves united in love, and are not so by the one power
that in truth makes them one---a higher life's spirit animating
both---the shared life resembles, for each of them, a life with a hollow
revenant that seems to be a full human being. The God-filled one loves, in
all that fills him with deep love, the infinite nearness of his God; and
whoever rejoices deeply and inwardly in anything glorious at all, in him
there stirs a feeling, rising from the soul's depth, of something
invisible and eternal, of whose heavenly glory he is, as it were,
reminded, as the Same speaks to him out of the outer object. Therefore he
feels himself refreshed to his inmost by what he beholds, as by an eternal
beauty's infinite presence in his soul. This it is that we prize, and to
which we attach ourselves, in the beloved; and the joy in the beloved is
then already a true joy in the Lord.

Meanwhile, the Divine living in both lovers, and its mutual cognition, are
not alone enough to make captive the drive and longing that are in motion,
and thereby to call forth and ground love. Inward esteem and
admiration---indeed a love of another kind than love [Elskov]---can
thereby be awakened and confirmed. But only when a joy in the personal
individuality, the peculiar idiosyncrasy, is added---when we are not
merely enraptured and inwardly gladdened by the excellent and glorious in
the other person, but also by this person itself and its personality
transfigured by the Divine---can a real love develop out of the delight of
contemplation, and love's longing find its fullness. Whatever it may be in
the beloved that can have been the first attracting and awakening thing,
and moreover the first incitement to love---be it the beauty of form, the
expression of the features or the grace of the movements, the energy of
uprightness and the sure eye upon life, or the mind's fresh liveliness and
the spirited play of the imagination, the deeper sense for the inner being
of existence, or a quiet peace reposing and resting in itself, or whatever
else it may have been that first struck and attracted the heart, and that
perhaps also shows itself, to the more distant observer, as the most
salient thing in the beloved individuality---yet that which now, in the
fullness of the feelings, occupies the lover and has his mind and heart in
its possession is not this or that single and particular thing alone, but
is the whole individuality itself, as it utters itself, in a whole and
full expression, out of the beloved's whole existence and all its
expressions. The first impression may be the effect of these or those
allurements which have carried away the youthful, open sense and taken the
mind up in rapture by their kindling or enlivening power. Some single
thing is often that which first, as incitement, draws the eye to itself
and awakens the desire and the longing in the soul that longs for love's
sweet delight. Thereby the mind is carried away and gives itself over to
the joy of contemplation; but the single allurement now becomes that which
spurs us to press into the inner wholeness and deeper essence of the soul
that meets us---it becomes that on which we, as it were, lean, and which
leads us, in the yearning delight, to see through the beloved soul
penetratingly and to pursue it in all its expressions. But now it comes to
the question whether its inner wholeness so answers to our longing that it
can still and fill it. Only then, when not merely the allurement that
carries us away in the moment, but also the fullness quietly living and
working in its depth in the beloved being, exercises upon us a lasting
influence, does it come to a really genuine love. The beloved's ever
self-same inner essence---which, with its abiding fundamental character,
constitutes the individuality's inner fund or well, from which the outer
manifoldness and movement of life proceeds---this inner essence it is that
we must have beheld and cognized, and by which we must feel ourselves
inwardly attracted, before we can be said in truth to love. For which it
is required that the first vehement movement subside, and that the mind
come to such composure and clarity that the lover can collect himself and
tell himself whether it is only the fresh, luxuriant expression of life,
or whether it is the spiritual fundamental shape and inner fullness in the
beloved itself, to which he feels himself drawn with so great a power, and
which with ever-fresh delight he longs to behold and behold again.
Genuine love can, in general, not develop out of passing impressions, but
only out of an enduring life, in which two, by living much with each
other, come to open their whole inner idiosyncrasy to each other. One must
live oneself into love [Kjærlighed].

It is thus the soul itself that we love, when love is genuine, in another
individual---the inner animating spiritual whole from which the
individual's whole life and idiosyncrasy proceeds. Yet the soul's inner
glory exerts its strongest influence in the manner in which it utters
itself in the feelings, and in the emotions and moods of mind that spring
from these. The soul that utters itself in the whole manifoldness of these
constitutes the proper mind, heart, heart's-disposition, temper [Hu,
Hjerte, Hjertelag, Gemyt]. By this we designate the inner fundamental
movement in the soul, from which the warmth of life proceeds---the inner
deep well for the streams of the feelings that spread out to every side.
Depth and clarity of cognition, and firmness and energy of will---the
greatest advantages in respect of talent and strength of character---do
not yet produce genuine, inward love; whereas a good heart, a deeply
feeling and pure temper, with which a good, though not precisely a
vigorous and strong, will always goes, is often to the lover full
compensation for lacking spiritual power and natural gifts.

When love attaches itself to the temper, and in general to the
individuality, to the soul itself in the other being, it attaches itself
to the abiding and imperishable, and an inexhaustible pleasure can then
proceed from it. The soul is the eternally young in the human being;
indeed, self-same, it can still, out of the aged outer form, shine toward
the lover in its unchanged fundamental shape. But whole we do not behold
it in any single moment of the other's life. Only when this life, in its
manifold expressions, spreads itself out before the gaze, does the soul
shine forth from it in its inner wholeness. One might compare this inner
spiritual individuality with a being who, wrapped in a garment, so
betrayed its shape, through innumerable movements, to the observer eager
to see, that at last, without stepping forth from its cover, it stood
whole before his imagination. The individuality utters itself in moments,
but its total expression is only the whole individual itself. Hence there
is here, for the lover's eyes, the perpetual novelty in the beloved's
manifold movements and expressions, when the life in this one is, by love,
awakened and constantly fresh. The beloved fundamental shape he rejoices
to behold, perpetually the same and yet perpetually new, repeated in the
countless movements of life, and finds in this a delight as inexhaustible
as life and individuality themselves.

The beloved's lively expressions are sheer allurements, since they all
reach the fundamental image lying in the lover, which is so dear to his
heart, and thereby awaken his delight and joy anew. The deeper and clearer
the love, the more inward the pleasure, the more the dear image unfolds in
our inmost. In lovers of deep feeling and deep mind there often develops,
out of their joy, a true study---a word that must be taken more literally
than at first glance it may seem. For is not all study the more genuine
and the more beneficent, the more wholly we sink ourselves into the
contemplation of the object, so that we thereby let it, in its full
wholeness, enter into us and stamp itself in us---whereby also, at the same
time, our own spirit's fullness first develops from within?

The lover is, in his love, a seer, who is able, through the outer cover,
to behold into the heart in the other being. A fortunate seer's-glance
belongs to the awakening of genuine love. But every youthfully moved soul
easily has such inspired moments. Yet there may sometimes really be need of
good fortune for another being's inner [self] to open itself to us, and
for a ray from the heaven that lies therein prepared for us to strike our
eye in an hour when our mind is free and open enough to receive a pure
impression with purity. Not seldom did two live long side by side before a
fortunate moment brought them to open their inmost being to each other and
mutually to cognize each other's worth, so that now both felt themselves
as touched by something Divine. Sometimes, too, the heavenly content lay
locked and hidden in a human heart, like a diamond in the rock, until a
fortunate finder cognized its worth and called it forth into the day, so
that now, in love's sun, it might radiate---that is, radiate back---an
inexpressible glory of light. It also often happens that an individual who
seemed to others, and continues to seem to them, of little significance is
loved with a heartiness and a pleasure and an abiding contentment and
delight that prove the inner genuineness and depth of the love too well
for there not to lie, at the ground of it, something particularly
genuine---which, however, only the lover, who has fortunately grasped and
beheld the deeper fundamental shape in his beloved, has been able to
cognize and appropriate so heartily. To be sure, love, although by its
inner nature and essence it makes seeing, can yet sometimes bring an
individual to be ensnared in error and blindness---just as a philosophical
pondering and striving not seldom led to fanaticism or confusion of ideas,
although philosophy and philosophical pondering is nonetheless precisely
something that opens the spirit's higher eye and makes genuinely seeing.
But by no means is that love always called forth by a delusion which seems
so to others because they do not see in the beloved anything so lovable
that it can seem to them able to yield so much delight and fulfilment. A
human soul often conceals much more in its inner depth than what lets
itself be shown to others---than to any but the one whom love makes a seer.
The Divine does indeed really live in every good human being, and gives
the human being essentiality, even if overshadowed or pressed back on
account of earthly defects, which yet easily vanish for him who cognizes
and beholds that Divine in its inner purity.

Nothing is more powerful to call forth love than that beauty of the body
which is the full expression---purely and clearly stamped in the sensuous
form---of a beautiful and heavenly inner [self]; for beauty fills our whole
human sense, the higher together with the lower. The whole human being
feels himself enlivened and refreshed by it. Yet in the outer form that
works most in which the soul utters itself most immediately: the beauty of
the eye, of the features, and of the movements expressed in the whole
individuality---this last being the grace spread over the whole individual.
But beauty is a good fortune only seldom granted to a lovable individual.
It is only seldom that, in the press of life, a being blossoming forth with
youthful freshness is granted room and freedom enough for a calm and
all-sidedly penetrating unfolding of its individual fullness, whether in
the spiritual or the bodily respect---not to mention that this is so often
already one-sidedly disposed, indeed downright suppressed and cowed, in the
first germ, or in the earliest growth, before the individual has yet come
to a right consciousness of itself.

Yet for him who with inward delight appropriates his love and his beloved,
she is always beautiful. And it is quite natural that the beloved form
exercises the power of beauty over him. All beauty rests on this: that
something which springs from life's higher source and bears its stamp in
itself steps fully forth in an outer, sensuous or temporal form, filling
and penetrating it, so that this bears the full expression of the former.
But what the lover beholds as lovable, and moreover as heavenly, in the
beloved's inner essence---this same thing he sees, precisely, with inward
joy expressed and uttered in the whole beloved individuality and every one
of its expressions. And thus the beloved will also, in respect of her
outer [aspect], make the impression of beauty upon the lover, and his joy
and pleasure will be of the same kind as that which beauty calls forth. Is
it not, too, a ray from the true heaven that falls into his soul through
his beloved, and does this not awaken the infinite feeling of the eternal
Beauty's heavenly glory? But whatever is able to awaken this feeling in us
thereby itself stands transfigured for us; and even where the individual
completeness and fullness is wanting, we prize as beautiful that being at
whose contemplation a remembrance of the eternal Beauty perpetually lives
up again in us. Hence there develops, out of love, a pleasure and a delight
not merely in the in-itself excellent and essential, but also in the most
trifling among the beloved's expressions and movements, because the lover
therein perpetually meets the beloved individuality, and takes delight in
seeing it stamped down to the very surface of the form. Yet this pure joy
is felt most fully in the total impression, in which everything single and
manifold melts together into a single calm and whole intuition, wherein
the beloved's whole being fills our soul.

Unspeakably comforting and enlivening is this pure contemplation, this
calm, undivided absorption in the beholding of a lovable being, and the
mild, clear influence with which the image of this being then enters into
the beholder's soul. Awakening in our inmost the kindred spiritual fullness
in the pure feeling of the Eternal's inward loftiness and glory, it calls
forth in us a higher life's free, purely spiritual sensations. As in the
sun's light and warmth organic life unfolds in Nature, so a new spiritual
life unfolds in love's warming light. The individuality cognized and
beheld with inward delight brings our own spirit's inner fundamental shape
and essence to unfolding. The soul's inner eye, as if touched by a
divinity, now opens to behold, living, the essentiality veiled under the
earthly cover in life and Nature; and, as love's far-extended realm, the
whole of Nature lies spread out before the spiritual gaze.

Yet, in love's first awakening and working, the new world that is to come
forth from the soul's inmost still reposes only in one great total
impression. The first moments may indeed work as a spur; the enamoured may
come into much outward movement, show themselves elated, ready for jest and
merriment; they are ready for anything that can bring them near the
beloved's person, they are alert on all sides. But it easily happens soon
that the loving soul, feeling the strangely deep quality in what moves it,
withdraws into itself. It begins to feel life's high earnestness; it is as
if something Divine (θεῖόν τι) had touched it; everything
personal, all desire falls silent; the shifting play of the feelings melts
away into an unspeakably deep and quiet longing for the absent, and an
equally deep and quiet joy in the present, beloved. Still only in a wholly
enlivened existence, not in a definite desire or overture, does the soul
venture to open itself to the beloved. Even when it feels unspeakably sure
of a returned love, it is yet as though it feared, by a definite outburst,
to disturb or disquiet everything that still stands so heavenly-fair before
it. Yet soon the increasing full joy in the other being's existence
reaches, awakening and spurring, into the whole personality, and calls
forth the most vehement desire and striving, with courage and
self-confidence. Thereby, too, the yearning, burning element can gradually
enter into love. But how fair stands, in and for memory, love's first mild
awakening, its calm dwelling in its inner delight, and the quiet, deep
longing and joy in which a presageful foretaste of a great, still-slumbering
life stood in the loving one's inmost, and in which a single word---the word
that uttered itself in the beloved's whole existence---surpassed for us all
the manifold words in which, later on, yearning, we took delight in hearing
our happiness uttered.

In an infinite devotion love breaks forth. With mighty power it tears the
human being out of the narrow circle of mere I-hood. However much the
individual may formerly have regarded itself as the nearest centre of its
existence, referred everything back to itself, and let the thought of
itself accompany all its receiving and striving, it will yet soon, when
love fills it, behold its life's centre in the beloved. The thought of her,
who now seems everywhere to surround us in an eternal nearness, accompanies
all its other thoughts. The beloved now becomes the one to whom the I
refers everything back that it may live, about whom all its delight and
desire turns, on whom it hangs with all its mind and all its feelings, as
the human being ought to hang upon God. How a youthful lover feels, we all
know. It is for him as if the beloved were the divinity from which his life
goes out and to which it goes back, just as if life's whole fullness
streamed to him from this source alone. In and by the beloved being he
feels himself blessed and seems to himself to live. And it is not sheer
delusion and confusion when the feeling lives in him with this power. For
from life's true sun comes, nonetheless, the ray that, bringing a new life,
falls from the beloved into the lover's soul; and what in the beloved's
being enraptures him so unspeakably is indeed the Eternal itself, which
comes to meet him, and thereby at the same time awakens in him the deepest
remembrances of the same Eternal's glory stirring in his own being. Hence
his love has this power over him, that he would feel himself as annihilated
were he to lose its object, as if he had been thrust down from his heaven.

Yet the happy lover perceives this devotion, and this foreign influence
which now lives and cries out in him, by no means as something ensnaring,
restricting, or conditioning. What the lover feels when he casts all life's
anxiety from him under the new life that wholly fills him, and when his love
becomes for him a fair refuge, a glorious possession to which---in whatever
may befall him---he hastens trustingly, because he knows it, and in it the
fairest thing he owns on earth, sure and dependable---this is precisely the
genuine feeling of freedom, because it is the feeling of a genuine and
essential existence, loosed from all lower bonds and raising itself above
the earthly. In love there are spirit and life; but the spirit everywhere
makes free. Inwardly to hang upon a beloved being, and to give oneself
wholly over to it, and to enjoy life not in the narrow circle of one's own
I but in the other being's worth and gladness and happiness---this is to
lose life in order to regain it doubled. But thus does the lover desire to
lose himself in his beloved. All love has this power to annihilate and to
baptize away the I-hood, when one loves from a full, undivided heart. Never
is the deep poverty of the I-hood felt sooner.
Whoever loves is humble at heart. How little does the loving one, with the
feeling of his own defectiveness within him, seem to himself worthy to be
loved by an adored being, worthy of the happiness he has found, as if he
were chosen before countless others. It is natural that it is often quite
incredible to a happily loving being that a blessedness which, enrapturing
the imagination, seemed at home only in the world of poetry has now really
become its own, and that it lives in the midst of the possession of that
toward which a youthful heart stretches out its inmost longings. It is
natural that it may sometimes be for the loving one as if it were all only
a dream. And yet love is felt again as the fullest, most real
possession---with an inward dependableness, security and certainty; not
such a one as, resting on these or those deliberations, would be merely a
mediate certainty, but a certainty as immediate and total as love itself.

The deep readiness to give up everything for the beloved, and the deep
feeling of humility when we consider the unspeakable glory that, as a pure
gift of heaven, is revealed upon us in our love---this feeling corresponds
exactly to the human being's inner nature, just as it also lies precisely
at the ground of love's sweetest sensations. The isolated, self-reposing
existence is indeed our first, and is on earth still so predominant; but
the genuine existence is nonetheless that in which a universal life,
streaming down to us from a higher source, enters into us and fills us.
Herein ceases and vanishes the feeling of finitude that lies in all
feeling of personality in the being that does not have life's source and
spring within itself. The individual then feels as if it were vanishing
away, or as it were dying away and being dissolved into the universal
eternal life, and therefore feels itself, in contrast to this universal
life, as little or nothing, since it feels itself as receiving so
infinitely much. Yet the personality does not herein perish, but is
penetrated and filled, and moreover raised and transfigured. Thereby,
then, this unspeakable devotion and humility is so inwardly deep and
sweet; it is the fullest feeling of the higher influence's pure presence in
us, which nothing in us resists. Precisely when, in humility's deep
feeling, self-feeling is overwhelmed and sinks away, do we feel ourselves
inwardly moved and refreshed, and feel ourselves raised, just as we feel
ourselves in our littleness and lowliness. This sweet feeling of humility
now lies, then, also in love, and belongs essentially to it; indeed,
precisely here it is so penetrating and sweet, because here the sense and
the earthly delight too are filled and refreshed. Yet it lies in the nature
of the sexes that woman, without doubt, feels this hearty readiness to
give everything up, and this humility, with a wholeness and moreover an
inwardness which a man is perhaps hardly able to grasp to the
full---wherefore, then, woman's love can also, without doubt, attain a far
higher inner strength than man's. What the fable of Tiresias utters holds
surely also in the highest spiritual respect. Precisely because woman is
the subordinate one, and is at the same time the one whose whole mind and
striving can go up into the conjugal and domestic existence---whereas man
still seeks a sphere of activity outside the house, and moreover turns his
interest to two places---precisely therefore woman must, and also can, give
herself over with an absoluteness and a wholeness that must give her
feelings so unspeakable an intensity.

The one lover is to the other as a god, so wholly does the one seem to
itself filled with new life and lifted up to a higher heavenly kingdom
merely by the other's glorious nearness, by what this yields us by its mere
existence. But what the one so inwardly adores in the other is nonetheless
in itself the divinity that, everywhere where two are gathered in its name,
is present in the midst of them. Therefore both can feel themselves so
infinitely humble, however much they are to each other. The more the
loving one is exalted and prized, the more deeply he feels himself as one
who is nothing, and who must himself take it as an infinite gift of heaven
that he can be so much to his beloved. Therefore love also fills us with
the sweetest gratitude, because we receive as a pure gift so unspeakably
much---a gratitude that does not, as it often otherwise can, have anything
conditioning about it, but that is as total and infinite as the very
beneficent power of love that calls it forth. Spiritual gifts are
everywhere more glorious than earthly ones, and are the more beneficent the
more, without these or those considerations, they enter in their full
wholeness into the receiver's soul. But of all that an excellent soul can
bestow, nothing is so beneficent as what it can give immediately by
itself, expressing the highest and most divine in its immediate existence.
And what the beloved thus bestows on the lover---now, moreover, giving
herself over wholly to being, for him especially, this all---is of all
gifts the one that awakens the inmost gratitude, however little there be
here any thought of an obligation. Hence, too, the inwardly moving quality
in love's feelings, which sometimes call forth tears of a pure joy. ---
And here it also shows itself that the familiar word, ``it is more blessed
to give than to receive,'' does not hold everywhere. On the contrary, he
who knows all the feelings that move the human heart might well as often
find occasion to say that it is more blessed to receive than to give---which
also entirely accords with the human being's inmost nature and essence,
according to which the total absorption in a higher life is the highest
living. Where what is given is only something earthly, it may indeed often
be ensnaring for the mind to receive; but to give, too, is then not seldom
likewise ensnaring for the mind. Where, on the contrary, what is given is
something Higher, or where in general something is given in a hearty manner
and out of full and true love, and now therefore a full, undivided heart
can come to meet the gift and open itself to it, there the one who receives
will most often enjoy life's highest fullness, enjoy love, more inwardly
and deeply than the one who gives. Already the fact that the one who gives
must act and do, whereas the one who receives can undividedly perceive, and
moreover undividedly and wholly give himself over, makes it so that this
latter can in a more total manner feel himself filled and seized. But
blessed we are only when we are rightly taken up and seized. Anyone can
become aware that the word of life by which we, hearing, are carried away
is more beatifying when we are wholly carried away by it than the same word
when we give it to others, even if it be given with its full power. The joy
that the idea awakened, conceived, and germinating in us brings with it is
quite another than that which the idea brings when we utter it and awaken
it in others. Indeed, even then joy enters rightly and fully only when we,
in giving, feel ourselves so carried away that in this very rapture we
ourselves receive a stream of life coming to us from a higher source. ---
And thus all the love that is yielded to us in the love-relation, and which
we receive with inward delight, has too something fine and sweet,
refreshing and inward in it, to which the love we yield does not always
attain.

Out of love's deep devotion unfolds its fair readiness to serve, and the
homage which, especially in the chivalry of the Middle Ages, took on so
fair a shape. At the beloved's feet the lover lays down every prize he
wins, and ascribes to her the honour for every victory. And rightly; for
she is the one who animated him with high life, and from her came to him
the courage and the strength by which he conquered. A Muse she is to him,
who inspires him in all his researches, under all his endeavours. Whereas
her mind is set on laying herself down, instead of all won tokens of
victory or accomplished works, at the beloved's feet, and giving herself
over wholly to being a faithful handmaid to him. What an inward joy and
delight it is to be able to sacrifice for the beloved, every happily loving
heart feels. Toil, work and difficulty become inwardly dear to our heart
when we have them for the one we love; indeed they become for us a need and
a necessity, that we may satisfy and, as it were, cool our forth-striving
desire for a total devotion. Only the strangest confusion in the
representations and feelings---the consequence of a deviation from Nature
and its significant intimations grown to an unhappily high degree---can,
especially in woman, have been what could bring so many lovers to let
themselves be led, by those who must seem to love without being able to,
into believing that they, on account of a refinement exalting them above
the natural way of living of others, had to let go of love's fair
readiness to serve, the sweet solicitude for a beloved individual's most
personal necessities, and the inwardness with which love is felt and
enjoyed when we work, sacrifice, suffer for it. Love is inwardly tender.
Out of deep solicitude that the beloved may fare well, may have everything
for the best, it gladly forsakes even its own [comfort] from the heart,
indeed is even fine and careful. Tenderness is deep solicitude, which keeps
the beloved in honour and holy, and precisely therein enjoys its feeling's
depth and sweetness. --- But again the heart is moved anew by an inward
delight to give itself over to sweet enjoyment for the beloved, that he may
feel the full pleasurableness of his possession, while he himself with
tenderness makes much of his beloved. Hence that fondness of love that
would enjoy by being enjoyed.

But as love is an unbounded devotion, so it is also an unbounded desire.
Just as inwardly the lover loses himself in his beloved and hangs upon
her, just as penetratingly is his mind set on meeting the same devotion,
and being able wholly to appropriate the adored being and take it into
possession. Both---the devotion and the desire for exclusive union---lie in
the drive, proceeding from love's inmost, toward fusion and complete
oneness. As need and longing and desire, love is just as exclusive and
ready to crowd out everything foreign, indeed just as inextinguishable, as
its self-giving and self-losing delight is inexhaustible. Both as a delight
to take the beloved up into itself, indeed as it were to devour her, and
also as a delight to be dissolved and to perish, indeed as it were to die
away in the other, love's divine madness (μανία θεία) utters
itself in a thousand even mimic outbursts.\footnote{``Divine madness.''} By
this exclusive desire the vehement, burning, restless element, shifting
between joy and pain, can now also enter into love's feelings, and the
loving one come to experience what Plato (Symposion, 403) depicts as
Eros's condition: that he, born of Poros and Penia, Plenty and Poverty, as
a being intermediate between the gods and mortals, bears the stamp of his
double descent. First---so it is said of him---Eros is ``always poor, and
far from being fine and fair, as most believe; on the contrary he is rough,
of unkempt appearance, runs about without shoes and without a home, lies on
the bare earth without covering, sleeps at the doors and on the roads under
the open sky; he has his mother's nature, and always dwells together with
Want. But, again, after the father, he pursues the good and the beautiful,
is manly, bold, enterprising, a formidable hunter, always contriving
stratagems, eager for insight and full of devices; he philosophizes his
whole life through, is strong in sorcery and in mixing magic potions, a
mighty sophist; he is by nature neither like an immortal nor like a mortal;
but now, on the selfsame day, he blossoms and lives, when it goes well with
him, now he dies away, but lives up again, according to the father's
nature; what he acquires perpetually runs away again, so that Eros is
neither rich nor poor, and likewise stands midway between wisdom and
folly.''

A deep feeling revolving about the individual's own I---a deep egoistic
feeling, then---goes, in love, into one with a deep sympathetic feeling.
Already in the merely organic respect the drive is, in its origin, an
individual and egoistic one, for it is an immediate consequence of the
organic tendency working toward perpetual self-formation and individual
formation. And afterwards, by this drive, not only is sympathy set in
motion, as it draws the one individual to the other, but next there also
arises the feeling of deep spiritual privations. Here shows itself an
inwardly strong need that demands fulfilment. Moved by inward joy over the
beautiful and divine, the attractive and refreshing, that is beheld in
another individual, the human being soon feels himself as one whose
isolated life is only a half---as one, therefore, who, in order to come to
full existence, must grow together with another being, so as now, in an
unbroken shared life with it, sharing everything with it, to enjoy life as
an unbroken loving. Therefore he desires a firm and inward and exclusive
union; not in this or that respect alone, but in all respects. He now
desires with all his might not merely to give himself over, but also to
appropriate the beloved and exclusively to possess her; not merely to love,
but also to be loved in return with the same inwardness.
But what the loving one desires is, however heavenly love's blessedness may
be, nonetheless always an earthly good. For it is an earthly being's
returned love and an earthly union. It is, on the whole, a good fortune
that is not always granted even to the one who longs most for it; it can
also be given and lost again. How dependent is love's happiness on earthly
collisions! Therefore the self, dissolved and vanished away in happy love's
deep feelings, cannot continue to be lost in a higher life's undivided and
unstaunched streaming-forth. The feeling of I-hood cannot rest. In a
thousand ways it is stirred up again, and, under a perpetual striving to
appropriate, hold fast and preserve what has been appropriated, or to
assure itself of the undivided keeping of its possession, all the heat and
ardour in the soul that lies in all egoistic existence is now easily
kindled, when the individual comes to feel sharply that it has to do with a
possibly opposing surrounding world. By the great excitability
(irritability) in the being softened by love there arises then, all too
easily, in the midst of the heavenly joy, a sensation coming from
jealousy's terrible hell. In all earthly possession there is danger.
Moreover, fear stands ready with the most hostile feelings. Already by its
mere existence the surrounding world can then bring vehement tempers into
passionate movement; and even if nothing hostile is stirred up in the soul
from that quarter, the struggle of individualities that must arise, when
two organizations already to a certain degree fixed are to melt away and be
loosened so that both may grow together into a new life uniting both, can
easily call forth a restless ferment and kindle resentment. Indeed, it may
perhaps be good if this ferment sets in early, under the feelings' still
richest fullness and deepest inwardness, when love is still most luxuriant
to crowd out all stoppage and to unfold its inner essence into outer
beauty. Such as the human being's conditions once are on earth, the pure
love wholly absorbed in its inner delight can hardly hold out, even in the
most innocent hearts, without their coming once to feel and cognize that
the gloomy darkness's joy-hostile might can reach even into the most
heavenly sphere of earthly happiness, when the human being does not guard
his heart. But what a deeply feeling temper suffers when it first feels the
possibility of a breach in the pure wholeness of the shared life is surely
often far more than anyone believes. Yet we will later find the place to
elucidate further love's passionateness---sometimes quite natural, yet again
often so highly unnatural. For now we will continue to consider love in the
fair shape in which it lives in the shared life of two happily loving
beings.

Genuine love, when it is a happy love, knits heart to heart, and makes of
two beings one. The more it grows and mightily reaches about in the inmost
of both, the more the shared life rises to the sweetest concord. With
unspeakable security and dependableness both hang upon each other. A deep
communication and confidingness, inwardly refreshing to the heart, is
between them. The one's whole life and inmost lies open to the other's eyes
with the purest truth. Already as soon as we begin to love and to hope for
returned love, we feel the inmost desire to unfold everything that lives in
us to the loving one, without reserve or reticence. Because we feel
everything heavenly and beautiful as something kindred and akin to us, the
heart feels itself drawn toward it, and longs to open itself wholly to it,
drawn by a deep sympathy, moved by an inward trust. Unspeakably refreshing
and beneficent is love's deep confidingness. A true liberating power is
there in this pure out-streaming of our inmost existence, this loosening of
every bond that otherwise so often draws us back into ourselves. As painful
as it is, holding back, to shut oneself up in oneself, so deeply refreshing
to mind and heart is this free openness, this undivided out-streaming of
all our feelings. It is not good for the human being that he be alone: so
ran the Lord's words already about the first human being.

How much does the human being not need to communicate himself and open his
heart to another---both in order that he may possess and enjoy himself and
his inner spiritual property with double security and delight, when what
livingly fills him, communicated to another, returns confirmed more clearly
and fully to himself; and also in order that he may free himself from much
in which he otherwise easily ensnares himself, that he may lighten his
heart, transform the inner conversation into an outer one, come out of many
an inner discord or disquiet with which he might otherwise grind himself to
pieces, and cast a clearer glance on his own life as he unfolds it before
another's calmer eye? But sweetest of all is this confidingness between
lovers, as those who are so bound together that whatever, with joy or pain,
strikes and moves the one, the other must feel as if it struck the dearest
part of his own existence. Therefore each longs to bring to his beloved
what livingly fills himself---only then able rightly to enjoy it, if it be
gladdening, or to bring it to clarity within himself, if it be disquieting,
when he shares it with the one his soul loves. Inwardly the human being
becomes, in this fair confidingness, confiding with himself, as he longs to
lay out his inmost purely and wholly for a beloved being's contemplation.
And it is more than a merely mediate communication through words that is
here. Immediately, as the two live, the one's life and whole soul lies open
to the other. Love's confidingness is not a confidingness in this or that;
it is a boundless confidingness, which words cannot exhaust, as if both had
an infinite secret to confide to each other. It is the expression of the
perpetually fresh striving to give the other one's whole heart. Thus in the
fullest degree there is truth between the two, since they inwardly rejoice
to walk and live before each other's eyes.

As the human being enjoys his life in another being, and again lives life
for this one, he feels himself as it were repeated, or set once more in the
other being. A oneness in a twoness constitutes the fundamental character
of the human being's person. It lies in consciousness. Whatever we live, we
live it as having to do with ourselves and referring it to ourselves. The
deepest expression of consciousness and personality, conscience, is
precisely a love with which we have and possess and embrace ourselves,
inwardly holding fast our true being and preserving it unlost. But in this
respect too we receive back, purer and clearer, the heart we lay down in
love's sanctuary. Whoever's genuine, true weal---a purer inner being---had
not become of heartfelt concern to him before, must come to this solicitude
for himself by love's power, when it is deep and genuine. What was to him a
deep reminder perpetually to guard his heart and keep himself pure and free
becomes so to him now with double strength. Just as nothing can be more of
concern to him in the beloved than the heart's guiltlessness and
purity---because only the beloved's true worth can give love itself truth
and subsistence---so must now, when he loves genuinely, his own being's
guiltlessness become unspeakably holy to him, since a beloved being beholds
in him the glorious image of its God. He would otherwise himself spoil his
love's pure beauty. Conscience's disquiet would now not merely be the
feeling of the breaking of his own inner peace, but would now also become a
feeling that love's pure peace was torn asunder. Thus love doubles the
power in conscience's inner voice. ``The human being is wedded to his
conscience,'' Goethe somewhere lets someone say, when the talk is of
married people who, even if they seem to live to each other's burden, are
yet as little to be parted as the human being is to be parted from his
conscience. But we will not see this as though marriage thereby got a
darker look, but on the contrary feel conscience's inward loftiness and
beauty by thinking of it as the one bound to us faithfully in life and
death, who is to be, for each, what in the fairest manner his beloved is to
him---the one for whose sake he feels himself driven, with deep devotion and
inwardness, to keep himself guiltless, so that he may always walk pure
before his love's countenance. And so love then lifts up the human being's
inmost self, and precisely as he gives himself wholly over, he regains
himself doubled; precisely as he mirrors himself, as he lives in and for a
beloved being's beholding, he feels his own existence's worth and
significance doubled, in the same moment when precisely thereby, at the same
time, humility's deep feeling is called forth. For the one of these
feelings does not abolish the other. The living presence and vigour of
both, working into the whole existence, belongs to full and vigorous
existence. And how free does love's deep feeling not make the otherwise
perhaps ensnared and bound soul, when the human being has in it a
dependable refuge under life's anxieties or assaults? Love gives strength
both to act and to endure.

A holy refuge---both under life's outer dispensations and under the
anxieties and sorrows a temper may have to do with in its own inmost, in
its struggle for a true existence---is a human heart that receives us purely
and with inward love: the fairest of all earthly temples, in which the
human being steps nearer the Eternal One, and comes to feel life light and
free, as he gives over all burdens to the one who can and will receive all
burdens that become too heavy for us, when they are laid down with inward
trust with him. The fairest sanctuary in which the hard-pressed heart can
find rest, the torn-asunder heart peace, and the burdened heart can pour
out what oppresses it, is a beloved being's loving soul, before whose
beholding eye the beloved stands, in his inner wholeness, purer and clearer
than before himself in his restless moments. How much and how often the one
beloved becomes the Eternal One's representative and priest, or priestess,
for the other, everyone who loves experiences; but it was doubly felt in a
time when gradually almost all temples were closed or lost their holiness,
and each single one saw himself more and more thrown back upon himself, and
saw himself without a community answering to his inner mind, because the
Spirit, however much it lived in individuals, had yet departed from the
congregations. Therefore, in this time, there was prized with double
strength the glory in that relation of life which still always had to
continue being poetic, and to be cognized in its beauty, because Nature
here speaks so strongly to every human heart. Indeed, it sometimes came so
far that the sanctuary's belauded guardians were adored as the very
divinity whose holy fire they kindled and nourished in the hearts.

But however high genuine love lifts us above the earthly's press and life's
lower sorrows and pursuits, however near it brings us the Eternal---indeed,
lets us feel the highest divine beauty's immediate presence on earth, and
opens for us the Eternal One's heaven; however unspeakably deep the
feelings of devotion and the confidingness of hearts that it calls forth;
however spiritual and pure in general the fair intercourse is between
lovers---yet love does strike its roots right deep down into the soil on
which it rests, as on its foundation, while in its blossoms it is ethereal
and breathes fragrance up toward heaven.

And it thereby shows its inner competence. For it leads to a marriage; a
family is grounded by it, and a house from which a new generation is to go
forth on earth. Not spiritual alone is the relation here, but also right
sensuous, and precisely therefore right sound, uniting both spheres of life
into a fair and vigorous oneness. The human being does not indeed live by
bread alone, but it does also live by bread, and cannot live soundly
without it. It is every finite being's determination, which must never be
forgotten, that it shall not merely strive to be blessed in its God, but
also fill its place in its world and, in the most definite manner, grow
into it. Otherwise it wanders about wayless, cast unsteadily to and fro in
life's manifoldness. Only the most exactly determined earthly existence
gives life complete steadiness. For here the human being is a part ordered
into a whole, and must feel himself as such; indeed he will the more fully
become a partaker in the life that lives in all, the more firmly he grasps
his place and attaches himself to it. The tree cannot raise its top toward
heaven without striking its roots down into the earth; whence else would it
get nourishment? And the human being, who, as has aptly been said, bears
his root within himself, must cultivate the field and nourish himself from
its fruits, and is thereby just as fully ingrown into the earthly life as
the tree is into the earth's lap. Thus love too has its sensuous
rootedness, and can precisely therefore hold out through a whole life.
Indeed, its inner beauty and power shows itself the more fully, the more
thereby life's lower occupation and desire gets deep significance for the
lovers, as they precisely thereby grow together daily more and more, and
enjoy the fair intercourse with each other, and test and experience love's
ever-more-unfolding fullness and riches, when it, on a thousand occasions,
makes life vigorous and sweet. In manifold respects, from the highest
spiritual to the lowest sensuous, the human being feels that he is not
determined to remain alone.

The natural drive that, under love's feelings, is in motion and gives them
their peculiar stamp lies, even where love's delight and longing has
transfigured itself into the highest spirituality, still constantly at the
ground, and plays into the whole of the feelings with no slight strength;
and it is not seldom evident how much it contributes to the
holding-together of two people bound for a whole life, under so much
occasion for irritation, and to the smoothing-out again of every breach so
easily.

We will here more closely consider what the doubleness of sex, and the
drive that knits both sexes together, has to signify; yet so that we,
instead of all merely physiological considerations, will call forth only
that to which, after all, one must always at last ascend---that to which a
glance over the whole of Nature leads.
Through the whole of Nature the organizing striving is directed at this:
that more and more that which is the true being, the inner spiritual, may
come to step forth, to utter and show itself as such. But this comes about
the more, the more prominent and definite (the more decided) the
distinctions and oppositions are in the manifoldness in which the inner
being steps forth. For the more does this show itself in its inner oneness,
as a unifying principle, and moreover as animating and vigorous, as ideal
and spiritual. Thus, in Nature's organizing striving, a perpetually
distinguishing and individualizing striving runs through the whole, and
what belongs wholly together steps apart from one another in the most
definite ways. Nature distinguishes and individualizes everywhere, but in
order the more fully, and in the higher and more spiritual ways, to unite,
the more definitely it has distinguished---so that life's inner One, its
inner spiritual being, may the more step forth as free living in a rich
fullness. And now the individualizing distinction goes even so far that,
under the rising of organization, the chief oppositions in every kind step
forth in two self-standing but wholly co-belonging individuals, in the two
sexes, of which each expresses its kind's nature only by half, and now, for
full existence, needs the other. The individuals here show themselves,
then, as the scattered members of the great organism of Nature, in which
the inner one Nature, under perpetual distinctions, exteriorizes itself;
they show themselves here as parts, belonging to each other, of a higher
wholeness---but a wholeness which human beings, by free self-determination,
are to bring to reality. Therefore human beings too feel themselves as such
members belonging to each other; they feel that the isolated, self-standing
existence is an insufficient one, and feel themselves driven, uniting
themselves with others, in communications with them, to fill up the
privations, abolish the defectiveness in their own existence, and come to
enjoy a more total human existence. The first, the nearest and most natural
step to this---that to which Nature most strongly spurs---is now that by
which the opposition which, of all oppositions in human existence, is the
most manifest, the opposition of sex, is dissolved into a higher oneness,
wherein, however, personality does not perish, but is precisely heightened
and furthered to its completion. When each has found its half, each is
first a completed human being.

The fusion of both into such a oneness happens by a union whereby there
comes forth a new, a third individual, who is to be, and also more or less
is, an expression of both fused individualities. But the new individual now
nonetheless always belongs only to one of the sexes, and the closed chain
is opened again; for the third individual too comes to need a fourth, and
the indifferentiation happens only by a series without end, so that life
shall not come to a standstill, but that, together with the individual
relativity and need, love's fulfilments may perpetually return, and love
may perpetually unite the distinguished. How spirituality here attains a
higher step in Nature is clear. Already in the animal a high sensibility
here comes more and more forth, and a deep instinct shows itself. But
sensibility and instinct are, in the animal, the most soul-like, that which
most nearly approaches spirituality. Instinct is a universal Nature's voice
speaking in the animal. But in the human being there comes here to
utterance that which everywhere is the highest spiritual, the highest
unifying: love.

The complete union in love of the two individuals drawn to each other by
love's strength, and a shared life springing from this unition, by which
life rises in power, significance and fullness for each of them: this is,
even with respect to Nature's purpose and in her own economy, always the
first and foremost; this is here the main thing. Procreation, on the
contrary, is only the second and secondary. For the sake of the coming
generation alone, the present cannot be something merely cultivated to
maturity. The coming generation does, after all, itself become once a
present one; and should that then again exist chiefly for the sake of the
next? No, the nearest is here also the first. The present time belongs to
us. In order itself, in its own existence, to express that fullness of life
to which then afterwards the next generation too is to come; in order
itself to be made partaker in this fullness and to enjoy it; and now at the
same time, under the earthly life's manifold striving, to be ripened forth
toward an eternity---for this the present generation lives nearest. Nor does
that tendency to procreation proceed from any other striving than the
individual's striving itself to exist to the full. Its origin was precisely
the individual formation-drive. Indeed, so accordant with itself is Nature
here, and so rich is life, that all that by which the next generation is
called forth and fostered serves precisely in the highest degree to give
life content and power, beauty and allurement for the present individuals.
On the lower step of organization we so often see the individual's fairest
blossoming-time, the highest point of the individual formation and
fullness, go immediately before---indeed coincide with---the point of time
at which the germs of the new individuals are formed; but, when these are
laid, the fair shape must now also often fall away or blossom out. In the
human being it is, on the contrary, by no means so, not even in the organic
and bodily respect, let alone in the spiritual. Man and woman stand in the
most luxuriant growth and fullness---indeed attain it only then---while the
next generation germinates up before their eyes and to their delight. And
while the father works and provides and gathers, while the mother watches
and busies herself and preserves for the children, indeed while both even
sacrifice and deny themselves much for these children, both come only then
to enjoy their own life, their own gathering and working and busying and
life's manifold property to the full; and they foster in themselves a
higher and more glorious life while they foster---a proper fostering, a
loving [Elsken], then, be it always---the same life's first stirring in the
little ones. Giving themselves over to living, not for their own self, but
for a beloved circle belonging to them, they now first rightly live. Life
is everywhere the more a true and genuine living, the more it is a true and
genuine loving [Elsken]. But by procreation the fairest love's sphere is
widened.

So let no one then, when under love's feelings the natural drive too is in
motion, hold that for unworthy or lowly which in Nature has the deepest
significance. Yet does this need to be said? Only in an age in
which---because the full power of religion and poetry no longer lived in
life's weightier and great relations---a strange sentimentality not seldom
brought feeling souls to want now everywhere to taste something holy and
poetic out of the objects, and to live solely on this everywhere, in all
relations---only in such an age could it be needful to say that one should
not disdain love's competent naturalness. There may indeed be peculiar
circumstances in which it can be a right natural, and therefore vigorous,
resolve of lovers to live merely spiritually knit to each other. But only
in a time in which souls sought, on the strangest ways, a replacement for
the lost communion with God---without which, after all, no one can hold out
life---was it to be feared that a sickly sensitivity had brought, and could
bring, this one or that to hold themselves too good for a sound, natural
life. Yet even a sound and pure temper may indeed, especially under youth's
fiery enthusiasm for the great and high in human life, sometimes come into
confusion by thinking of how the deepest and holiest feelings stand so near
the most sensuous. But let no one look askance at love because it is so
total, because its power stretches from the highest to the lowest in the
human being's existence, and casts nothing out in his nature. Precisely
thereby is it so competent and glorious, that it easily unites and
penetrates into wholeness that which otherwise so easily falls apart in the
human being's life, and which so easily makes it so that he lives in
moments, and as patchwork, the life that ought to be lived whole and as of
one piece. For among the things that confuse human beings much, and are
with difficulty brought to clarity, belongs this: how they are to carry the
higher life, which dawns for them in the fullness of ideas in their inmost,
into the daily earthly life, and, in the midst of this, to find again and
still possess and have the former. Love resolves the task easily; for it
casts its heavenly gleam over the whole earthly life. The lowest and
meanest in the daily earthly life becomes beautiful and refreshing to the
heart, because it belongs to the shared life in which the lovers live for
and with each other. Love hallows what it touches. A new and deep longing
strives to get its object, and if it gets this, life becomes, in a new
sphere, anew an unceasing caring and working for beloved beings, a new
service of love. Therefore let no one misinterpret the strength with which
longing often works, and privation strongly proclaims itself, in the inmost
of the noblest womanly beings. It is not longing's strength that here makes
a difference, but the manner in which it presents itself in the mind's
clearer feelings, and what here the conscious mind aspires toward. What in
noble woman-hearts comes to definite sensation and consciousness---and that
too only late---is precisely the longing for Nature's fair call to woman,
the longing for the joy that is won through pains and nourished in
sleepless nights, the joy that a human being is born into the world, as
Christ's words run in John 16:21.

But Nature must here indeed be natural. It is never so when it is not pure
and moral; for it is not so when, in wild outbursts, it breaks asunder its
own calm subsistence and inner peace. All genuine love is pure at heart,
and therefore also reliable, when only it is in itself tolerably clear. But
it is not always so, and thus there can be danger where there should least
be any. Yet, since the otherwise Natural is then so little natural and
little human, love too is, when only it is not all too rash or spurious,
itself in a high degree protective. In all good moments love is felt as an
infinite solicitude and safekeeping, and in its full, fair rapture every
desire dissolves into the most spiritual. The power of such moments must
then, by means of inner reflection and confirmation, spread penetratingly
to the rest. The more love's deep feelings come to clarity, the more full
the desire is with which the lover appropriates his love's whole fullness
and in all ways attaches himself to it and lives in it, the greater will
love's inner purity be. Let it then also, even under the closest union's
rightful freedom, preserve this purity, which keeps all
ἀκολασία away!\footnote{Licentiousness, intemperance.} For here
there is, even then, a limit which, for love's own security and peace, it is
dangerous to overstep.

But there springs, moreover, from a moral ground something finer and
nobler, which shows itself beautiful in deep love, and must always be kept
in honour. It is its bashfulness and shyness, esteem's tenderest
expression---a certain holy distance, answering to the natural opposition,
and persisting even in the inmost nearness and closest union, which keeps
the united ones out from each other, so that neither the self-own
subsistence nor love itself may perish under the strong attraction, and
both, in wild rapture, devour each other. Everything that touches us as
holy, and in which, as it were, a divinity comes near us, has something
arresting in it, that surprises and tames the mind's free delight. Partly
the I dares not, however much it be attracted, grasp at it, because it
feels itself as little and lowly, and now fears to be cast back; partly it
also itself strives against, so as not to be overwhelmed, until it finally,
nonetheless, with whole heart gives itself over, and now regains itself by
willingly and gladly losing itself. Indeed, the stronger the resistance
was, or could have been, in the first moment---the stronger in general the
person's self-ownness to hold itself fast---the stronger will the feeling's
intensity become, when now the heart must nonetheless give itself wholly
over. But in the ensuing alternating strife between the personality, which
will maintain the self-own subsistence, and the encountered Divine's
preponderant attraction, the heart must often first close itself yet more,
in order to be able to collect itself, so as to open itself right fully for
sympathy and love. It then even casts itself back in order to hold itself,
while it feels itself surprised and captivated, and thereby feels itself
embarrassed and unfree. Therefore there follows, with the bashfulness in
intercourse between the two sexes, in youth often at the same time a
certain shyness and severity---indeed a certain feeling of inner pride, to
which the mind as it were takes refuge, when the point of time nears at
which the heart is to feel itself drawn so strongly to another heart that
it is to lose itself to this. This pride becomes pernicious only when the
heart persistently will shut itself up in itself and close itself to love's
fair sympathy, and will root itself in its self-own subsistence. Then it
strives against Nature and against life's inner being, which only
distinguishes in order the more fully to unite. But even in the inmost
unition, distinction and opposition must not perish. Love is a oneness in
twoness, and in order that it may live with power, both the oneness and the
twoness must step vigorously forth. Therefore let both the self-ownness and
that holy distance remain and preserve themselves under love's devotion and
sympathy! Even under this latter's deepest power, a vigorous mind will yet
maintain its peculiar nature and its personality. This let the loving one
feel precisely purified and raised by his love! Indeed, love's inwardness
is, when only there is loving in truth, so far from suffering under the now
possibly ensuing resistance, that precisely often love became faint and
lukewarm because a certain resistance or counter-working, a certain
resistance, was wanting. But together therewith let there then precisely
preserve itself that tender, even bashful attentiveness, which will keep
the beloved being's inner personality in honour and esteem---indeed keep
it, like a divinity, holy. There is, however competent, and free of reserve
and bashfulness, the sound relation may otherwise show itself, yet in all
love something which, like the tender blossom, demands the tenderest touch.

Outwardly too that careful fineness utters itself, as well as that
in-itself self-withdrawing shyness, in the intercourse between lovers, when
noble love is to itself tolerably clear. The more inwardly its holy
divinity is perceived, the more will a finely feeling heart shun all such
audible outbursts or unconcealed utterances of its love, which could
profane it by giving it up to the indifferent onlooker, or by showing the
beloved too much in the drive's power; just as the finer feeling will also
lead to being everywhere attentive to the esteem due to every third
person---indeed to the wrong it can be against many a temper longing after
love's sweet joy, to stir up its longings by one's own happiness.
Everything deep and holy in general a noble temper will, though it may well
let it step forth through the word, yet only seldom let speak in immediate
feelings in the midst of daily life's intercourse, where the most different
thoughts and inclinations move the different people. Through the word it is
only spoken of; in feeling it is itself present. And it is not to be so
where it cannot be so to the full, and cannot move all with the same
feelings. Otherwise it is easily given up in a manner that either offends
or provokes to mockery. Much loveliness is needed in order not to offend
against all grace, when the human being is carried away by vehement
feelings' overwhelming outburst. Yet it is not precisely such reflections;
it is the immediate feeling of the holy in a deep love that makes pure and
noble souls bashful and reserved with their love in social intercourse.
But where lovers are gathered, they cannot easily be apart from each other.
Perpetually it delights them to be about each other and to occupy
themselves with each other. What way out does feeling then find, there
where the heart's deeper movements shun being seen? In thousandfold jest
and banter the feeling's inwardness veils itself. But while the lovers now
banter and seem to break themselves loose from their love's superior power,
they precisely enjoy the sweet intercourse with each other.

Yet also otherwise, on the whole, love's merry, often so comely wantonness
is just as natural as it is quick at hand. Indeed, precisely when, in its
fiery youth, the feeling of its deep earnestness is strongest, the mind is
just then most easily ready, in moments as it were to tear itself loose
from it, stirring itself freshly in the most unbound movements, as if it
delighted to drive a loose game with the heart's best possession and to
dissolve the feelings altogether. It is then as if the finite being cannot
endure so much fullness, but must sometimes, as it were, loosen itself from
it in order to be able to bear it. Indeed, so far does this delight go---of
enjoying under jest, and inwardly entering into one's love---that it is even
frequent and natural that a loving one takes delight in having something to
laugh at and to find comic in the beloved, which Tieck, in Phantasus, even
holds to belong altogether to love. Natural is this delight in finding
something to laugh over in the beloved; for it is natural to occupy oneself
with the beloved's idiosyncrasy down to its most individual peculiarities,
and perpetually to return to and dwell upon its contemplation. But for him
who, with this living eye, with this interest, beholds an individuality,
there is something stirring and irritating in it, because in every definite
individuality there is something peculiar and characteristic, different from
that in others---and that a something which, because it belongs to this
individuality, steps forth as something absolute, but now therefore as it
were jars and contrasts the more,\footnote{The Danish here, ``contrasterer
besmere,'' is uncertain in the source (the word \textit{besmere} is verified
from the scan but is non-standard; see the transcription note); rendered by
sense.} and thereby irritates one into bantering with it. Everywhere in
friendship and love we feel precisely a delight in finding, in the dearest
persons, something strange and droll; and let it be unwelcome to no one to
be the object of a laughter with which only the one who is loved is laughed
at. Even that which in and for itself is earnest and fit to call forth a
deep esteem can become the object of this jest, when it is regarded as
something personal, and now, at one time the contrasts, at another the
coincidences, between the wholly personal, peculiar or individual and the
universally valid and exalted above all concrete shape, become an
inextinguishable object of amusement. On the whole, there lies deep in the
human being's nature a delight---uttering itself under manifold shapes---in
an annihilation of everything that is something definite, so that, free of
regard and unbound by all form, the free life may live alone and stream on
unstaunched after its inner delight. All play is, on the whole, an inner
loosening or relaxation in the spirit, so that this, letting go of
everything that otherwise strains and draws activity to one side, may,
under an alternation in which no point holds another fast and everything
seems to want to go apart into sheer moments, feel a loosened existence's
free mobility. Under play and game, therefore, the spirit is enlivened and
refreshed, as it, unbound by every single and definite direction, stirs
itself freely and loosely to all sides and in all directions that present
themselves, without giving itself over to any single direction's straining
and moreover earnest-making power. Even in love, then, this delight now also
utters itself---in letting go of the enrapturing, moving and deep, and
giving itself over to an inner relaxation, under which nonetheless the
lovers, in the midst of the momentary feeling of unboundness (as if they
had now set themselves at liberty), come with double delight to feel
themselves as those who are bound and lost and are not their own, but also
to feel the sweetest blessedness in being thus captive. Yet, along many
ways, love seeks its way toward merry jest and wantonness. For in many ways
lovers are an incitement to each other---to laughter and merriment, as to
rapture, indeed to tears.

Yet let love guard itself against setting all its delight in these
incitements' perpetual alternation, as if it could possess and enjoy its
fullness only in fleeting moments and under the play of stimulations. Let
love guard itself against burning away in sheer craving, as well as against
scattering itself into a perpetually unsteady amusement; let it also guard
itself against what would be still worse: being torn to pieces under
perpetual irritations, in an alternation of embitterments and
reconciliations. Yet the perniciousness of these irritations the
consideration of love's passionateness will, further on, bring us to weigh
more closely. Life's deepest enjoyment is calm and quiet peace, the
undivided, pure, full feeling of something eternal and unlosable, that
under all alternation is abidingly present. The stimulating, however much it
enlivens and awakens delight, the enrapturing and moving, however much it is
the deepest feeling of the sympathy that therein moves and shakes the
mind---both let us here nonetheless feel ourselves as finite and needy, in
order that we may feel ourselves as blessed in our glorious possession. But
there is a higher soul-joy, which Jean Paul, in speaking of the Greek
poetry's plastic calm, so beautifully calls higher ``than Lyaeus's pang,
than rapture's warm, tearful tempest.'' What stimulates or moves strongly
works only transiently, and on every stronger awakening or straining there
follows a descent and relaxation; indeed, if nothing else lay in the
moment's enjoyment than what the moment brought, there easily follows an
emptiness. But underlying the perpetually alternating there lies, in genuine
love, something deeper: love's own inner living. This is the
non-alternating, the lasting, abiding, which gives the joy inner
unchangeableness. Love yields more than gladness in a thousand fair moments.
It is a single unbroken inner enjoyment of an inner glory that is near us
everywhere and present in all moments. So let it then now also be such a
one; let there be calm and rest, security and dependableness in it. And let
the secure feeling of certainty and dependableness not be perpetually
irritated and assailed at every occasion where this or that might want to
assail. How fair is not, in Schiller's Wallenstein, Thekla's inner security
and peace in her love, when in the midst of unrest for three days she does
not see her beloved? Whoever loves, let him learn to love with deep
inwardness, without the perpetual craving and chasing, burning and
thirsting, that does not even grant us the long and deeply after-working
after-enjoyment which first makes the hour's enjoyment a genuine one. There
are hours in love on which one can live whole days; but one enjoys them too
only inwardly when one lives whole days on them. The perpetual insatiability
and impatience that never treasures the granted happiness to the full does
not enjoy it from the heart's ground either. Under all love's alternating
moments---under wantonness's merry play, under shared dear occupations and
sweet confidingness, under tenderness's willing sacrifices, under rapture,
longing, melancholy and pain, indeed even under doubt's bitter torments,
under gladness as under tears---that which every lover feels and enjoys is
nonetheless always his love's depth and strength, and the essentiality in
the life that has dawned for both in their joy in each other. So let it then
now be so! So let the deep feeling of this essentiality, the mild, quiet,
refreshing feeling of a fair good's secure, calm possession---let this go
heightened, fuller and more inwardly forth from all love's alternating
moments!

\begin{center}\rule{0.28\linewidth}{0.4pt}\end{center}

We have thus considered love's glorious power and fullness from the most
different sides. In all the ways in which an object can bestow gladness on
us, love shows its vigorous and beneficent influence. Much in life yields
pleasure and joy through the beautiful and glorious that is therein
immediately presented for pure contemplation alone, and for the enjoyment
merely springing from this. Much again works deeply by awakening our
fellow-feeling, so that that which sets us in motion is another's joy,
another's worth or happiness---indeed is also the mere feeling of the union
that knits us to another being, and makes it so that we give ourselves
wholly over to this, and that nothing can be foreign to us that strikes or
concerns this being. In all these cases the feeling is the stronger and
more inward, the more we can forget ourselves in the object, the more
everything personal can fall silent and lose itself in the full devotion to
the moving influence. On the other hand there is much that, awakening the
feeling of a personal need and necessity, fills us with desire and craving,
and the enjoyment now consists in a satisfaction. Finally, all that
enlivens and heightens our personal self-feeling is of no slight power in
calling forth joy. In all these ways now works happy love's beneficent
power. So comprehensive is the fullness of feelings in which love spreads
itself out. Here is refreshing contemplation; inward sympathy and devotion,
deep fulfilment, and finally also high self-feeling through love's both
uplifting and liberating power---whereby even the unspeakable humility works
along and gives the feelings depth and inwardness.

But there is here yet another opposition of feelings to observe. To love
belongs not merely an infinite attraction, but also that distance between
the two lovers by which the twoness and opposition is preserved, without
which love's unifying and fusing power could not be perpetually fresh and
new. To this corresponds, in feeling, the esteem essentially belonging to
all love, which cognizes in the beloved being a worth and an idiosyncrasy
that is to be kept in honour and to enjoy its full right and furtherance.

To live fully there is required, next, both a sense for what life brings and
power in working back. Most fully does he enjoy life who both meets its
enlivening manifoldness with open sense and light mobility, and also, at the
same time, with a penetrating mind grasps its inner being and depth; who is
enterprising, vigorous and enduring in deed, but yet, under all that stirs
one to grasp with energy into life, preserves the mind's inner equilibrium
and calm. These are the moods everyone should strive to draw forth in
himself and to preserve. With a light mind and with a deep mind, with a
vigorous courage and with an even, composed courage, life must be lived.
But thus does the happy lover live it. In love there is a sense for
existence's liveliness and for merry jest, and a sense for life's deep
earnestness; here is energy and zeal in willing and striving, and at the
same time here is a calm security in the full feeling of a glorious estate's
dependable possession.

But to this power and fullness love comes only gradually. Like everything on
earth it must grow up from a germ and unfold little by little, and, in
time's advance, attain its full blossoming---to which, however, it can come
only as we ourselves draw forth its power and fullness, and ourselves let
its inner warmth penetrate the manifold members in which our inner life
stirs, so that its moving and enlivening power may from the heart spread out
to the extremities.

All inward love is in its origin genial. A deep nature stirs of itself in
us, and the power of deep ideas lies at the ground of its working. Its inner
dark movements call forth love's strong longings before it finds its
object; and when it finds this, then the human being is, before he himself
yet rightly knows it, in the midst of love---seized, taken up and captive by
the strong feelings that have of themselves filled his inmost. He has
immediately lived himself into his love, and loves because he cannot do
otherwise, according to an inner irresistible necessity---yet such a one as
is at the same time a high freedom, and is inwardly felt as freedom, because
it is felt as that which proceeds from the individual's own inmost being.

But what has thus of itself sprung up in the human being, and, as a mighty
natural influence of a higher kind, moves him in the strongest emotions,
this let him now come to inward clarity upon in himself, and let him
appropriate it from the heart and to the full! Let a clear composure step
here, as everywhere, to the genial, so that the possession given by a
higher inspiration may not be possessed with blindness, or rapture's fair
uplifting pass by like an intoxication. But the genial (τὸ δαιμόνιον, as
Plato calls it) the composed one appropriates, partly
cognizing, partly willing.

In a cognizing, feeling and willing has, on the whole, the spiritual life,
and all that lives therein, its subsistence. Therefore there must now also
come, to love and the fullness of feelings in which it utters and moves
itself, both a cognition and a will. By both the whole love, with all that
it contains, is as it were confirmed and posited (\textit{poneret}) anew.
The human being must with attention attend to his love's content and
many-sided power. He must tell himself what it is that has now happened to
him; he must seek with clarity to grasp and cognize what has rightly been
allotted to him, and thus work toward its being able to live and spread
itself in him with its full strength. Then he first comes to enjoy it to
the full and rightly to preserve it for himself, as it now knits itself to
the whole remaining content of his inner life, and, while by the clearer
consciousness it is itself clarified, now at the same time shows its
beneficent power in all directions.

But just as, on the one side, by a cognition, so must we, on the other
side, by a firm resolution, appropriate the new life that dawns for us in
our love. We love because we must love; but what we thus irresistibly must,
and feel ourselves driven to, this we must now also will. The loving one
wills his love---so let him then now also will it from the heart's ground,
with the inner freedom of a self-possessing mind. We live ourselves into
love. But the inward, exclusive union for the whole of life is to be
expressly entered upon. What has happened of itself, the human being is,
with free self-determination, to confirm---therefore also first to test
himself whether he too can, from the heart's ground, will, and really does
will, what he is now himself to resolve. Let the relation then become a
genuine moral relation, founded by a free will's and resolution's firm
determination! In this firm, holy resolution---now also to will from the
heart to appropriate the won happiness, and to hold it fast and wholly give
oneself over to it, and thereby to fix oneself and the easily moved delight
and drive---fidelity has its root and its seat.
Fidelity must constitute the kernel in the glorious companionship in which
two beings, bound for a whole life, perpetually grow more and more
together. In genuine love it perpetually increases, under the shared life's
gracious course and under the feeling of both parties' indispensability to
each other, and, fastening itself firmly in both hearts, it fastens both
hearts to each other. By it love first gets the stamp of essentiality.
Nothing is sweeter, deeper and holier in love than the feeling of this
fidelity on our own side, and the certainty of it on the beloved's.
Whereas the thought of infidelity is so revolting and so terrible, because
infidelity makes it so that even the higher, more glorious in the earthly
life shows itself as something that has no subsistence, no genuineness,
but, subjected to earthly upheaval, is nothing, like everything else
earthly. All that love can yield and bring, joyful and beautiful, it is
only under the presupposition of this certainty, security and reliability.
Therefore the human being feels in his fidelity the firmest and strongest,
and in the certainty of it from the beloved's side the most of concern, in
the whole of love. What could be more terrible than, in the midst of love's
fullest feeling, already to have to fear the worm which, hidden under
life's up-blossoming bloom, should undermine even the most spiritual and
pure. In our inward fidelity we press, in spirit, our beloved to our heart,
indeed bear the beloved as it were enclosed in our own breast, as the
holiest part of ourselves. Therefore we perceive love's feeling as an
eternal feeling, and call it so. How wasting, how ruining, on the contrary,
is infidelity. Let one imagine a genuine, inward love-relation; let one
imagine that fair fullness of life in which love spreads itself, how
everything that surrounds the human being gets a new, more lively look, and
a thousand trifles become means for the utterances of that love and
heartiness with which perpetually the hearts open themselves more and more
to each other and knit themselves to each other. A rich estate of the
fairest memories the human being takes with him into the future; life is
enjoyed not in moments alone; the deep, lasting feeling that goes through
all moments gives the whole life continuity. And now let one imagine love
vanished, the relation dissolved, be it also only inwardly and spiritually.
It is now as if a hostile demon had stepped between the human being and his
fairest life-period; to all memories of this attaches the thought that it
is all vanished, and disturbs and embitters them. The human being is as
cut off from his present's root, the fair past; the glorious property he
possessed in a rich memory is gone, his life is as if cut through, and its
sweetest part cast away. The whole long, fair youth-time is even less than
a vanished dream; for on this one can still think with delight, but no
longer on that. So ruining is infidelity, that it not merely dissolves the
present into a nothing, but even tears asunder the preserved image of the
whole past. To be sure, there are cases in which a human being, in order to
preserve fidelity toward himself and toward God, must---however much the
heart may bleed thereby---tear asunder a bond that knit him to another human
being. We are faithful to the beloved because we are faithful to the true,
good and beautiful. Fidelity in love has its root in a higher fidelity, and
must fall away when it cannot be nourished by this, but contradicts it. Yet
such cases we shall come to consider more closely when we come to speak of
unhappy love.

But even outside these cases we still see, all too seldom, love call forth
the union it could and should bring with it: a true, heart-refreshing,
inward soul-union for a whole life. How fair and comforting it is to see
people who lived with each other for many years---indeed until age made them
infirm---still love each other with inward devotion and deep confidingness
and trust, indeed still with youthful delight and pleasure rejoice to
behold each other. In youth everyone who loves with fieriness and
enthusiasm hopes, indeed is assured, that his own love will have this
endurance; he is sure that the spirit and the heart he loves in his
beloved, exalted above infirmity, will still, out of the aged form, speak
to him with undiminished power and still gladden and enrapture him. But
while the world is perpetually full of young lovers who are quite taken
with each other, it is at the same time full of married people with whom
marriage seems to have become what a current word even says it is wont to
be: love's grave.

The first thing in fidelity is the firmness and reliability that brings
with it the unspeakable certainty with which the one relies on the other,
and counts on a disposition in the other unassailable by all earthly
incidents and collisions, according to which both perpetually continue to
be to each other what they have been to each other. They feel themselves
knit together partly by the power of the good and glorious they behold and
cognize in each other---wherefore, too, with bad and reprehensible beings
such a pact can have no subsistence or security; but partly they are also
knit together, and continue to be to each other what they are to each
other, because they have been it; because life once gathered them and let
them, with and by each other, behold into its deeper being and enjoy its
deeper fullness. Whereby it shows itself that to him who has, there is given
in full measure, as every moment on which memory with joy dwells makes the
union more inward and firm.

But there is yet, in lovers' fidelity, another thing to remark, which does
not, like that first, hold of every soul-union, or every genuine and deep
friendship. This other thing is the exclusive element in the union: that
both not merely, with unassailable firmness, hold to each other and have a
sure and dependable home, an unfailing refuge, with each other; but that
also neither of them desires to be, or is moved to be, the same for any
third; that love, therefore, unites only two in this manner. Here too
shows itself the power of a firm resolution and will, and here too it is
needful. It is the natural drive lying at the ground of love that hereby
comes into consideration. This often easily movable drive, which can be
drawn to all sides, the human being must himself fasten or fix by a holy
will's power, so that it may not, stirred up to disorderly movements, drive
him unsteadily and waylessly about, and set the mind in unrest and discord
with itself. Here is something on which a bond and a bridle must be laid,
and where a protective feeling of duty, a feeling of respect for the holy,
must step to, be nourished and preserved. Only all too often did a lack of
spiritual power in this respect lead to that which, instead of stilling and
refreshing the heart and bringing fullness, left the heart empty and laid
the mind waste.

Meanwhile: is love to bring with it that life's outspread glory shall lose
its allurement for the loving temper? That love which with such mighty
power awakens the spirit and opens the eye for the beautiful and divine in
life---is it now to lead so far that the same eye, fastened on a single
object, shall be blinded for all others, or close itself to them? By no
means. However holy the union to which love leads must, in that one
respect, be kept as an exclusive one, yet let there for the rest---outside
that which belongs to the nature here in motion---reign the envy-free mind,
the ἀφθονία which the Greek so praised in his gods. There must be loving
with a sense of freedom. So far shall the exclusive not go, that now all
contemplation's joy and all spiritual communication, all living interest,
shall concentrate itself about the beloved alone, and the rich, widely
outspread world cease to be, to each of the lovers, what it formerly was to
them---a world in which heartiness, solicitude, interest and love found a
wide sphere for themselves, and the heart everywhere rejoiced to meet its
God. Do not, then, you who love, begrudge your beloved the merry, lively,
open mind, the fresh, ever-ready joy in life! Nourish much rather, in this
respect too, the inward, sweet solicitude, the hearty safekeeping, that
wills that the beloved shall have everything for the best! Behold it with
joy, if a higher joy in the whole of life, and in all beautiful and great
that moves in life, springs from and is nourished by your beloved's love!
Allow the beloved's happiness to become a true happiness for this one! Love
shall in truth not narrow; and what opened the free glance upon life shall
not afterwards becloud it. Not even the first, foremost and highest object
of his beloved's high esteem can the beloved always be, as little as the
sole object of the loving one's whole confidingness and its refuge in all
the affairs that oppress or move the mind. Love, which in so high a degree
can make seeing, shall never ensnare the pure, clear judgment. No one
becomes unfaithful to the beloved because he becomes faithful to himself,
and will not look down on or pervert what speaks as truth in his inmost, or
because he strives to fight the fight we all have to fight, with an upright
mind that grasps every help and support it needs. Only with respect to the
not-always-duly-ordered Nature's stirrings, let everyone guard himself and
have the eye watchful over himself! For here is danger; and where something
innocent and good lies too near something that is not good, let one rather
let go the innocent and good than expose oneself to danger and unrest, and
forget not to warn oneself or each other! For without inner peace and calm
the good does not, after all, come to our good.

To love with genuine sense of freedom, and let egoism rightly perish in
love, is, however, not easy to learn. But one consideration may perhaps
still elucidate what has been set forth, and open it an entrance to the
heart: a glance at the other life.

Everyone who has tested life and love, even if only tolerably; everyone
who, desiring and yearning, has felt that goad and spur of I-hood which is
in all desire, especially when it is vehement---because under all vehement
desire and yearning we come to feel ourselves in sharp opposition to a
possibly opposing surrounding world, to an object perhaps not immediately
accessible to us; everyone who has tested life and love and the fates of
both, and has experienced what both give, and how both are lived and
enjoyed, whether there is throughout in them clear calm and pure, quiet
contentment, or whether they can disquiet both, and how and why:---let
every such one ask himself how he will imagine a life in eternity, when he
is to think it as a life uplifted above earth's pressures, pure and clear
and blessed, but yet as a life in which all that was loved is found again,
and otherwise so quite humanly that we can speak in earthly wise of it. Is
that too to be a life in which exclusive possession is to reign? Is the
blessed life too to be a burning one, and is not even the homegoing to the
true home to bring calm and peace? Or must not rather the home of
blessedness be also that of peace, and must we not rather hope that in this
home, when life has raised itself to genuine freedom and full personality
and individual fullness, now the goad of egoism---with which we here were
plagued, while we were thereby spurred on---may fall out; that in the
kingdom of glory none woos and none is brought home, because there is no
talk of possessing? So that, to be sure, what here constituted a fair
shared life's inner being and fullness is by no means to be believed lost
or over for a following life; nay, that on the contrary in this respect too
life's spiritual property here won is taken along yonder; so that lovers
understand each other inwardly and hang upon each other and love each
other, precisely because they loved here---indeed that something quite
special may even continue in their love; but that yet herein one being does
not demand anything above another, or that in general a ``before'' is not
heard, or can even so much as come up in thought. For where a ``before'' is
thought, there already stirs the poor I, the paltry self-regard, under
whose ensnaring, arresting influence we have truly here on earth suffered
enough. Let one only ask oneself, too, whence yonder a desire for something
exclusive should come. In the pure joy of contemplation, the inward
enjoyment of the unspeakable glory that opens itself for our eye in a
beautiful being, there is, after all, nothing for which we should begrudge
anyone, whomsoever, being made blessed thereby, wheresoever. And as regards
sympathy and participation, and the devotion with which the one gives
itself over to the other, there is here, for the one who enjoys the
happiness undividedly to possess another being's whole heart, no ground
enviously to be angry that the same heart is just as beneficent for another
being, and gives itself just as purely over to this. For it is a pure
heart's power and its happiness that it can, in its full wholeness, give
itself over to a thousand, undividedly to each---when, namely, it gives
itself over in such a way that in and with the same (\textit{eo ipso}) it
gives itself over to the Lord, whom we meet on a thousand ways. --- Yet, to
this high freedom from envy, this ἀφθονία, only very few pure, undemanding,
happy natures are able to bring it, in a life in which the goad of egoism
can so hardly come to rest. Meanwhile, let everyone yet strive, on the
whole, to live and be stirred in that spirit, and with that mind, whose
full, undivided working he hopes to enjoy and possess in another world!

\begin{center}\rule{0.28\linewidth}{0.4pt}\end{center}

%% ============================================================
%%  SECOND BOOK
%% ============================================================
\chapter*{Second Book.}
\markboth{Second Book}{Second Book}

Love should and can, according to its inner nature, lead to a deep
heart-union, whereby two, united for a whole life, perpetually grow more
inwardly fast to each other in a fair shared life, to the pleasure of God
and men. Why, then, does love so seldom lead to the fair union to which it
can lead? Love is like a fruit-tree in spring, which, covered with
blossoms, sends out a fragrance that proclaims the coming summer's
fruitfulness and fullness. Marriage should be like the fruit-tree that
bears the whole year through, at once, blossoms and fruits, beginning and
full-ripe; for such is the spirit's power in the human being, that his
blossoming-time is not over because his fruit-time has come. Why, then, is
a marriage so often---even with quite worthy and good people---like a tree
whose fruits are some withered, and some worm-eaten, and hardly a single
one right fully ripe and alluring?

First and foremost because people so often go heedlessly into the important
relation, and so seldom remind themselves, when it has been entered upon,
that they must themselves take care to preserve for themselves what has
been bestowed on them. Already before the union that is to be valid for the
whole of life is entered upon, let one watch over one's heart, and not
believe everything to be secure and reliable because the awakened feelings,
with vehemence and strength, concentrate themselves about a chosen object.
The power of beauty and grace---of certain outer allurements in general,
working on the feelings through the awakened sense---is most often that
which draws love forth and captivates the gaze and the inclination. But one
forgets, all too often, to tell oneself in time even this: that a
love-understanding is to be a deep understanding, a genuine and deep
communication of hearts, and that lovers therefore ought to strive to enter
into each other's natures. One cannot repeat this too often, when one sees
how frequently the whole loving stops short at the delight and pleasure
with which two see and embrace each other; how frequently people fall in
love with each other and knit the union without their sympathizing, without
their understanding each other from the heart, without a deep and
heart-refreshing spiritual communication being between them, and without
their thinking to draw forth and foster the heart-intercourse that might
yet come about. The love that springs from a delight of contemplation---be
it moved even by the strongest drives---that does not lead to this deep
confidingness, to the inward heartiness and deep dependableness, is no
genuine love. Yet it is not seldom that a certain power of beauty and
grace, or a certain power of spirituality and quickness, so draws one being
toward another, and so stirs up all appetites, all drives, all inclination,
that the one feels captivated by the most vehement love for the other,
without there being true accord and sympathy, and without its coming to a
true understanding and communication of hearts. But if this be wanting,
then it comes to this: to be upright toward oneself and tell oneself that
it is wanting. One must always have courage enough to be truthful toward
oneself. Let no one ever forget that love is not merely to draw an eye to
an eye, but also to knit a heart to a heart. Let neither the inward delight
and joy of contemplation nor the drive in motion carry away anyone who
desires love's full and deep blessedness, before also a genuine sympathy
shows itself, and a deep and full communication has stepped in between the
two! Not even a sharp and cutting feeling of a privation and an unsatisfied
longing can rest so heavily and oppressively on the soul as the feeling
would that comes of being bound for good to a human being ``who does not
understand our inmost being, and yet unceasingly demands our regard. That
longing at least fills the soul; but such a relation offers us, without
cease, the everlasting emptiness.'' But how heedlessly do many people rush,
or let themselves be drawn, into such relations. And one's own blindness is
only all too often supported by others'. Especially is there many a case in
which one could wish a daughter strength of soul enough not to make a
solicitous mother unhappy by letting herself be made unhappy by
her---when the mother, out of a little-Christian anxiety for the future,
will by representations move the daughterly heart that is moved by no
inclination. The most childlike compliance and fairness does not cool the
blood, or satisfy nature, or secure against a hard struggle and torn
feelings; and an unhappy daughter is too great a punishment for a mother's
impatience for a tender and loving daughter not to spare her mother it by
steadfast resistance. Not to speak of the cases in which impatient parents,
instead of a simply unhappy daughter, get a corrupted daughter to bear on
their conscience---and moreover a doubly unhappy one. A marriage entered
upon, after many deliberations, with half a heart, is in truth not
concluded in heaven---and that, according to an old word, every marriage
ought to be. In the man, moreover, the heedlessness and folly is just as
great when he marries without being loved, be he even driven by the most
vehement love. And is it then to love, to make unhappy the one you care
for? Let no one rely on the returned love's coming after all. If it can
come, it comes easiest and soonest while the beloved's heart is still free.
Once the hard pressure of a feeling of duty lies on the subdued heart, it
comes more slowly. And what folly, to skip over the first and most
important thing in order to hasten to that which is, after all, only
something under the presupposition of this first, under the presupposition
of returned love. And is there anything so unmanly as, with a kind of
overwhelming and compulsion---the spiritual compulsion too is a
compulsion---to take possession of a woman who yet resists with all her
feelings?

Two periods in the human being's life are, moreover, here well to
distinguish. In life's first youthful up-blossoming nothing can satisfy but
the highest and most glorious. The deepest fulfilment for mind and heart is
sought, and a union without fiery love and inward captivation is here
nothing. To give oneself over for outward earthly considerations is in all
cases an almost still greater folly than that of selling one's teeth for
dainty dishes. What enjoyment is anyone to have of dishes when he has no
teeth, or of riches and earthly loftiness when he is unhappily married? But
in a youthful period the individual cannot even give itself over merely on
the strength of the cognition of the other being's worth alone, when an
inward inclination does not at the same time draw the heart and move it
with irresistible love. To give oneself over without proper love---even if
the heart be otherwise moved by feelings of genuine esteem and
goodwill---is in the youthful age always precarious. For it can then easily
happen that either love's fiery and lively awakening, and moreover its
whole soul-warming and enlightening and blessing power, is thereby
suppressed and stifled, so that the human being thereby excludes himself
from the happiness that might yet perhaps once have opened itself for him;
or---and the danger here is even greater---that a deeper need and delight
does once really come to awaken, but that it now awakens as a disquieting
one, that brings the heart into discord with itself and with the deeply
stirring nature, which cannot find fulfilment, calm and refreshment in that
which the calmly deliberating but not at the same time irresistibly drawn
resolution has grasped. On the contrary, later---after nature has stirred
in its youthful strength, after life has been tested and the feelings have
spread themselves in the individual's inmost with the power and in the
manner that the life allotted us brought with it---later, if the individual
in that first period did not find and come into possession of love's full
happiness, another state will set in. Whether the heart has loved but not
found union; whether it has loved, on the whole, with more or less strength
and inwardness and with this or that happiness; or whether it has also not
loved properly, but yet, under the longings' inner power, felt love's deep
sweetness; or whether it has lived on more calmly and been little moved by
love's delight or longing:---in all cases life has now formed and shaped
itself and taken on a more definite figure; the fullness of feelings has
worked as it could, and nature has gone through its first stage. The
individual now knows what love is. To be sure, the more definitely it knows
this, the more clearly it has learned to survey its state, the more it
brings a demand with it---instead of the young one gladly taking what is
given, as it is given---the less can it now find an object. What makes the
matter so easy in youth---that life itself is first to form itself
spiritually, and therefore easily forms itself with and by, and moreover in
harmony and agreement with, a foreign individuality, because everything is
still soft and in becoming---this is now no more. The more firmly life is
fixed, the less is it fit to be loosened. Meanwhile, let the drive too have
been held down---first by the youthful enthusiasm, later by the calmer
cognition's and will's power; let the mind for a time have subsided: nature
does not subside, it demands and presses. Its movements stir up again, and
the human being will now easily come to find true Paul's word: it is better
to marry than to suffer ardour. Now he is content if, instead of youthful
love's heaven, he gets nothing other than the earthly happiness to which
love can lead. And if his striving is pure, if he is not merely calm and
composed, but also faithful toward himself and truthful---so that he guards
himself against ensnaring and entrapping himself, or against hurrying and
stupefying himself, and lets every movement that now will carry away the
perhaps restlessly yearning nature work slowly, until it comes to clarity
and, when the relation is first surveyed, can find the calmer
contemplation's accord---then it may still happen that the inner natural
drive's strength can so loosen the heart to cognize and feel all the
beneficent that is beheld in another being, that he can still, in his love,
find that heaven without which even the whole earthly happiness is still too
little. So that yet a true inclination and heart's attachment to another
being can draw forth the deep confidingness and sympathy, as well as the
true contemplation's delight and joy, and a genuine spiritual fulfilment.
Let no one look askance at a love-relation that is knit in this calmer
manner in this period of life, in which youth's fresh, glorious hopes must
fall away. Let no one look askance at anyone because the Lord was less
gracious to him than to others, and because he learned to tell himself
this. Such marriages too can be---to use the old word again---concluded in
heaven, or concluded on the strength of so pure and sure a feeling's voice
and prompting power that the individual is conscious that it wills because
it must will, and that the outcome proves the feeling's genuineness. Only
here it is doubly to be remembered---what the youthfully moved soul yet
feels soonest---that no love-relation is anything when in it a heart does
not knit itself to a heart, and the love-understanding is really a deep
understanding and communication. All marriage must be entered upon in such
a way that an inner higher voice's dependable accord lets the individuals
become certain that it is not caprice's play and one's own fancy that
drives them. Every marriage must be concluded with a devout mind
(\textit{pio animo}). But it is herein never to be forgotten that what may
be in its right place in life's second period is not therefore at once so
in the first, youthful one.

Meanwhile, that marriage is not always concluded in heaven which vehement
delight carries two away to enter upon before they come to collect
themselves---that is, to reflect and come to some clarity about themselves.
With a union that is to be valid for the whole of life one need not, after
all---especially when the heart first feels sure in its cause---make haste.
There is, all the same, so much reason to go forward calmly and slowly. Let
one grant a young couple, and let it grant itself, a lively period of being
in love, and let it enjoy in slow, calm draughts the mild, refreshing
fragrance which, announcing a wonderful future whose fullness seems
unspeakable, breathes through love's fair May-days. It is especially for
the woman so beneficent to enjoy a long love-spring; for she is to go to
meet a calling, and must have calm and time to gather mind and reflection
for the wholly new sphere of activity, and to prepare the heart for the
time that awaits. And is it not, too, a pity for a womanly being in youth's
first blossoming---now already at once, after a new life has dawned for her
in her love---to have to be drawn into marriage's sorrows, indeed to bear
the mother's burdens and endure the restless nights? Not to speak of the
fact that the over-hastening can easily lead to the unnatural, when the new
germ is forced forth before the first fruit is full-ripe, so that the
mother, even in the organic respect, still herself grows out, while already
a new being grows under her heart. Whereby it is to be remarked that
maturity, even in the organic respect, is not yet at once there because the
full growth is there. To this comes that, to the organic full-maturity, a
certain, only later-arriving spiritual maturity must still come and be
awaited, so that the man may yet be far enough to be able to be husband and
father, and the woman to be able to be wife and mother. Parents do indeed
first learn from their children to bring up children; but they must yet be
far enough that their children's willfulness cannot make them themselves
willful. But unfortunately here too, not seldom, that vanity which ruins so
much---often also a restless impatience---gets its pernicious influence.
Nature's early up-surging movements and promptings one should not at once
take for a voice of Nature; for that this nature, in most people, comes
forth all too early on the whole, long before even the full growth is
attained, is well enough known. Yet also for lovers in later years it is
important and beneficent not to hasten with so important a union. And that
for a reason which is highly important for all individuals: that they must
first learn to live with each other, and to look each other right into the
soul and fully know each other, in order that they may in life itself test
whether they also belong to each other, or whether the drive and delight,
easily awakened and carried away in youth---often also in the later
age---has perhaps ensnared them into an error.
What the human being properly ought to have in all life's cases---and what
we, when we are not so reborn that we can have it in all, ought yet in all
important cases of life to await with calm devotion (but of course not
await under a torpid indifference, but under living working toward
it)---this it is especially so important to await before a marriage is
concluded: the higher verdict, by which a voice coming from our inmost,
which is not our own, says it to be good, what we have in hand, or else
disapproves it. Precisely then, when the individual's delight and
resolution finds this inner higher accord, is---but also only then---that
done which the word just cited so beautifully calls the marriage's being
concluded in heaven. How happy would human beings be if everywhere and in
everything they so had God before their eyes that, in daily thought of God
and prayer to him, they were accustomed to bring their affairs before him,
and to take life as a service of the Lord, in which it came down to willing
what he wills. They would then, to be sure, far more be what human beings
so often are not, and what it is more difficult to be than it can seem:
truthful and upright toward themselves and with themselves. Only all too
often do human beings hasten impatiently forward, without collecting the
mind and going into themselves, and uprightly giving themselves account of
their heart's and mind's movements; indeed they have not, even when an
inner feeling definitely warns them, the courage to listen to it, not the
courage to be sincere toward themselves. Meanwhile there are, to be sure,
cases in which a single moment's unspeakable certainty can so have the
testimony of its truth in itself that an instant can decide for a whole
life. But these highly genial moments are as rare as all the geniality that
anticipates time. And precisely in such a moment must the two, who behold
each other deep into the soul, feel themselves unspeakably near the Eternal
One, as they feel themselves unspeakably near each other; precisely in such
a moment must the union of the hearts be a highly religious one. And yet
even so it is necessary to wait, and let the new life come to settle and
come to a higher clarity. Everything on earth---even the highest spiritual,
when it steps into earthly beings---is subject to growth; and what is
conceived and born in spirit and heart is just as little carried to term in
one day as the human being himself was. Let one consider that life, in the
spiritual respect too, just as it has its crises, so also has its stages,
and that it is dangerous to hasten heedlessly over anything. It is no small
thing that is to happen when two individualities, fixed to a certain
degree, are to be loosened in order to melt together with each other. In
each of them, in certain respects, the firmer form must be dissolved in
which each had perhaps already set itself fast, so that each may open
itself to the new life's penetrating warmth, and both may rightly become
one. The familiar saying, \textit{corpora tantum soluta agunt} [only
dissolved bodies act], can also be carried over to the spiritual. But this
loosening happens not seldom more easily, mildly and conveniently before
the firm bond is knit, when the mind is still freer, and no feeling of
necessity ensnares or stirs. Let no one rely on that which is to come, but
has not yet come, being able after all to come after the union. No,
precisely because it is then to come, it sometimes comes with more
difficulty and more slowly; and here too the old proverb returns: that the
times are equally long, but not equally useful. Yet this is, to be sure,
not applicable in equal degree to all individualities; and when in the
midst of spring summer's heat and sultriness will present itself, then let
it just as gladly become summer at once. A very long betrothal-time is, to
be sure, still less desirable, unless it began in the earliest youth.
Love's first fieriness and fresh fragrance must, above all, not be over
when the marriage is begun; and let one rather preserve it by contentment
and moderation than have consumed the best before the best point of time
comes. But it always seems to be very well, and should never be
neglected---what even the Church's ordinance formerly required---that the
bond be, by a solemn act, knit provisionally by a betrothal or espousal
before it is knit to the full. Only the interval between the two points of
time should not, as formerly happened among us by laws, be restricted to a
certain unalterable brevity, but precisely widened to a span of a suitable
length, within which the marriage should not take place. And the
provisional union should be able to be dissolved at the one party's request
alone, because it was precisely regarded as a trial of souls. (A familiar
custom in many regions of Germany and other lands, the window-visits, can
only on the strength of a similar consideration have been upheld by public
opinion, which yet otherwise, in far less moral places, by no means is wont
to permit anything immoral.) The solemnity with which the espousal was
entered upon, the solemn manner in which a betrothal is still often
declared, knits both lovers more firmly to each other than before; what
they formerly desired, they have now thereby already provisionally got into
possession; and as their relation is now known and respected in the social
intercourse to which they belong, they can now also test themselves with
respect to this. Already before that first step, those who inwardly love
each other have not seldom a kind of anxiety, as if that which in their
hidden solitude was so sweet should now suffer by being given up to the not
always careful, sometimes disturbing influences of a crowd moved by so many
passions or moods. Therefore let one also grant young lovers, for a time,
the concealment under which the awakening love so gladly wants to dwell,
and so gladly dreams itself unknown. But since now a declaration must,
after all, happen---because lovers' shared life must after all once come
forth into the day, and as it were must venture into life's restless
tumult---let it yet happen first provisionally, before it happens finally,
or before it knits a marriage! The provisional bond is, after all, already
a bond; and when this kind of union lasts as long as it always ought to
last, it can become quite a good trial. It is better that the betrothal
become love's grave than that the marriage should become it.

There is many an egoistic nature which, desiring its love's object, in
passionate movement seems, in order to force it to itself, able to
sacrifice everything---seems to love with a fieriness, an energy and
firmness that promise the best; but yet soon after, when the secure
possession is won, shows itself otherwise. Sometimes it happens that such
an individual gradually grows lukewarm and becomes indifferent upon the
attainment, because it took it merely as another need's and delight's
satisfaction---thus as a transition, and not as the lasting and abiding
fulfilment of a privation lying deep in human existence. Often, on the
contrary, we see that the egoistic nature, which for a time was held down by
love's overwhelming power, comes up again and reaches about in love itself,
so that the individual gradually begins to plague the beloved with its
willful, imperious or self-willed mind. It is not all who can bear
happiness, bear the secure, calm possession of a good vehemently desired.
Even the mere thought that the bond is a bond plagues many a mind; indeed
it does happen that the one who showed great solicitude and attention
begins gradually to leave it off, merely because he is ensnared by the
thought that he is as it were bound to it, while he himself again begins to
regard as a right and a demand what he formerly, with joy and with love's
gratitude, merely took as a sweet proof of the returned love in which he
rejoiced so much. Indeed, so near to us lie the egoistic movements in our
inmost, so quick are they to come forth in the human heart, that even the
better individuality can sometimes find itself ensnared by similar
feelings---now discontented, now cold, indeed as it were frozen---almost the
instant after the fullest happiness and joy seemed to have set in. Such
things let one then know, and heed not the uncanny mood until it passes.
But if the natures are right egoistic, much strife easily arises, and both
collide the more strongly with each other, the more vehement the attraction
was. In all cases it is an anxious existence---an existence that can easily
become ruining and wasting for spirit and heart---to live under perpetual
strife and disharmony with each other, be the love even really vehement and
strong, and follow even upon every strife a love-reconciliation. Therefore
ought lovers, under a provisional union, first to learn whether they can
live with each other in the inseparable one. Precisely first when the
betrothal is entered upon, precisely when the desire has found its first
rest, can it show itself whether the relation has subsistence enough for the
next step to be ventured. He who, seeing that two do not belong together,
wishes to part them in the gentlest, mildest manner, acts sometimes wisely
in giving them together by a provisional union, so that they may now
themselves learn to get their eyes open and part themselves; whereas an
inclination struck down by a word of authority keeps its goad and does not
easily come to rest.

Meanwhile, when even only the provisional union is dissolved, this is
already a breach, and can make a deep and painful rent in existence. And
yet, if no attachment were to dare end with a re-separation, how could
quite reliable unions come about---unions that not merely the eyes, but the
hearts too, knit from the heart? Therefore let the intercourse between the
two sexes on the whole be quite free, and let there reign that
free-mindedness, that liberality in the way of thinking, which permits that
awakened hopes may come to nothing. For how are the feelings that can lead
to something to be able to get room to reach about, when the freedom does
not reign of stepping back again, however much one has approached? Many a
time it happened that the one who had come past the age in which what is to
happen happens quite of itself did not come to a union with another being,
merely because he would not venture to awaken uncertain hopes, and yet
could not come to give a secure and dependable one without exposing himself
to awaken the uncertain one first.

\begin{center}\rule{0.28\linewidth}{0.4pt}\end{center}

Those who for a whole life are to live with each other, to the pleasure of
God and men, must more than care for each other; they must also accord
together, they must inwardly understand each other and answer to each
other, and next also themselves watch that a good rightly granted them
becomes for them a true good.

Grounds why two who are each for themselves worthy and care for each other
yet do not belong together can, in individual cases, be many. In general
there are here three chief points to have in view: the difference of sex,
of age, and of station---insofar as this last brings with it a difference in
both parties' culture and whole spiritual mode of existence.

As regards the difference of sex, it is well to consider that the
man---however much his wife may otherwise be superior to him in gifts of
mind---must necessarily in certain definite respects have a definite
preponderance, if the wife is not precisely to feel herself unhappy under
the two's shared life. For precisely for the wife it is a misfortune to
have a husband whom she can and dare---indeed sometimes even must---look down
upon. Every deep, noble and finely feeling womanly heart will surely far
rather choose a tenth part of a truly manly, vigorous and spirited man's
love, if only she is completely and for good sure of this one tenth, than
all the undivided love and adoration that man can yield her whom she must
dominate in order to be able to live with him. There is something the woman
must have respect for in the man to whom she is to belong, be there ever so
much else in which he concedes her the advantage. What this something is,
is difficult to express in general words, however definitely one perceives
it in single given cases. Were we to express it, we might perhaps say that
it is his will, regarded as a whole as will, or as the personality's inmost
and most definite expression: it is his manly independence. A wife who
ceases to regard herself as the one who exists to serve her husband---indeed
even comes to a certain feeling of disregard or setting-aside toward
him---is a most unhappy wife. All that gives love and the love-relation the
depth and sweetness that seems, in a still far higher degree, allotted to
the woman than to the man is gone, and an unnatural relation has set in.
This every tolerably well-ordered and soundly feeling woman-heart feels.
Even the degree of influence which the wife most often not only has, but
also ought to have---and which precisely the most sedate and manly men most
easily concede---even this she does not gladly hear named or remarked; as
indeed it is not needful either that she be conscious of her influence in
the same degree in which she exercises it. A deep respect for the man, a
certain feeling of submissiveness, lies in every noble womanly being. The
finely feeling woman who heeds this feeling well perceives its significance
and its depth, and cultivates and nourishes it with delight. Clärchen's
situation at Egmont's feet, in Goethe, is therefore for a feeling
woman-heart a most attractive position; and a cultivated mind's most
glorious riches do not hinder, but much rather bring with them, that a
love-full womanly temper feels itself most deeply in soul made blessed by
being able to show a highly esteemed man how humble and submissive she
feels herself, from the heart, toward him. Yet these feelings are of the
tender and fine kind that they can come to speak only in rare, happy
moments. But in order that the woman may nourish these deep feelings toward
the man, she must, to be sure, not be all too superior to him in gifts of
mind. And here, then, lies one cause, springing from the difference of sex
alone, why two sometimes do not belong together, although each of them can
suit others quite well.

As regards the man, the first necessity for him is naturally that the woman
be, from the heart, prepared to be a woman---to which belongs being
domestic; for herein lies a great part of the woman's calling in the shared
life; she must be able to regard herself as a servant for her husband and
for the growing generation. For the rest, it goes without saying that here
regard is had only to that which is to be looked at in individuals of both
sexes as spouses. What is on the whole required in order to be upright and
worthy human beings, of that the talk is not here. Of imperious disposition
no human being must be, least of all a woman who is to be spouse and
mother.
Of differences of station there is only one that is so essential that a
union in which this difference is overlooked---leaving special exceptions
out of account (an outstanding individual, or a particularly inward
inclination, can bring a well-grounded exception)---becomes a true
mésalliance: it is that difference of station which is at the same time a
difference with respect to upbringing, cultivation, and the whole manner in
which life is regarded and enjoyed; a difference that brings with it that
the one cannot enter into the other's circle of ideas, that moreover
between the two no true spiritual communication can take place. Meanwhile,
this difference of station is only a purely human one, not a civil one. In
all civil stations there can be found the cultivated and the uncultivated,
nobler and less noble individualities. The true nobility, valid at all
times, which will also always maintain its influence in the State, every
human being of natural endowment can give himself, and every human being of
human sense can give at least his children, when these do not too much lack
endowment. All differences with respect to a merely civil station mean but
little where genuine love is found---especially in our time and our
fatherland, where one will hardly, save in single cases, regard a conjugal
union between two, of whom the one belongs to the highest classes in the
State, the other to a far lower, as a mésalliance, or as anything striking,
if only the two are not too unequal in spirit and cultivation. It is also,
not merely reasonable on the whole, but in particular wise, that the
highest stations use the freedom which the spirit of the age has brought,
to widen the circle for their choice, and sometimes to refresh their so
often rather withered stems with fresher saps, and therefore to profit from
the sounder human existence which yet still constantly seems chiefly to be
at home in the cultivated middle station, especially as regards precisely
the domestic life. Only princes must, in our time, still---as in so many
other human respects, so also with respect to a fair and pure love and the
union it strives after---suffer under their exalted position's many
inconveniences. No prince can yet easily act as a natural and competently
feeling human being in his family's circle; indeed, such as things have
stood, it is almost a wonder that we have in our time seen, in the royal
palaces, marriages as happy as anyone in the burgher-houses can wish them
for himself. It is, to be sure, in this respect, as in so many others, a
great good fortune that there is, after all, in Europe one land in which a
multitude of small princes have kept their princely dignity. Yet the time
will doubtless once come when a prince of clarity and power of mind will
find out the means of procuring for himself what every sound human being
needs, and of teaching the world that one need not pay the burden of being
a prince with the prerogative of being excluded from enjoying a pure human
existence. For the rest, it is quite natural if the lower station sometimes
shuns the union with the higher still more strongly than the higher is
sometimes unwilling to think itself in union with the lower. For the lower
here has more to venture. This holds especially when a girl of lower
station is to be united with a man of the highest; whereas it is less
precarious that a girl of high nobility be united with a man of the middle
station---provided, for the rest, that the two are mutually completely sure
of each other. For since the man draws the woman over to his place far more
than he is brought near to hers, in the latter case the removal from the
sounder middle position, for the one of the two who belonged to it, is not
so great as in the former.

Age's difference is still to be taken into consideration. One such
difference meets us here that is always precarious: that which brings with
it a difference with respect to the period in life in which the two stand
as regards love's feelings. When the one has already tested and enjoyed
love's whole fullness and glory, or has at any rate in one way or another
come so far that love can no longer, in him, have the youthful
enthusiasm's fieriness and liveliness, but the other individual, on the
contrary, still stands in life's first spring, then it can often, with
respect to this latter individual, be quite hard that the union is so
unequal from the two sides; for it then easily brings with it that the
younger does not come to find the youthful, lively, fresh, luxuriant
love-spring which he was after all determined to enjoy.

\begin{center}\rule{0.28\linewidth}{0.4pt}\end{center}

Although all love---and moreover love [Elskov] too---by its inner nature
awakens the spirit's higher eye and makes seeing, yet it can also often
blind, and make it so that the human being does not see what is to be seen,
because the loving one so fastens the eye on that which especially attracts
him that he overlooks the rest: a one-sidedness in which lovers often even
with delight---indeed as it were with diligence---set themselves fast, as if
it were a \textit{religio}, a holy necessity, for them, to find everything
beautiful and captivating in each other. But just as love on the whole
makes seeing, so it also accords completely with its nature to make seeing
in the particular. Precisely the deep solicitude for another being, the
inward delight in watching that this one may have everything for the best,
also opens the eye for the defects; wherefore, too, a genuinely loving and
solicitous mother is wont to have a sharp eye for her children's faults.
Love's inwardness and strength will not at once suffer by this cognition of
the beloved's defects; it will not suffer thereby, when precisely this
cognition is taken as something that heightens the deep and sweet feelings
of mildness and forbearance, and of a certain spiritual safekeeping and
solicitude. A vigorously feeling being does not even want that it should
seem, to the beloved's eyes, to be more or other than what it in truth is.
On sand one does not build even a hut, let alone a palace; and that is what
he would do who wants to build his whole life's happiness on an error.
Therefore it is precisely those who feel with genuine power and clarity to
whom it is of concern not to be more to their beloved than they are in
themselves. Truth is here, as everywhere, the only reliable thing; and to
be truthful both toward oneself and others is, like the fear of God, useful
everywhere. And precisely under love's youthful enthusiasm it is best that
both, with pure sincerity, see what is to be seen and will once be seen;
precisely then is it already time that the one be sincere for the other and
toward the other. For precisely then do both best learn to accommodate
themselves to each other's defects. Precisely then can they---if not
easiest, yet with least pain and displeasure---learn what they must learn:
that the glorious good bestowed on them in their love is an earthly good,
and that they, however natural it is that the one is as a god for the
other, must not too vehemently will that it really shall be so. To be sure,
no commandment is transgressed more easily and more often than the first
commandment; indeed, in love's first warmth it is transgressed even in a
fair, heart-refreshing and enrapturing manner, in that here, in a certain
respect, the one really becomes what so often the human being must be to
the human being: a god for the other. But let one yet guard oneself even
here! We see not seldom vehemently loving, but withal vehemently demanding
individuals, who with a strange confusion utter themselves and act as if
they willed that the beloved really should be what it seemed to them, a
superterrestrial being. Let no one stride to marriage with great
expectations or demands, but first learn rightly to grasp and remember that
we are all the Lord's debtors, and all together are on the way to his
judgment-seat, and all hope there to find grace. And let everyone who, on
the way, takes for himself an inseparable companion, take this one as it
is, with all defects, and hide these not from himself.

Next, let one regard marriage not merely as a happiness to be enjoyed, but
also as a calling to be satisfied. Especially is it important for the woman
to enter into marriage as an upright-feeling man enters into the calling
that the community of the State assigns him---whose servant he regards
himself to be, and for whose sake he considers himself to exist, while he
calls it a great good fortune to have found it, when it answers to his
heart's mind. What this calling in the community of the State is for the
man, that the married state brings the woman: the true sphere of activity,
determined by Nature, necessary to a complete existence. And let no woman
cease, when she now really steps into this sphere of activity, to regard
herself as that which the tenderly and deeply feeling beloved so gladly
dreams herself to be, the man's solicitous servant in all things. Let
neither of them forget, too, that love is constituted not merely by an
inward attraction, but also by that distance whereby both are held out from
each other, and to which corresponds that bashfulness and fineness which in
the union must be kept in honour. Let the feelings of esteem, and of the
attentiveness belonging to it, therefore preserve themselves and be
nourished! Therefore let the united ones not forget, either, that they have
duties toward each other! The word esteem and the word duty sound, to be
sure, to an inwardly and strongly feeling and loving being, often like
strange expressions, that seem to have little home, and to be without
significance, where everything, by the sweetest and most blessed feelings'
irresistible strength, seems to happen of itself---not merely without
resistance but also without reminder. It is also good, and precisely as it
should be, when---here as everywhere in the whole of life---a high, holy
life's unstaunched stream and breath so moves and animates the earthly
being's heart and its outward-going working, that now the whole life is of
itself lived as it ought to be, with full power and inward enjoyment. But
in our earthly existence on the whole, and likewise in marriage, not all
moments are God-filled or spirited and lively, even for a God-filled or
spirited individual. There come faint moments; there comes all kinds of
displeasure under life's unrest. Then it is good that in the
marriage-relation, as in life on the whole, the feeling of duty steps to,
guarding and preserving, so that the faint moments may not even pervert or
hold up the later, more spirited ones---indeed spoil the whole relation. The
human being must everywhere be watchful and guard his heart, and watch over
his beloved's too and bear solicitude for it. Spouses are, after all, to
bear life's burdens and sorrows with each other; indeed even these young
lovers look forward to with delight and courage, because these too are to
work toward knitting them more firmly and inwardly to each other; for to
those who fear the Lord, everything must come to good. But to that which
spouses are to bear with each other belong not merely life's weightier
sorrows, and the need and danger toward which a courageous mind precisely
feels itself called and enlivened to vigorous activity or endurance; there
belong also life's small teasings and vexations or tedium, by which a
temper only all too often lets itself be stirred up to a certain
willfulness or unruliness, which, when it is not watchful over itself, can
easily reach about. Those greater sorrows are borne more easily than these
trifles. Many a wife accommodates herself more easily to watching seven
nights by her husband's bed or by a child's cradle than she accommodates
herself to a needle's having gone astray at an inconvenient moment. But it
were yet too much if a lost needle should part those who have held out
weeks, indeed months, by each other's sickbeds. Therefore let them not
forget that they owe each other what each of them feels he owes the first
best human being who comes in at the door! But what they owe each other,
let them yield each other, then, as those who yield it out of inward and
watchful solicitude for each other! Perpetual small irritations can gnaw
more on life's and love's fullness than anyone believes in the moment. And
those who often irritate each other never know themselves secure. But one
place at least everyone must have, where he knows himself secure and feels
himself well, that he can always with delight hasten thither; and this one
place must be his home.

That the union to which love leads is not to be a union in this or that
respect alone, but in all respects, and that it is not to bring with it a
partial but a total shared life---this, in youth's fair period, every happy
lover feels; and while he himself strives to appropriate the fullness of
cognition and soul-enjoyment that has opened itself for him, and in which
an awakened and noble youthful mind rejoices so much, he desires at the
same time with all his might to live and enjoy the rich life with his
beloved. Both then find a thousand points of contact, and thousandfold
communication takes place between them, provided only they otherwise belong
together. But even those whose love fell in this fair youthful time---and
still more often those who first later found love's happiness---so often
forget, when the full union has set in and life's other claims and callings
present themselves again, the happy shape their youthful love had. All too
often even love is taken as a transition: the two united individuals step
gradually, in the spiritual respect, back into their former manner of
existing, and come, time after time, so far that both at last live, each
for himself, their isolated life side by side with each other, instead of
living one shared life, embracing both their whole lives, with each other.
Instead of their striving to enter into each other's individualities and
knit themselves, under manifold communication, to each other in many
points, the shared life restricts itself to a few points in which both feel
themselves as grown together; outside these it is as if they did not belong
together. Lovers are in many ways and in many respects an incitement to
each other. But to which side love comes chiefly to turn depends on how
they take it. And so it happens, then, that often great one-sidedness gains
the upper hand, and that whole regions in their shared life come to lie
waste. Often this state now sets in, to be sure, because both in themselves
really did not belong together, because the one cannot enter into the
other's circle of ideas; sometimes, on the contrary, from a certain
over-great passive contentment in the beginning, and from laziness---or at
any rate lack of attention and brisk power of mind---afterwards: a laziness
or inertia [Inertie] that even often leads to narrowing the judgment and
the free glance upon life, so that the lovers even lose the sense for the
manifold glorious that, different from that which they behold in each
other, moves about them, and now therefore lose the spur to their own
forward-striving and development. Meanwhile they then gladly give
themselves over to a certain kind of the movements in which love sets.
There are many whose whole love seems to stop short at the delight with
which they behold each other and make much of each other; and now, in order,
after all, to bring a certain manifoldness and variety into the relation
and drive away a feeling of tedium which caresses cannot always chase off,
they seek perpetual new arousals that turn about the same point, and the
spirit cannot over this come up. Herein many lovers sin, and forget that
even Eros's service must not be desecrated; forget that love is to lead to
a soul-communication and a spiritual relation; and now lose life's infinite
enjoyment over a highly narrowing one, which now does not even rightly
become an enjoyment from the heart's ground, and can easily become a
spirit-consuming one. Or else it happens that one lets life's thorns and
thistles gain so much the upper hand in the spiritual shared life that they
at last subdue the up-germinating seed. Sometimes, too, there was here to
blame an arrogance springing from the egoism ruinous in so many
respects---a self-willed and imperious delight---when the one who had given
itself wholly over was yet now too proud to want to give itself wholly
over. But not seldom even a shared life luxuriant from the first suffers,
and falls or withers away gradually, because a sufficient nourishment is
wanting, until at last the attention to the nourishment that might yet
present itself is also weakened. The surroundings then sometimes contribute
not a little to enlarging the cleft between the two. There is often shown
too little attention, or too little sense for marriage's holiness, and it
is not always considered what each one, in this respect too, owes to
others. The dearest people are not always welcome to us, and cannot be,
when they step disturbingly into our circle. To nourish a genuine
intercourse with each other there is required, among other things,
sometimes also solitude's calm. In a circle in which an extended social
shared life easily forms itself, not seldom the most excellent people
suffer most in this respect, because the most excellent soonest enclose
deep within themselves such feelings as, by coming forth into the day,
could easily be profaned.
Spouses on the whole all too often do not live enough with and for each
other. But in order that a full spiritual shared life may be able to take
place between them, the shared life must, to be sure, have an outer
substrate, a basis on which it can rest, and from which it can draw
nourishment. What is required for all life is also needful for this shared
life. But all life on the whole steps forth as living by stepping---working
and vigorous---into something other than itself, which other now becomes
the basis or substrate for its living utterances. Thus the soul, as living
principle, steps forth in the body, in that it forms, cultivates and
maintains this body, which meanwhile is yet something quite other than the
soul that constitutes the unbroken living therein. Thus, too, all that
livingly moves in the human being's spirit and heart---the whole fullness
of higher influences which, under the form of ideas and feelings, stir in
the human being's inmost---all this steps forth in the human being in an
aggregate of representations and states of mind, which are something quite
other than that inner fullness of spiritual living. The representations and
states of mind are the momentary, alternating and outer, which only is and
means something by that inner spiritual being which steps forth therein.
They are the substrate, the ground for the spiritual life, in which this,
by perpetually calling it forth, determining and cultivating it, shows its
inner vigour and essentiality; herein it steps forth as spiritual living.
And thus, then, the shared life too, in which two lovers in a fair and
inward union are to grow together, must have its substrate or its basis.
Not merely in the outward and civil, but also in the inward, purely
spiritual respect, marriage must have a ground for its subsistence, a
constant source of nourishment. There must be something about which both
parties' united striving, working and living so turns that it can become
for them a constant vehicle for the hearts' communications, and, as it
were, a material by which life's inner power and content is perpetually
awakened to utter itself. The human being is, in the spiritual respect, to
grow the whole life through. This growth happens as life's inner power
perpetually more and more cultivates the circle of ideas and the aggregate
of feelings and resolutions that the human being bears and preserves in his
inmost. But lovers are to grow with each other. They must, therefore, in
union with each other, cultivate the circles of ideas in which each of them
expresses and enjoys his inner life, so that thereby both parties'
spiritual life comes in many ways to reach into each other.

Happiest in this respect is the position of those whose calling attaches
itself so immediately to their domestic life that the wife must immediately
share the man's whole sphere of activity and his interest and work
springing from it. To be praised in this respect is especially country
life, in which the man's whole working and the wife's whole working reach
into each other. It is unspeakably sweet to work in the company of a
beloved being. But this happiness, which in the lower station is often had,
is in the higher far rarer. Here the man's calling most often lies in a
quite other sphere than his domestic existence; both spouses can indeed
work for, but not with, each other. Meanwhile it is everywhere natural to
share with the beloved all that livingly moves and fills the heart. Even
if, therefore, the man's calling and occupation lies even far from the
circle of ideas in which the woman's life alone is wont to move, it is yet
natural that he desires to share with her the interest and the feelings
that in this calling move and fill him; just as she too must desire to take
part in it, in order to be able to rejoice in what gladdens him, and
inwardly to enjoy this joy with him, and to be able to enter into his cares
and grasp them to the full. So let nothing, then, hold the lovers back from
living thus with each other, and sharing all that reaches into the life
each of them lives! Even if the woman is little fit for certain manly
occupations, she can yet well be able to grasp their worth and
significance, to understand and with living interest regard their purpose.
On the whole there is in the womanly nature's idiosyncrasy nothing that
hinders the woman from having a sense for all that---at least on the
whole---which the man is able to cognize, but which he, to be sure, must at
the same time be fit to think out and carry out. Spouses' shared life can,
then, so far, very well be a genuine one, embracing both parties' whole
lives.

Next after communion with God, and a union with other human beings for
whom a human being can open his heart, there is hardly, for the one who has
come to a certain spiritual maturity, any spiritual necessity more
essential and more pressing than this: to have a definite sphere of
activity, a calling, that answers to his peculiar nature and inclination,
and which he is able to fill out. The woman is in this respect exceedingly
happy; for when she first finds the one to whom she can give her whole
heart, she finds, at once and in one, fulfilment for two of the most
essential necessities. The union to which love leads opens at the same time
the fairest sphere of activity for her---that to which she is precisely most
of all called by Nature, that which is able in the most refreshing manner
to fill out mind and heart. Something, as Schiller aptly lets the chorus
sing in The Bride of Messina, something the human being must have to care
for from one day to the next, so that he may endure the long life. But
surely there is, among this caring and busying and working so needful to
the human being, not much that, when it is as it can be, is fairer and more
refreshing than a woman's unbroken solicitude that a beloved man and dear
children may have everything for the best. Life is on the whole in a higher
degree a genuine living, the more it is a genuine loving; but to love is to
give oneself over for something other than oneself, and to enjoy life and
rejoice in life in the devotion that turns about the beloved object.
Therefore a genuine living is always a genuine service of love, an inward
solicitude that that to whose service the human being has given himself
over may enjoy its right and its fulfilment and be at its best. Thus every
proper man loves his calling and feels himself in its service with delight,
and lets it be of heartfelt concern to him to satisfy it, enjoying his own
personal life in the faithful and hearty fulfilment of his calling. The
simile that a man is wedded to his study and his calling---that it is to
him as a bride whom he inwardly loves and bears solicitude for---is a simile
that is of significance, and that expresses the man's relation to his
sphere of activity and to the idea by which he is led, and for which he
works, as this relation should be, but as it, for many reasons, to be sure,
is not always. But how fair is especially the woman's position in this
respect? What unity, steadiness and energy can her calling not give her
life? What the man so often must live and possess and enjoy in two
different circles, that she lives and possesses and enjoys in one and the
same circle. Precisely while she lives for her calling and lives for the
idea and serves the Lord, she lives immediately at the same time for the
beings, and to the service of those beings, who are nearest and dearest to
her heart. To her it is given, precisely while she fulfils her nearest
determination, as part of the whole, to be able to concentrate her working
about the one circle to which she has given her heart. But, to be sure, the
same thing whereby the woman's lot can be so fortunate a one also makes, at
the same time, that her position is in another respect more precarious; for
just as both so essential spiritual necessities find their fulfilment with
each other and in one, so both must also be dispensed with together; and
the woman-heart that is not so happy as to be able to give itself wholly
over to a beloved man must now at the same time forsake the calling
determined for her by Nature: a determination of nature that can only half
be replaced by anything else---be it, then, by a quite decided talent for
one of the activities that reach into the whole of human society.

That the woman lives to serve---and to serve personally, and to the
satisfaction of others' personal necessities---and that this service is a
fair service, and that a doubly fair one when it turns immediately about
highly esteemed men's personalities: this every awakened and fiery womanly
being easily feels, in whom the natural human vitality [Menneskelivvær] has
not suffered too much under a perverted upbringing's or an unfortunate
station-relation's harmful influence. Young women of sense and feeling find
great joy in personal services, and learn also gradually, without
difficulty, to cognize that---since only a genuine working is a genuine,
enjoyment-full living---it is a good fortune that their working can and
shall turn immediately about a beloved circle's personal well-being, and
that the feelings of esteem and devotion have so much right and so much
occasion to utter themselves at once in action, under personal solicitude
for their object. The Goethean saying, ``Let the woman learn to serve,'' is
easily complied with; at least the soundly feeling woman complies with it,
on the whole, with delight. But then let there present itself, in the
conjugal union, also the other thing he adds! precisely by serving, let her
come to the disposal that belongs to her in the house! No one rules better,
and comes better to a dominion, than the one who is first himself livingly
and vigorously governed by a true idea or by a pure feeling; but let the
one who is animated by this then also enjoy the disposal and the freedom in
his circle! Above all, let the wife enjoy it then in the house! Every
soundly feeling wife desires this, when she is first conscious of her
competence. And rightly. She can only serve to the full when she can
dispose to the full. And if the true feeling is there---the genuine, deep
love that makes watchful and makes seeing and makes conscientious---then for
the man too nothing is more natural and convenient than this. But let one
then not forget that, since it is only marriage that opens to the woman
that which the man has already found elsewhere, the peculiar sphere of
activity, so let the woman be granted what the man in his sphere so gladly
wishes to find: right and freedom, in carrying out what is to happen (while
a generally valid order therein is uprightly and faithfully observed), yet
to do it in her way, and that rather now and then less conveniently, before
she rightly gets herself formed to her calling, than that the developing
sense and idiosyncrasy should be cowed. But then let the woman also, just
as little as the man, forget that---while it may indeed with reason be
called an earthly heaven-kingdom, what is prepared for a happy couple by a
genuine and youthful love, when they take the happiness granted them in
such a way that they themselves let it become a heaven for them---this
earthly heaven-kingdom is yet by no means an earth-kingdom. Just as Christ,
for those who with all their might would see in him an earthly king, did
not become the heavenly king he was, so all too easily can the fairest and
most heavenly in the love-relation become perverted or confused and
disturbed for those who take the union as if it were earthly glory they
were now to dress themselves up with. It is so far from being the case that
earthly abundance and pomp should further a beginning conjugal union's
blessedness, or heighten the forming shared life's allurements, that on the
contrary a certain restriction in circumstances---when only the way to
acquire is free and open---best accords with the newly knit couple's whole
existence, or, as one calls it, best suits them. Precisely this restriction
makes the beginning shared life so cosy. Whereas a quite complete and
abundant domestic establishment and possession, into which they are set at
once, finding everything already ordered for them, robs them of the fair
joy of---by shared striving and shared deliberation---time after time
calling forth the whole domestic surroundings, and cultivating and
completing it with an interest, solicitude and love attaching itself to
each single part of it. Riches and abundance is just not always a blessing.
Whereby it is also to be remembered what Goethe, in Wilhelm Meister's fifth
book, says, in drawing attention to the wrongness in the festivities with
which conjugal unions are begun and, as it were, consecrated: ``there is no
feast that ought more to be solemnized in humility, stillness and hope than
the wedding-feast.'' And yet, how often it shows itself here too how petty
selfishness or petty considerations are ready to come forth at every
occasion and to govern human beings. Especially let the woman here be
careful, and not cease to be love-full and humble from the heart because
she has become happy. It is indeed natural enough that the womanly being
who, while she still saw the coming happiness in hope, perhaps showed great
self-restraint---that this same being, when the fullness of happiness has
come, can no longer so well preserve the equilibrium. Happy love makes the
woman free and authoritative; it makes it so that she, more than before,
comes to feel herself as one who has a personality, has worth and
significance, and moreover has a will. It is indeed also most often good if
the young heart, in which now a world of feelings comes forth and works
itself forth, gets room and freedom right from the heart's ground to break
loose for a while. Let no one demand that what today still goes fermenting
up and down in the mind shall have settled to clarity and calm tomorrow.
Meanwhile let the woman guard herself: first against the self-will, the
utterances of an imperious mind, which now easily stir and get a desire to
come up; next let her guard herself against the unsteady. The best shelter
here is the nearest, that which is especially able to bring the mind to
collect itself with power: the hearty and faithful preparation for a
calling that really is a calling. At the same time also a watchful delight
in living, right from the heart, to the beloved's pleasure and joy. On the
whole let every loving one consider---what yet so many often seem to
forget---that returned love is not preserved in any other manner than that
in which it was acquired, and that one cannot easily continue to be loved
when one ceases to be lovable. Let no one regard love's promise as a debt;
merely as a debt we do not, after all, wish that love should be yielded to
us. Let no one be vexed---on the contrary let it be dear to everyone---that
the love with which we are loved has its ground in a higher love, and that
we therefore must watch that it may constantly be nourished thereby, that
together with it that higher love too finds nourishment; so that we, then,
in order to preserve for ourselves the love that was bestowed on us, must
preserve and nourish in ourselves that which makes the human being
captivating for the human being; and not merely preserve it, but also let
it step forth and utter itself so that it can rightly work beneficently
outward. To be loved with a blind love is not to the profit of anyone. But
to love and be loved in such a way that a pure love for the true and good
and for the beautiful, and a living desire to express it in one's person
and in one's life, can be confirmed and furthered by our love---that is to
love with great advantage.

Yet the pettiest utterances of an egoistic delight lie not seldom---one
would hardly believe how near---by love's otherwise egoism-overwhelming
feelings. All too often people utter themselves as if---when they have
scarcely come so far as to have attained a blessedness that should make
them willing with delight to let go of much that might otherwise be
attractive to them---as if now, called forth by this happiness, all petty
appetites and desires swarmed forth from all the heart's corners, and tried
to be furthered and to obtain some nourishment by the new life that has now
opened itself. And what such appetites, what all kinds of vanity, do to
ensnare and confuse, that innumerable considerations of fashion, and of
what most people are wont to do, or of the judgment of these and those, do;
so that it is not seldom difficult enough to come to enjoy a
straightforward, natural, sound life in one's own way. It even goes so far
that people, out of all kinds of foreign considerations or appetites,
undermine that which is, after all, the basis for their happiness---a secure
and reliable domestic establishment with respect to livelihood---while they,
precisely thereby, instead of getting any genuine enjoyment from it, spoil
the sweetest and most refreshing thing in their shared life: the tender,
deep solicitude for each other, and the cosy withdrawnness and love-full
restriction that first makes the shared life a true idyllic life, something
whereby marriage in its beginning is so much beautified.
\begin{center}\rule{0.28\linewidth}{0.4pt}\end{center}

But in the worst, the most spirit-consuming ways can the mind---ensnared
and entrapped in itself and in its considerations, the mind that everywhere
works so ruinously---disturb and destroy a conjugal union, when it leads to
sheer irritations, especially to jealousy's plaguing unrest. In fiery
people's first youthful love the painful, life-darkening and disquieting
irritations come up only all too easily and often.

It is not at once and easily completed, what is to happen when two
individualities---which each for itself has yet always to a certain degree
fixed itself and taken on a definite shape---are to be loosened in order to
grow together into a new life uniting both. A certain struggle between the
two sets in, which in vehement souls can easily lead to all that
alternation of calm and unrest, gladness and pain, hope and anxiety, in
which youthful love nonetheless so frequently comes to utter itself. It is
in many respects a fermenting movement in which love sets the temper. Love
is indeed---though it be a passion in the higher sense of the word, in which
every rapture, even the genuinely genial, is so called, as one also speaks
of a divine passion---yet by no means, according to its inner being,
necessarily a proper passion in the word's usual lower sense; for it by no
means lies in its nature that the spirit's free life and sound
forward-striving should suffer by it; there by no means belongs to its
essence so one-sided and vehement a desire that it should, quite
captivating the individual, becloud the mind, or at any rate make it
unclear, and, in strife with other feelings, abolish the inner secure calm
and equilibrium in the spirit. But in its youthful beginning fiery love is
yet almost always passionate, and can then, so far, bring something morbid
with it---though, to be sure, in otherwise sound beings loving with genuine
power, no other morbidity than what any of the crises can produce which
belong to a developing life's natural course. Fiery love naturally carries
one away for a time in such a way that the individual, torn out of its
former track, swarms about in the new region, without firm steadiness with
respect to life's other many movements, indeed even without sure and secure
power with respect to love's own utterances. The soul's wings are indeed
loosened and unfold in the rising love's warming sun, but they are also now
easily entangled and ensnared, as they vigorously strive to move.
Especially is it here to be noted how the same feelings which with so much
power can tear the human being out of a selfish existence's narrow circle
yet, in another respect, so easily stir up and hasten the individual's
desire to maintain its peculiar subsistence and let its personal
idiosyncrasy enjoy furtherance. Since love is not merely an infinite
devotion, but also itself draws vehemently to itself and desires with all
its might to meet the same devotion, it is natural that the loving one,
seeking and expecting this devotion in the other, now also makes his
personality count. A certain perpetually recurring mutual resistance, a
certain resistance, sets in here naturally under the advancing
growing-together; indeed it belongs to the relation's cultivation, so that
the harmony---instead of becoming a genuine accord, which answers to both
parties' mind and peculiar nature and yields them both true fulfilment---may
not spring from faint compliance, and become more an outer than an inner,
more an apparent than a real one. The inner emotions and endeavours
revolving about one's own I can even not rest altogether, since life only
forms itself and reaches about in and with a self-asserting personality.
Even in order to enjoy one's happiness to the full, and rightly to enter
into it with delight and power, the own self must come up in the soul. The
new life can only fully come to the individual's good when it feels its
inmost I, its most peculiar personal nature, thereby clarified, raised and
furthered, by means of a true and deep accord between its own heart and the
new, self-opening life. Let the union with a beloved and adored being set
no one in discord with himself and his nature! Therefore let the compliance
and yielding, which to be sure in many a point must set in, be no weak and
heartless one, but a hearty and love-full one, that yields because the
individual from the heart wills it---from the heart wills to let its love
have this power with it, in that it does not feel anything true and genuine
in its being thereby cowed, but precisely in the yielding feels refreshment
for its mind, and willingly gives itself over to that loosening in the
mind, without which yet the own self, the freer soul, cannot come to feel
itself raised and widened by the penetrating love's power. But not at once
or easily are found and held fast, under the individualities' now often
ensuing collisions, the right fixed points for the so movable moods and
feelings---the firm points about which the shared life can form itself into
a true harmony and a firm and genuine subsistence. Much uncertainty,
wavering, unsteadiness and blind one-sidedness can now come forth, and it
shows itself often quite plainly that Eros is not of heavenly origin alone.
Since the blessedness the heart has found is always so far only an earthly
one, as its possession is not exalted above all upheaval, there often plays,
under the feelings' alternation---now a certain anxiety, now a certain
tempting desire to feel and behold one's influence confirmed, and to assure
oneself of one's happiness's unchanged keeping---so into the shared life
that thereby many an irritation arises, which gradually easily leads to that
alternation of irritations and reconciliations whose precariousness we have
already in the foregoing come to speak of: an often highly spirit-consuming
alternation, killing the deep calm and security, of tensions and strainings
of spirit and returning mildenings and meltings. Even into a sheer play
with the feelings this state often degenerates, so that it is as if the
lovers---though half unwittingly---call forth the frequent love-quarrels
with a kind of delight, or from a kind of bad habit, merely in order, under
that alternation of feelings, to feel and enjoy their love right strongly,
and thereby to bring a certain movement into a shared life which, from lack
of genuine communication and vigorous liveliness and a comprehensive mutual
interest, can sometimes become drowsy, indeed tedious, and quite come to a
standstill. But often the ground of such perpetual irritations lies deeper.
Not seldom in a disharmony originally grounded in the individualities,
between tempers which in themselves cannot well accord together because they
do not answer to each other, however much both feel themselves really in
love, drawn to each other by the power of certain allurements. And even if
there be no such disharmony, yet a certain strife between the selfish
emotions and the sympathetic, between the I and the devotion in which it is
to lose itself in another being, can persist long before the individual
comes to clarity. Indeed, as we have already before had occasion to remark,
precisely the feeling of the love-bond as a bond can here become ensnaring
and conditioning, and thereby tempting and stirring. Both then
become---while they are in a high degree an incitement to each other for the
strongest love-feelings' perpetual enlivening and up-surging---yet now at
the same time also an incitement to each other for sheer embitterment;
often, right as if they also needed this, and were precisely also by this
incitement drawn to each other: a state, however, in which there is no
refreshment and comfort. Only all too often lovers forget in time to attend
to how they stand with each other, to enter with attention and interest
into each other's natures, and now, telling themselves uprightly what they
want, to let fall what must fall, and draw forth what is to be drawn forth,
and hold fast that which must be held fast, so that confusion and unpeace
may not gain the upper hand, either in one's own inmost or in the shared
life. Yet, to be sure, with youthful lovers not every love-quarrel, every
embitterment of the kind that love alone makes possible, can fail to
appear---not even easily every groundless one. The ardour would not even
seem great enough if it could not lead to them. Meanwhile let one always
guard oneself against a play of the feelings which, while it gnaws on the
power of spirit, can work toward constantly more fixing---indeed
enlarging---the cleft between the two, who now, instead of growing together,
let the sharp edges by which they collide with each other be fixed more and
more.

We meet here jealousy in the first of its shapes, as an ardour rising to
passionateness, ensnaring and entrapping the mind, in the maintenance of
one's own I in relation to, and in opposition to, the beloved---a
passionateness which, however, presupposes love, and only in and by this
takes on the stamp it has, perpetually anew nourished and kindled under the
inclination's strong feelings. For we do not here speak of the mere
self-willed and imperious disposition, which, with a setting-aside of the
other---as if this other were a mere means---coldly and frozenly will only
have its will carried through, and does not ask about accord. The eager
struggle for one's own self which often sets in under love's most vehement
feelings, and which kindles strife and heat, precisely wills and demands
accord, and that an accord of love; it demands the other's undivided
attention and readiness and compliance everywhere and in all things, as
sign and pledge of love's unbroken keeping and unweakened strength.
Precisely with love's doubt and torments, anxieties and embitterments, does
the agitated mind here tear itself asunder. The selfish emotions, so easily
coming forth in all human beings, so ready and quick to utter themselves,
cast themselves on the love, and seize it as that in which they set
themselves fast, and now rage with double strength, instead of love's
getting its power and influence, liberating from petty egoism, in the
loving one's heart. The egoistic tenacity upon oneself and upon one's own
worth, influence and significance, as one that now perpetually lies in the
individual's thoughts because it is its own---this now embraces the whole
love too. Egoism's devil rages with double strength. The unclean spirit
which for a time had left the human being, and had gone about in deserts
and in dry places and had sought rest and not found it, now returns, as it
stands written, back again, and finds the house it left swept and adorned,
and goes and takes seven spirits with it worse than itself, and enters in
and dwells there again; and the last becomes, for that same human being,
worse than the first. The holy and beautiful in love now itself becomes
nourishment for the malady gaining the upper hand. Precisely the
unspeakable tenderness, the boundless devotion with which one heart sinks
wholly, out of love, to another---precisely this now presents itself to the
confused mind in such a way that just the thought of it, which seems with
unweakened strength in every moment to have to show itself and come to meet
the demand---indeed not even to let it come to a demand---spurs to the many
resentful movements, to the perpetual doubts and embitterments, and seems
not merely to justify them, but even as it were to hallow them. So the
temper confirms itself even---right as if it were with diligence and
zeal---in the almost raging inflammation and vehemence which gradually, as
it spoils love's own beauty and lays the mind waste, now also spoils and
destroys the whole shared life. Precisely the individual's happiness leads
it into the unhappy condition; and the fairer the heaven that seemed
prepared for the happy one showed itself to be, the more vehemently is
kindled, at the thought of it, the terrible hell. To wrath and the most
hostile sensations, to hatred and cruelty, can love now lead, when at last
the adored one seems no longer to be worthy to live. At every occasion
there is now stirred up anew the whole pent-up resentment in the soul that
perpetually looks back on its significance and influence, unsteadily cast
to and fro---the soul that is always in such a mood as if it were never
secure and sure of its cause, but perpetually asked itself and perpetually
wanted to find out whether it still counts, with the beloved too, for what
it is to and will count for. And in what condition is the beloved
individual now set? If this one is of mild and yielding disposition,
inclined to calm, it is easily cowed, or at any rate easily becomes so
careful and timid that the shared life's freshness, liveliness and
merriment thereby suffers---indeed that it at last comes to a standstill,
or, if movement remain in it, only moves anxiously under perpetual
tensions. To such unlovingness comes at last the vehement jealous love:
that it itself works to strike down and ruin what it precisely should have
of concern to keep in honour---what precisely perhaps chiefly makes the
beloved's person captivating and attractive, indeed perhaps even hastened
love itself. All of which also sets in; but to which yet many a vehement
strife, wrath and bitterness, unrest and unpeace, and moreover many a
destructive disorder in the tempers and in the shared life, comes, when the
beloved individual has too much firmness of character and strength of
spirit to let itself be cowed down. A mere appendage to another's
individuality---as if it were set in motion only by the other's life and
soul, quite formed off and fitted after this other---such a mere addition to
the life in which another being moves, without peculiar, free, vigorous
life-utterances, an energetically feeling being cannot endure to be, and
will not be. It cannot let all considerations perish in the one
consideration for the beloved; it cannot let all its life-movements be
captivated by the beloved, which would then be an entrapping and morbid
one. A god, for whom everything possible should stand back, everything be
sacrificed, no human being---not even the most highly loved and
esteemed---can become for another, however much it, in a both foolish and
presumptuous perverseness, behaves as if it would be it, indeed thinks it
ought to be it.

We thus see---still without a third party's disquieting stepping-in working
along, without any regard to a surrounding world---a jealous one's
unsteadily circling fire burning in a temper that does not let love's
devotion get power to overcome and quench the so easily stirred-up selfish
emotions. And this disturbing vehemence is quite nearly akin to the more
proper so-called jealousy. By this---by the passionately confused and blind
jealousy over the beloved, in opposition to, and under an ensuing relation
to, a third person---there is only widened the sphere in which the same
feelings rage. Only the occasions change. Here too there reveals itself,
where there is not in truth danger and ground for watchfulness, a mind
ensnared in its own self, never forgetting its I, but perpetually fearing
and yearning for it. Only that here it is a third person against whom, at
the same time, jealousy's vehemence is kindled, as soon as the beloved
shows obligingness and a loving heart toward such a third, or finds
pleasure in this third, who now becomes the occasion for unrest and
displeasure, and for perpetual new doubts about the returned love's full
strength and dependableness. Here too passion works against itself, and
against the best and fairest in the whole love-relation.
Is not freedom that element of love without which it cannot come up and
stir as true, deep and inward love? Indeed, is not that which the
lover---vehement and burning even to jealousy---wants to meet and know
himself to possess, precisely a free love from the beloved's side, one
called forth by the heart's irresistible inclination alone, and therefore
precisely also unassailable? No lover, but least of all the one so eager to
this degree, is content with mere love-utterances called forth by a feeling
of duty; an immediate devotion, springing from the heart's free movement
alone, he desires; and the thought that what is bestowed on him is yielded
him only on the strength of a representation of duty would precisely
irritate and outrage him. But what perverseness, then, to annihilate or
work against the necessary condition for that which he desires---the
condition without which love cannot be the one he wants, or without which
he at least can never know whether it is such a one, namely a wholly free
one? His feelings' natural strength and consequence should lead him to will
that everything compelling be removed, even the very least thought of duty
and right, as if here there were neither an obligation nor a demand. Every
compulsion, every reminder of a demand and a must, would indeed strive
against the wish to be loved from free inclination alone. This even the
human being loving with morbid, passionate vehemence often feels himself;
and, confessing that it has nothing to demand, it is only against itself
that it will rage, full of bitterness and cruelly, when jealousy assails
it. But even then it does after all set the beloved in anxiety, and lays
even then painful and constricting bonds, and drives away the freedom which
alone can assure it the pure, secure and total feeling that it really
possesses, in the beloved's heart, a heart giving itself wholly over with
free delight. Precisely love's natural consequence should lead to the
genuine sense of freedom. To transform all that is a necessity into a
matter of freedom, and let the spirit and freedom of spirit get room
everywhere---that does he who has sense and mind and gift to live life in
the genuine manner; but even to take that which is a matter of freedom, and
can only have worth by its freedom, in such a way that it is treated as an
object of binding necessity---that is a great folly. And yet, how often do
human beings draw that in which the freest heartiness would reveal itself
down into a conditioning demand's sphere, and let (to apply a passage of
Goethe in his \textit{Aus meinem Leben}) each other atone enough for having
once been so happy as to be able to be a joy to each other.

And how little, too, is that morbid love, in many another respect, a
genuine love. Not to speak of its being so often already rooted in the
sensuous, in a low and unworthy manner: even where it, as a deeper feeling,
as proper loving, is directed at the spiritual, it has, though in root and
origin a genuine and energetic one, yet not so penetrated through that it
has come to unfold itself with a true love's all-sided and deep power. How
little does the perpetual struggling and yearning, the doubting and
forsaking, the perpetual reflection on itself and on its possession as its
own---right as if the maintenance of one's own I's significance could never
be forgotten---how little does it permit the sweetest feelings in love to
come up and to work with comprehensive strength! The sweet feelings of
true, self-forgetting solicitude and tenderness, of mildness and
forbearance, do not come to work. The ensnared soul does not even come
rightly to behold or to enjoy what is given it. That which should be to the
beloved a true joy, that which precisely with delight should be called
forth, maintained and nourished---the beloved's open sense for life, and
the heartiness with which a pure and free temper, a happy soul, gladly
comes to meet every warm human heart and opens itself to it---this now
becomes, for the unhappy passion, a never-ceasing source of pain and
unrest. What power has not a true and unfailing love---precisely the love
everyone desires to find---to open and as it were loosen and free the mind,
so that it now with fresh liveliness and hearty merriment gives itself over
to every fair impression, to every soulful life-movement, and to every
awakening feeling of the universal love for all human beings. But all this
the untamed passion works, in blind delusion, to stop and crowd out---or at
any rate to weaken and displace---in the beloved, so that this one's beset
temper at last, out of anxiety, dare not utter itself and widen itself with
freedom. On a desert island the human being living under jealousy's
vehemence would like to live with the beloved, shut out from all the world,
so that the two alone could altogether be enough for each other, and that
above all neither of them should find another, or a freer and more widened
heaven, than that which the passionate soul then thinks to prepare for
itself and its beloved: a heaven that in truth would become narrow and poor
enough, if any could come forth at all, and jealousy would not rather rage
all the same under that first shape of it, for which no third party's
stepping-in is needed. So presumptuous is the human being ensnared in this
derangement of mind, and so ready to ruin the happiest and most captivating
nature in the beloved. And now it at the same time, in its own loving, lets
perish that which is otherwise the fairest thing in love: the inward,
pure---and therefore calm and secure, mind-widening and freeing---unspeakable
trust, the pure consequence of a true and full devotion, the immediate
assurance reposing in the pure loving alone. So far can jealousy's strange
derangements go, that something most base can step in place of the highest
and most glorious, and that the loving one, lost in a self-lacerating
madness, would rather live in trust in eunuchs than in the deep,
heart-refreshing assurance, the infinite trust in the beloved's hearty love
and fidelity. And what, after all, is a love worth in which these deep
feelings of trust do not live? How profaned here becomes the enjoyment of
the holy! In truth, if love is to lead to mistrust's---or the perpetually
stirred-up doubt's---unceasing pain and plague, and to the never-resting
anxious self-regards, and if thereby the sweetest, noblest, freest and best
in love is to be constantly cowed and never come to animate the heart with
full power, then it is after all more sound not to love. Better than to
suffer and be ensnared so much under one's love's blessedness is it, after
all, to cleanse oneself of such a spirit-pain by tearing its root out of
the heart---the vehement passion from which all this suffering of mind
comes---or at any rate to let it grow lukewarm so far that a more bearable
existence, precisely because it is lukewarm, can set in. Indeed, the one
who suffers from that spirit-consuming malady is perhaps nearest to his
blessedness when he can bring it so far that he forsakes the vehement
desire with a religious disposition, and, confessing himself weak, is ready
to accommodate himself to the most painful; for then perhaps the evil demon
may soonest give way, so that the Lord's power can show itself strong in
the weak. Clear cognition and calm contemplation---whereby the unnatural in
jealousy, the destructive, and that which altogether strives against true
love's essence in it, is grasped and seen from all sides---must work on the
one side; but on the other side, the compulsion, so often needful
nonetheless under the spirit's childhood here on earth, which the
individual lays upon itself. Do you suffer so much under your painful love,
under the perpetually returning doubts and thoughts that grant you no calm
and rest, then resolve within yourself that now you will---however
vehemently you love---nonetheless accommodate yourself, if it should be so,
to the returned love in your beloved not being stronger or more secure than
in the bitter moments it appears to you. Try once what a devout resignation
can accomplish! If everything stands in the Lord's hand, then let your
love's happiness too stand in it. Perhaps it will then also be fulfilled in
you, that he who rightly understands how to lose himself shall win himself.
Genuine, resolute surrender to God's will is everywhere of wonderful power.
In the highest love lies that which can give all other love rest and peace.
But here too let one at the same time consider that what is to come forth
and grow in the spirit's life---just as much as what grows in the outer
organic world---attains its needful strength only gradually. Let no one
expect that an evil which has got power in his inmost always gives way at
once because he condemns it. It can come back many times before it
altogether sinks down. Therefore let the one who has something to work on
in his inmost not become impatient or dispirited, but let untiring striving
toward the better work calmly and slowly!

The terrible state to which jealousy in its highest degree can lead is just
not so frequent. But emotions that play into it---feelings that as it were
taste of the tormenting pain and of the unruly wildness, up-surgings in
which the quiet beginnings of that soul-confusion and that bitterness
stir---are not so rare, and lie doubtless near enough to most souls loving
with youthful fieriness to be easily awakened. It shows itself here, what
is so often seen, that the attained possession of an earthly good
vehemently desired is just not always at once a blessing of heaven.
Precisely as an earthly good, love's happiness lies so near those loveless
emotions; for as an earthly good it is an object of exclusive possession,
therefore also of danger and fear---indeed of burning yearning and chasing,
of doubt and anxiety, when it never seems secure and dependable enough, or
never high, pure and divine enough. It shows itself here how dangerous it
is to have other gods besides the Lord. Precisely that---so wholly to set
one and all in one's love, or to give it so high a place and significance
that everything else in life seems to want to be consumed by it, and now to
drive it to an idolizing that does not leave the heart room for deep love
to God and for hearty devoutness, as if the beloved were divinity
enough---precisely such an excessive adoration of the beloved easily leads
to one's now, willing and demanding the same of this one, also setting
oneself up as a god, and never seeming to find undivided love and attention
enough. So high, on the whole, must love not be set, so holy must it not be
esteemed, that it should lead to stopping the inner life's free development
and fettering the spirit that is determined to ripen, on earth, toward an
eternity---the spirit whose god no one can be but the true God, to whom the
whole life is to strive to come.

\begin{center}\rule{0.28\linewidth}{0.4pt}\end{center}

%% ============================================================
%%  THIRD BOOK
%% ============================================================
\chapter*{Third Book.}
\markboth{Third Book}{Third Book}

In truth Eros may well be said to be of mixed origin and nature, and not
all have reason to praise the mighty demon. Unspeakably beneficent can love
become, when in an inward union under fortunate circumstances it knits
together two tempers that, in fair harmony, know how to enjoy from the
heart's ground the happiness granted them. But how often has love worked to
ensnare the spirit's power and forward-striving life, and to hold the best
in the human being captive in unrest and confusion. Even the most spirited
and lovable beings sometimes did not come to live a fair, spirited and
enjoyment-rich life with each other, merely because an unhappy
falling-in-love stepped disturbingly into their relation. A pure
contemplation, filling and refreshing the soul, of a glorious
individuality's depth and fullness; an inward spiritual communication and
participation; sweet confidingness, trust and devotion; and the full
assurance of possessing, in the other's heart, a dependable refuge under
all life's happenings: what sources of unspeakable joy in a content-rich
shared life. And yet, how often did these fair goods not come to the good
of the human being for whom they daily stood open, or became for him only
an impure, unclear and half enjoyment, merely because they called forth
love's yearning, burning desire, with an unquiet mind's stormy movements!
Be it now---for in many ways love can cause confusion and unpeace---be it
that it stirred up the egoistic in the hearts, and jealousy's vehemence and
ardour; or that it drew two individuals to each other who could come to
fulfilment for a desire only on sorrowful and misleading ways---a desire
which, awakened and nourished but unsatisfied, perpetually disturbed and
disquieted; or that it was kindled only in one of the two hearts, and now
cast this into a hopeless love's spirit-consuming alternation of passionate
moods, and into the perpetually renewed uprising of the unsteadily
tossed-about feelings. Even so far can it come with love's errings, that
it---called forth by certain allurements' inciting power---knits together
natures which properly and at bottom do not harmonize at all; indeed that
it captivates a lovable being to an utterly unworthy one. Add to this,
further, that love sometimes, like a hostile demon, steps disturbingly into
earlier dear, peaceful relations, and brings the human being, for his
beloved's sake, to forsake father and mother in a sense of the word in
which it becomes heavy and oppressive for the feeling heart! Against all
such sorrowful entanglements---into which an otherwise, in human life, so
glorious, significant and beneficent love can lead the human being---that
love is almost not to be called an unhappy one whose outer fate is so
unfortunate that the lovers, in the earthly respect, are torn from each
other, while its inner being, its inner purity, beauty, strength and
fullness, remains unassailed, and the spiritual property can therefore
remain unlost, pure and clear, however much the heart may suffer under the
painful earthly blow.

\begin{center}\rule{0.28\linewidth}{0.4pt}\end{center}

If any love is unhappy, then surely that one is, by which two
individuals---who either do not accord together at all, or else, over sheer
disorderly emotions, cannot come to a genuine peaceful harmony---are
brought to enter into a union in which they perpetually become occasion and
incitement for each other to a thousand feelings that weaken and consume
the spirit's power. It is difficult to say what in such relations most
gnaws on the soul and ensnares or burdens it: whether that a real love
continues to work with vehemence, or that the feelings that knit the
relation grow lukewarm and die away. The first becomes spirit-consuming by
the many penetrating and burning incitements, the many tensions of spirit
and painful anxieties, the constant unrest and the unceasing regards, among
which the soul, as on a raging sea, is cast waylessly to and fro under an
unquiet alternation of heat and cold. If, on the contrary, love is
vanished, what faintness, emptiness and loathing easily presents itself;
how spirit-dulling and deadening the shared life becomes, when that which
should set it in motion is gone, and other fresh life-movements are
constantly stopped, ensnared or pressed down by the never-subsiding feeling
of a chain that hangs fast and rests heavily on the soul. Doubly assailing
for spirit and heart the state becomes if the ensnared soul now comes
moreover to cognize that it is bound to an unworthy one---be it that the
allurements, notwithstanding (possible it is), still exercise their former
influence, or that it is over with the love. The continuing union becomes
without peace or without fullness, and even the liberation from the anxious
or oppressive outer relation will, if it be granted, still not free the
ensnared soul, which, with half and torn feelings, perhaps under inner
discord with itself, seeks to save itself by tearing itself loose.

Manifold are the different cases that here meet us when we ask how far such
bonds are to be kept in honour or torn asunder. Every peculiar case must be
decided for itself according to its peculiar character. There are relations
in which a purely and definitely feeling, upright temper firmly cognizes
that it cannot and dare not think of a breach; there are others in which
the definite feeling of separation's necessity---even if the heart has
perhaps not yet torn itself loose---is just as firm and reliable. Happy is
the one in whose inmost, under such unhappy relations, a pure and vigorous
feeling---not lured forth by foreign inclination and passion, but springing
from true esteem for the higher love and true fidelity toward the Eternal
One, and which in its calm and firmness bears witness to its
reliability---speaks with decisive voice, and bears and nourishes an
unshakable resolution, whether it now bid one remain or flee. But not
always is the mind, drawn to and fro among so many considerations, thus
soon and easily, by a higher verdict, fastened to a definite point, and the
disquieted heart raised up above the fearful doubts and the anxious
wavering. Then it is especially important not to hasten, but to tarry and
let the first movements subside, and not to hurry from feeling to will,
from resolve to action, so that slowly the nourished feeling may test
itself under life's advance and manifold collisions, whether it will hold
out under them. The imagination too must have room to paint out beforehand
the future toward which the human being's mind begins to be set, whether
this future will indeed continue to show itself in the shape it first
seemed to want to take on. Often people first learned, by being separated,
how much they were to each other: so let the mind then make a trial with
itself, by setting itself livingly into the representation that the
separation had already happened.

All the considerations which the question of the greater or lesser
permissibility of divorces calls forth, when they are regarded from
legislation's standpoint in a State---especially in a definite period of
time---we will here not enter into. Only that which, under such
entanglements, stirs in noble and pure human hearts; that which, for an
upright and vigorous mind, animated by a higher love, offers firm fixed
points for the perhaps wavering resolutions---only this will we make the
object of such general considerations as, in given cases, can now give the
decision to the one side, now to the other.

Provisionally we remark in general: although marriage in itself always
ought to be regarded as a union insoluble according to its inner
nature---because it must always be knit with the thought that it is to be
valid for the whole life---one may yet in truth not go so far as to oppose
every divorce which gross breaches of the marriage have not made necessary,
and may on the whole not restrict the permission of an amicably resolved
separation within the narrowest bounds. There is many a divorce over which
every sedate person who stands near enough to see what called it forth can
rejoice. Even if it is indeed not, in and for itself, gladdening that
divorces exist at all, and even if their frequency is a sorrowful sign of
the times, yet this sign is not at once to be regarded as the sorrowful
thing, but much rather that of which it is a sign. If the divorce is not
gladdening, the marriage was often still less so. It is herewith as with
legal disputes or lawsuits, whose increasing number is not a proof of a
State's increasing health, but which nonetheless had rather increase than
that the number of unappealed violations of right should rise, and
injustice go its free course. Directly, the holiness of conjugal unions is
surely not furthered by keeping in honour such marriages as are just as
little lived to the pleasure of God and men as they in any respect come to
the good either of the united persons or their children; even for the
children's sake the parents' separation can not seldom be the most
advisable. A calm, peaceful, secure existence is the basis for all genuine
growth in the good; under sheer tensions of spirit and unquiet emotions, or
in a stifling and oppressive surrounding, nothing can thrive in the human
being's spirit. If one lets that which the preservation of bodily health
requires be of so great weight, then let one consider also the spiritual.
There is in the human world not so great an abundance of spirituality that
one need not be watchful against all that by which the soul's health is
stormed in upon; and let one not regard it as a side-matter, in considering
a relation's character, whether an awakened and lively spirit is therein
consumed away and comes to ruin, or is nourished and thrives. If often
surprise and weakness, frivolity and a certain youthful arrogance, or
temperament's vehemence, bring people---before they have yet come to
clarity about themselves and about life---to knit unions that never ought
to have been entered upon, then, when such a thing is now, too late, become
aware, a loving solicitude---and such a one should after all take place
everywhere---should not let them suffer for it a whole human life through,
and suffer in a manner that can cow down their whole forward-striving to
higher goals. The human being must indeed have too much reverence for the
holy to let himself be easily and soon carried away to wish holy bonds
loosened because they seem to press, or because something more comely seems
to offer itself; and let no one regard himself as cowed or ensnared because
a relation conditions and limits, and brings duties with it that are not
always easy and pleasant. But in a life in which even the best must daily
pray: Let us not fall into temptation!---there are yet not so very seldom
cases in which people so daily tempt each other to countless hostile, evil,
bitter, dark thoughts and feelings that it is better for them that they
part. Paul's words in 1 Cor. 7:5---that Satan may not tempt you because of
your incontinence (διὰ τὴν ἀκρασίαν ὑμῶν)---have application in many a case
outside the one he speaks of.

Meanwhile it is thus, to be sure, merely for our sins' sake that something
must be permitted and called the best which therefore just cannot yet be
called good. And so it must at the same time be remembered that precisely
also for these our sins' sake many a bond must bind more strongly than,
with respect to the single case alone, would be needful, were the general
frivolity and unruliness not there. Many a disharmony that now leads to a
breach in which there is no blessing would perhaps not have come so far had
it not been held fast and furthered because the thought of separation's
ease made [people] less attentive and careful. When the marriage-bond first
shows itself as easily soluble, it also soon shows itself as less holy, and
gradually it then comes so far that already at the marriage's entering-upon
a greater and greater frivolity sets in. Therefore, in order to maintain
the principles and their effect on the whole---in order to maintain
marriages' holiness in general---that must sometimes be denied in the
single case which, with respect to the two individuals regarded by
themselves alone, would perhaps be the most advisable. Yet let one not
forget that to the whole the two individuals too belong! The necessity of
maintaining reverence for holy bonds can demand much indirectly; but let
one always consider that marriage, like the Sabbath, is there for the human
being's sake, not the human being for marriage's, and that individuality
has a great right. Better than prohibitions, which cut off all free
self-determination and one's own resolution, might well also such measures
work as, by setting a holy reverence's and religiousness's feelings in
motion, can bring the tempers with freedom to step back. It would surely in
this respect be of effect if marriages were dissolved just as well by an
ecclesiastical act as they are entered upon by it. As then, on the whole,
divorce is indeed an action of such a nature that it ought to happen with a
solemnity that can remind the human being that he ought in all things to
have God before his eyes, and do what he does with a devout disposition.
People who have once so solemnly been knit to each other, and have been so
much to each other, ought not, after all, to go from each other for good,
except with a mind attuned, by religious representations, to reflection and
earnestness and to peaceableness.

So much provisionally in general. For spouses in whom the thought of
divorce comes up, and for the one who stands near enough to be called to
work on their mind, there is still much to weigh. First: what it means,
from now on to have to think of oneself as a stranger to the one---indeed
perhaps to shun seeing the one, or at any rate, without former confiding
address, to have to pass by the one with whom one has been so inwardly
closely united, and toward whom perhaps in earlier times the heart many a
time has beaten with delight and longing. Especially must it not be easy
for a but tolerably deep and finely feeling womanly being, without coming
into inner unrest and confusion and feeling anxious at heart, to unite the
thought of the distance that should now set in with a living remembrance of
all that it once was to the man it now will tear itself loose from---how it
has been bound to him, how closely and how nearly. It can surely not
easily, when it paints out for itself the relation's whole character,
endure the thought that such a devotion---a life in which it has thus
belonged to another---should now be regarded as a terrible dream that it
must shun to think of or hear spoken of. And now, to the ruined or at any
rate embittered memories out of a perhaps very fair past, still the many
relations, together with the one relation, torn asunder or confused; above
all this: that, if children are there, these too now, right as other
things, are to be divided and torn away from the one or the other, and,
under an unnatural disturbance in their inmost, either become strangers to
this one or that, or else find themselves in an unclear and tense position
between the two. In truth, there must after all be much needed to be able
to will and resolve something that is so unnatural and offensive to
feeling.
And yet there are cases in which a human being, in order to preserve
fidelity toward himself and toward something Higher, must will this
unnatural thing. But then the human being wills it too only where a higher
fidelity demands it. Let the question here be not of a right, but of a
duty; not whether one may, but whether one ought to tear oneself loose.
Only when a holy feeling of duty moves the mind and drives to the
resolution, let the human being give this room and power with him. To the
actions morally arbitrary and merely permitted---by which, as soon as only
the outward occasion offers itself, delight and inclination dare rule over
the choice---a divorce can nonetheless never belong. And herein we have,
then, something that can orient and guide in many respects. Only when the
individual feels itself called and prompted by a feeling of separation's
necessity---springing from love for the good and from a faithful
holding-fast to this---does it grasp at resolving upon it. A feeling of a
holy obligation, proclaiming itself with definiteness and clarity, can here
alone offer a firm fixed point for the perhaps otherwise unsteady thoughts
and misleading wishes. Thereupon let the heart here test the arising
desire, whether it indeed has courage to call that which it desires a duty
for itself. But let it herein be upright toward itself, and guard itself
against deceiving itself. What is not genuine and good cannot lead to true
peace. The genuine keeping of the truly Good is the foundation for a
blessed existence, first in inner, then in outer respects. Holy laws, when
we transgress them, turn as hostile against ourselves; inwardly they let us
have no calm, and outwardly we dare no longer henceforth in life hope for
protection from them.

But when there is now only talk of an ought, and of that which a holy
necessity demands, then let one consider, for the first, that the fidelity
which belongs to love's essence is itself something in a high degree holy.

Although that which from the first knits two human beings to each other is
something higher than both, and they love each other because they love the
good and the beautiful, yet love attaches itself wholly to the person, as
person; and because this person was to us and is to us what it was and is
to us, we feel ourselves bound to it. Therefore let feeling not at once
believe itself entitled to step back, and perhaps to turn to other
quarters, because the beloved begins to seem less captivating or less
lovable! Even if the personal attraction and the personal pleasure decrease
in inwardness and sweetness, yet the faithfulness and the proper love and
solicitude by no means decrease! Precisely then let the heart feel, and
nourish the feeling, that it was not, after all, for nothing that it gave
itself to this definite individual; that it would not merely receive and be
gladdened, but would also communicate and show itself active in solicitude,
and in all things would share fate with the beloved. Least of all dare
foreign inclination---delight in this or that which seems more captivating,
or an arising new passion---get room and influence. To let such things get
power over oneself is to cast oneself into a whirlpool that can at last cast
us against the most terrible cliffs, or cast us out onto the most desolate
shores; it is to call all that is holy up against oneself. Even that which
most of all might be able to seem to be inviting---the conviction of the
co-bound one's moral unreliability or moral unworth---can still not at once
be enough for a vigorous and soundly feeling heart to want to be parted
from the one to whom, in so firm and significant a manner as by a marriage,
it is knit. As long as it still has that power over the other heart, that it
can work fortunately upon it and can keep itself pure, it will precisely
feel itself invited to remain---indeed would seem to itself as one who would
flee a holy calling, if it thought of tearing itself loose. Here let Paul's
words in 1 Cor. 7:16 hold: ``how knowest thou, wife, whether thou canst not
save the husband; or how knowest thou, husband, whether thou canst not save
the wife?'' Only when the human being feels itself unable to work, and now
at the same time perhaps in danger of coming into tempting and heavy
struggles out of which it perhaps cannot save itself with a pure and whole
mind; or when it sees itself degraded or maltreated, or, without benefit to
anyone, bound and limited in its striving toward the better---can the case
set in in which it feels itself called to demand release. Holy laws also
demand their holiness maintained thereby, that the human being shall not,
while it itself tolerates the unworthy, tolerate the holiest to be drawn
down and desecrated. To such desecrations can then also sometimes the
vehement jealousy be reckoned, when it in its wildness destroys or spoils
the best in love. There can be cases in which it---when otherwise nothing
will help against the consuming sickness---can become a holy duty, both for
the one who suffers under his own jealousy's ungovernable vehemence, and
for the one who suffers under the beloved's, and that for each of them a
duty both toward himself and toward the other: to grant himself and each
other rest, by dissolving cleanly the relation that causes them so much
pain and plague, and is to them so great an occasion or temptation that
very much is stirred up in their inmost of which it is never good for the
human being to be filled. As regards the solicitude for the common
children, if there are such, this can in many cases just as definitely
demand the marriage's dissolution as it can, in other cases, give the one
or the other of the parents strength and firm will to accommodate
themselves to everything in order to maintain it.

An amicably completed divorce is just not always---even if both parties'
outward consent is had---done with both parties' true inner accord,
especially not always from the first. Where, on the contrary, both right
from the heart wish and desire the divorce, there seems, with respect to
the two individuals in and for themselves, to be no scruple about it, and
there seems to be no question of a duty of fidelity, where both with equal
delight release each other from it. Meanwhile: fidelity here has in many
respects its ground in a higher one, and hangs in so many ways together
with this higher one, that one must not believe everything to be clear
because both with equal delight and willingness come to meet each other in
their wish, and because there is no third either who could thereby suffer.
Let the human being be also faithful to himself and his resolutions, and
hold faithfully fast to his youth's best memories and feelings, as to a
faithful friend. What in its life's best moments filled it, let it have
respect for that in the many faint and darkened ones. A breach with the one
who was once dear to the heart---even the gentlest and mildest breach---makes
a disturbing rent into the whole life; and it does not, in many respects,
come to the human being's good to break off its life's sequence, or to cast
a whole elapsed life from itself. Such a change is always precarious and
dangerous, and can easily, for the future, make the life---which is now come
out of its former direction---uncertain, unsteady and wavering, unless firm
and vigorous feelings of a higher kind determined the resolution, and
therefore step in in place of the lost ones and give the life new
continuity and steadiness. Of mere delight and fancy let one also here not
make the holy dependent; here too only a certain necessity, a prompting of
moral nature, can bring about the authorization. It is surer to govern a
mind that unruly will grasp after this or that than to give it the rein. If
we are all here on earth as on a way with each other to a better world,
then we must hold out with each other, when we have first resolved
faithfully to go along together, and not at once go from each other to the
right or the left, merely because it seems more convenient.

For the rest, it is not indifferent of what character the relation is and
was that is broken, neither what kind of union it is, nor how it was knit
and began. The merely provisional union, the betrothal, binds naturally not
so strongly as the marriage; and a relation that avowedly was entered upon
without proper love may sooner be loosed than one which real
falling-in-love knit.

When it has not yet come to marriage, there naturally falls away a part of
the considerations that make the marriage's dissolution so precarious. Not
yet have both given themselves wholly over to each other, and forsaken
life's former course and manner in the degree that the final union brings
with it; not yet is the last and decisive holy vow given. The greater ease
of the dissolution is already a direct consequence of this: that the
provisional union can be regarded as a kind of trial, even if the united
ones did not regard it so, or would have regarded it so by others. So holy
a relation as a marriage should, after all, not be entered upon with half a
heart either, and it would be so if the desire to become free of it had
already arisen. The one who feels this desire cannot, after all, so long as
it is not struck down again, give the holy vow---as it is to be given---from
a full heart without inner reluctance; so that faithfulness and uprightness
themselves seem to demand that he not conceal his heart's true feelings and
thereby get the appearance of promising what he is not able to yield. But
with all that, an orderly and vigorously feeling and willing temper---which
does not let every arising emotion or mood at once get room to rule, or
even only to speak---will be slow to step back, or to utter the inclination
to it, where it has bound itself by a vow which, even were it only a
provisional one, was yet a holy one, and in the heart already a decisive
one. It must here too feel that it must remain faithful to the once-chosen
being, and remain faithful to its life's good moments, so long as no
feeling of a higher nature demands the dissolution. The consciousness of
having grasped disturbingly into another human being's inner soul-state and
outer condition is too oppressive for life to be lived freely and lightly
while one bears it with oneself. Meanwhile, under the conflict of motives,
those who demand the relation's preservation cannot here, to be sure, work
so strongly as at the marriage. And once it is clear that the two do not
harmonize and cannot well hope that they will come to live a blessed life
with each other, then it is best that both unite about separating, so that
they may not seem to want to tempt the Lord, their God, by yet casting
themselves into the relation and giving their life's peace up to be
sacrificed. Also this difference always remains between the marriage and
the betrothal: that, so long as the soul is filled with doubt, or at any
rate no dependable ground of determination has come to mature strength in
the soul, so long there ought not to be acted; so long, therefore, ought a
marriage neither to be dissolved, when it already is there, nor to be knit,
if it is not entered upon; the betrothal ought, therefore, so long not to
go over into a conjugal union. If, however, the one resolves, for the
other's sake---notwithstanding the inner reluctance---to hold fast, so as
not to burden itself with the feeling of having disturbed the other's peace
and weal, then let this resolution first come to maturity and strength, and
let what is resolved happen now with devotion to God and with a whole
heart.

A union which was never a real inner soul-union---in which never a heart
wholly, from the heart, gave itself over to another---may naturally soonest
of all unions be dissolved, as soon as it was at the same time avowed to
both that mutual feelings of love did not draw them to each other. Since
here only outwardly and in outward respects fidelity was promised, so here
only such an outer fidelity is broken by the dissolution: it is no really
genuine one. When a being, by surprise or a kind of overwhelming in early
youthful age, became a prey for another individual---who well knew that it
was taking possession of that other, either without returned love, or else,
though a certain returned love was stirred forth, took possession of it at
an age in which this could not at all yet have come to that point of
maturity which it is a holy duty to await before a holy vow is
desired---when a being has thus become another's prey, it is natural that
it, already merely for that reason alone, later on, when it comes to clarity
about itself and life, can demand the relation's dissolution, prompted to
this demand by a just indignation, so that the cunning, or violent, or at
any rate high-handed and always unlawful and presumptuous
taking-possession may not continue to bind the unworthily treated one, and
count as lawful and good because it succeeded. A proceeding that in itself
was unlawful and presumptuous ought indeed to bind the one who used it, but
yet not to bind the one who was the object of the unworthy treatment.
Meanwhile a nobly feeling heart will yet, even in such a case, be slow to
resolution, and cautious. For however much that which is to happen is
freedom's and deliberation's affair, yet that which has happened, so far as
it has happened, gets, for the one who suffers under it, the appearance of
a higher dispensation. And the heart that sees itself fettered cannot think
of itself as one that has, without a governance it must respect, come into
its present condition. Even its earlier blindness or lack of strength it
can now reckon hither; and therefore it must here too guard itself against
resolving something whereby it shakes the ground on which its existence
rests, casts itself into an unsteady and unfit%[FLAG: rendering of "uskikker", printed p.149 — perhaps "uskikkelig/uskikket"; see transcription note]
course of life, and sets itself in discord with itself or its nature. Let
it therefore await the point of time when a pure and holy resolution comes
to maturity and firmness, and when it feels itself sure of itself and of
its resolution's purity and goodness.

On the whole, under whatsoever of the many unhappy entanglements a temper
may find itself placed, let no resolution be taken in vehemence, or while
the mind still unquietly works with itself! Where important resolutions are
to be taken in matters that turn about one's own weal---especially where it
concerns the maintenance of one's own I in opposition to other people---the
selfish or egoistic emotions easily come up, under all kinds of shapes,
even if it is clear what ought to happen, and make the mind unclear, and
defile or pollute the motives that alone should animate the human being. An
upright temper can very well perceive whether something that is not good
and pure plays in among the feelings; and let everyone be ready to confess
to himself that he is confused and assailed by that which gives us all
enough to do with in our inmost, and which no one can call himself secure
from. It is not for nothing that we should daily pray: Deliver us from
evil. If we were exalted above all its assaults, why should we pray thus?
So let the mind then first bring itself into the clear, and come to inner
clarity and calm, before it begins outwardly to work and act, in order to
get accomplished what it feels to be right and necessary! It is not
indifferent with what mind---and moreover not either at what point of
time---that happens which perhaps in any case is to happen.
\begin{center}\rule{0.28\linewidth}{0.4pt}\end{center}

Leaving an object at whose contemplation one can dwell with some delight
only so far as one can hope, by throwing light on it, to prevent or resolve
confusions, we will at once go to another object of which the same holds:
the unhappy relation, or collision of relations, under which the heart
cannot follow its pure, free inclination and, after its delight, dispose of
its future with respect to an individual who desires a union with it,
without offending against the feelings it long calmly nourished toward
other persons---especially toward parents---and making a disturbing breach
in the relation to these. Also under such collisions in life it often
happens that the mind is cast to and fro between opposite representations,
and sees itself, on neither the one side nor the other, in a position at
once to come to a pure, clear, calm- and peace-bringing condition.

Two different states and relations it is here of importance to distinguish.
When the heart, moved by love, vehemently desires to knit a union contrary
to the wish and will of those by whom, from youth, it has been accustomed
with holy reverence to let itself be led, the relation is quite another
than it is when there is talk---because these wish it and spur to it---of
having to enter into a union without loving. In the latter case the
resistance is far surer than in the former.

Already one's own personality must be too holy to the human being for it to
let its spiritual welfare---perhaps in the highest and most important
respects---be set at hazard for the sake of lower calculations and
considerations. But still more must a vigorous feeling for the conjugal
union's holiness give strength to resistance, where it concerns whether we
are to tolerate that mankind's fairest possession shall be trodden down or
perverted in our person. Every human being can and must in such cases feel
himself as mankind's representative.

To enter, with suppression of the heart's best feelings, into a union which
according to its nature ought to be knit from a free and total inclination
of the heart, is in itself always a profanation of something holy, a
degradation of that which everywhere ought to be held in honour and
respect. The conjugal union must be esteemed too highly for it to become
the object of any kind of compulsion, or of such foreign influence as binds
or ensnares. To this comes that, where the persons for whose sake a human
heart should sacrifice itself are led merely by such considerations as are
of self-interested or vain origin, or are led by all too great anxieties
for the future---which yet always stands in the Lord's hand, and with
respect to which no one can, after all, be his own providence---it is even
for these persons' own sake precarious to give way. To comply with others
in order to support them in low earthly cravings and plans---indeed in that
which may perhaps one day fall heavily and oppressively on their
conscience---is in truth, after all, not from the heart's ground to wish
them well. What is not of the good is also not to the true benefit of
anyone. This consideration can also find a place there, where it is just
not low considerations or petty scruples, but more natural and human
feelings, that are in motion. Unhappy children are, after all---as we have
already once had occasion to draw attention to---so great a misfortune for
loving, or merely humanly feeling, parents, that precisely the children
themselves can, from this consideration, draw strong motives for firm
resistance, in order in time to secure beloved parents against that
heart-rending misfortune. Meanwhile: in a life in which the soul that
seemed most called to stir freely in a higher sphere, and most longs for
it, is not seldom pressed by spiritually narrow circumstances; in a life in
which the human heart's fairest and best desire and striving so often does
not find room and freedom for a fair and harmonious unfolding; indeed in
which it can sometimes even be anxious-making, beside the need in which
others find themselves, to live in beauty, luxuriance and fullness---in such
a life there may well be peculiar cases in which the heart dare not even
hold fast to that which it is otherwise in general a holy duty to maintain
and hold in honour. Such as human beings once must live with human beings,
such as we all once must suffer pressure and afflictions together with each
other, and not regard any human need as foreign to ourselves, there may
well set in relations in which beloved beings' need can step so near a
heart that it can feel itself called and prompted by the best feelings to
sacrifice its inner spiritual well-being in order to save those whose weal
must be to it in a high degree of concern. In the midst of the full
cognition and feeling of what it is that one determines to bring as a
sacrifice, this sacrifice can yet in such cases be that which most answers
to a devout love's deep feelings. But always such particular cases belong
to the exceptions that can seldom set in; and let it never be weakness, but
the moral strength, that brings one to give oneself over thus. What ought
to be holy for the whole of mankind, let it be so also for each single one
in his individual position. Let the human being not too easily and soon
give up his spiritual welfare to the troubles and unhappy confusions a
marriage can have as a consequence which is entered upon against the
heart's inclination. Let every human being regard himself as a plant in the
Lord's garden, and keep guard over himself as such a one, that it may grow
and blossom to the pleasure and glad contemplation of God and men. There
are sacrifices which, even if they seem to disturb the growth and the
fullness, yet precisely---when it happens rightly---further life's
luxuriance and beauty; indeed the most painful are sometimes only
transplantings to a better soil, under a fairer sky. But there can also be
sacrifices which, even if they happen with good will, yet destroy or cow,
because they attack existence's inmost root.

Of quite another kind than the cases now spoken of is the collision of
conflicting motives and promptings that sets in when there is talk of
knitting the union decisive for the whole of life contrary to beloved
parents' definite wish and will. What in the former case can in general
summon to firm resistance will here, on the contrary, sooner lead to
letting go the heart's desire---at least make it far more precarious to
follow the inclination: precisely the feeling for the conjugal union's
holiness and high significance. Just because it is so fair a union, it is
natural, and a consequence of inward feelings' penetrating strength, rather
to forsake than to knit it, so long as all is not in the clear, but
something disquieting is left behind that will disturb its pure peace and
beauty. It is too holy for it to be entered upon while precisely thereby
other dear, holy and significant relations are confused, and the soul
cannot undividedly and wholly feel itself well and give itself over with
full delight. There are indeed positions under which the ensuing breach in
the former relations does not signify so much; for not always have, in the
single cases, the relations in reality that significance and worth that
they, according to their original nature, should have. But where the
desired union would become really heart-rending for esteemed and estimable
parents---who perhaps, if their children thus forsook them, would regard
them as lost---one ought yet not merely to heed their aversion or warning,
but also to esteem even the desired blessedness too much to buy it thus, buy
it in so imperfect a shape that it would yet only half become that which
the heart desired, and which alone can satisfy it. Here too it holds that
something holy is profaned when the conjugal union is entered upon with
suppression of the heart's best feelings. By forsaking, there is not made
either, precisely at once, a decisive determination for the whole of life;
if the love is genuine and dependable from both sides, then the inner
soul-union can still remain, and with this---which is the most
essential---the hope too can still be held fast. Whereas the step that is
taken, when the union, notwithstanding all that stands against it, is yet
entered upon, binds for good, without being accompanied by the blessing
that is so beneficent for mind and heart. The human being will so gladly
dispose of his fate alone; but at the most important steps let him not be
too hasty and too self-willed, but consider that the good must be appointed
us by a higher governance if it is to become for us a good, and that the
human being must take important and significant hindrances as a hint to
thoughtfulness. If, on the whole, every marriage that is to become blessed
on earth is to be resolved in heaven, as we have before a couple of times
reminded, then here especially ought the inner higher accord so completely
to set in that it lets us with complete peace and security do that which we
do.
\begin{center}\rule{0.28\linewidth}{0.4pt}\end{center}

We come now to the kind of unhappy love which, by this expression, one is
most apt to think of: the hopeless one, which does not find returned
love---indeed often does not even dare to demand or desire it, but yet
continues, and now causes unrest and pain. For even then, when it would be
wrong to desire real returned love, and this wrong is also really cognized
and felt---indeed when it is even love's own truth and depth that brings
with it that the one in love cannot wish it---even in such cases the
inclination is yet still easily nourished, and the sweet delight and joy at
the thought of the adored one. Under all the pain and sorrow under which
the mind, enclosing itself within itself, suffers, it cannot yet tear
itself loose from the source of all this suffering, but holds fast to its
inclination, and keeps it as an unspeakably sweet possession. The deep joy
in a fair being, the infinite devotion and participation, opens and widens
mind and heart; and, right before the entrance to a heaven-kingdom, the eye
cannot leave off perpetually beholding into it, even if the entering-in is
not granted. The fruitful germ of a whole new, great and glorious life
will---even if it cannot come to the unfolding at which it aims, indeed
precisely then even the more---captivate mind and contemplation and fill the
soul with expectation and longing. It also worked---however precarious and
in general little desirable it often is---sometimes with unspeakably
beneficent power in a human soul; and where it is a human being's
dispensation that its spiritual life is to be nourished and warmed by
suffering and pain, surely especially the one whose suffering springs
especially from a hopeless love can come to perceive ``these sorrows'
secretly forming power, deep in the heart,'' as Goethe calls them in one of
his little poems. Yet let the one who bears so fruitful a sorrow and
pain---fertilized with so much heaven---in his inmost, hold himself fast, and
flee the distraction of the mind or a wandering grasping after this and
that, and concentrate with power the fullness of his feelings. Let him
learn: not to forsake his love's inner content, its true inner being, even
if he must forsake the inclination's earthly striving! That which in all
love is higher than the desire that draws to itself---the joy in the
beautiful in another being---a truly philosophical, that is, a truly
Idea-loving and for-the-Idea-living temper will hold fast in the midst of
the loss. But in order that thus the higher love may yield its heaven, the
desiring must fall wholly silent, and the appetite to draw to oneself and
possess must not perpetually incite. What perpetually incites and
irritates, and, clinging fast to all representations in their circulation
in the mind, like a fixed idea perpetually returns again, burdens thereby
and presses the mind, and leaves not even room and freedom to enjoy the
highly adored object when it offers itself to contemplation for full
spiritual enjoyment. The Psyche thus ensnared does not then even come, in
the midst of its right element, when the mildest air fans and lures, to
unfold its pressed-together wings. Perpetually stirred to grasp, without
being able to, the mind does not open itself for the refreshing
contemplation and a pure communication's pure joy. If the heart is not to
perish pitiably in the glorious one's glad nearness, in the midst of its
in-itself heavenly influence, but openly to receive the full, pure
impression's refreshing power, then the goad of appetite must not
perpetually sting and burn. With a vigorous self-determination and
self-mastery it must be suppressed or torn out of the soul. For which, to
be sure, many an outer means too is necessary, or at any rate helping and
furthering. Sometimes there already helped a free and fresh turning of the
mind toward life's other manifold glorious and fair possessions; for life
is rich and can in many ways set [one] in new movements, when one rightly
grasps into it, and it then offers material enough for the mind, filled
therewith, to be able, even in the loved being's nearness, to move itself
in freer, more glorious, enlivening regions---whereby the enjoyment of the
loved being's beauty, spirit and heart can precisely be heightened and
become purer and fairer, while the spirit now loosens itself from the more
petty thoughts that formerly perpetually anew came up and ensnared the
mind. Of itself this does not, to be sure, happen, but only on the strength
of the acceding brave and vigorous resolution that indeed dares to let the
fair Heavenly get its full beneficent influence, without the delight bound
up with it getting free play or playing the master. It is something so
unspeakably refreshing, enlivening---something that with so beneficent a
warmth penetrates the soul---to rejoice in a loved being's beauty and grace,
heartiness and liveliness and spiritual captivation, that one should just
not at once forsake that joy, as [one would not forsake] bathing the soul
in so mild and enlivening a sunlight, even if the appetite will show itself
unruly; but that one should, on the contrary, so compel this that it dare
not confuse or pollute that soul-joy. Only that it not steal itself too much
forward. It has always a treacherous tendency. Like a greedy worm it will
undermine and hollow out the fair life and the sweet feelings, when it is
not cowed. It is a fire that consumes and lays waste when it gets power to
blaze up. It will cast [one] into a life full of unrest and unpeace;
indeed, when the way to satisfaction for the appetite goes only through
wrong, it will cast [one] into a life full of dissembling, false and
mendacious conduct, and yet not come to its goal or attain fulfilment; for
the enjoyment, which cannot become an enjoyment from the heart's ground,
will---disturbed, half, and without totality---as it were vanish under the
hands and keep on burning and consuming where it should seem stilled. Such
satisfaction---which even in a criminal manner goes out to ruin that which,
for feeling, should precisely show itself as something inviolable in the
beloved, without which love itself could not will to subsist---such a
satisfaction is yet, in and for itself, already a wretched satisfaction.
And therefore let one then take strength to flee altogether, if it must so
be; for this means too is often necessary and in itself surer than that
other way out, which demands a more philosophical power of spirit. But the
adored image let one yet always take with [one] as a refreshing, sweetly
enrapturing possession, which even by the removal from the object can gain
in captivation and beauty, because its deep fundamental features can step
forth purer for the more calmly beholding and surveying gaze. It is, to be
sure, not easy to withstand strong allurements' inciting power. The surest
is always to flee from them. Separation does not, to be sure, easily loose
a bond between hearts, nor will it so easily---even where no mutual
inclination knit a common bond---abolish the power in that allurement
whereby someone is drawn toward another by means of the cognition of the
heart's true, deep worth in this other. That, then, separation shall not do
either, but only abolish the disquieting, anxiety-making and painful in
that allurement's power, and make [one] able more to have good of that
glorious cognition. From an unsound inclination, on the contrary---one that
was not even called forth and nourished by such a cognition---separation may
well free [one]. Where it was only the merely outer allurements, be it the
bodily or the spiritual, that captivated and ensnared the soul, the removal
works sometimes more and more quickly than one should have expected. For if
the loved figure has no genuine content in it, then the image of it, when
the momentary refreshings of the appetite no longer perpetually return,
will not be able to hold itself in contemplation, because it is not able to
fill this. Yet even after a long lapse of time an again-ensuing approach is
always precarious, and the well-known [saying] that old love does not rust
can only all too easily confirm itself. Those allurements worked so
strongly because they went into league with something in the human being's
inmost that is always strong and ravishing: the idea of beauty and the
feeling for the beautiful; for this was awakened by the contemplation of
those alluring merits, and these can now, by recollections and by calling
this idea of beauty forth anew, easily again exercise their former power.
Especially can the newly ensuing approach be dangerous when the loved image
in itself was fair and glorious and worth being loved so highly. Yet this
image itself let one not be robbed of---especially not let it be perverted
or spoiled by any self-numbing or bitterness.

Even bitterness and resentment are easily stirred up by the unhappy love,
not merely toward a more fortunate third, but also toward the loved being
itself. Embitterment and wrath, and a certain desire---even sometimes rising
to cruel thoughts---to assail in a hostile way, indeed as it were to tear
asunder, the being that cannot be moved by so much adoration, easily
reaches about alongside the love; because, when the strong craving and
yearning comes up, then the egoistic is in motion, but this is by its
nature just as near to uttering itself in a hostile and repelling way as to
showing inclination. The vehement incitement stirs and stimulates, but
therefore also irritates and embitters, when it, because the satisfaction
does not come, becomes painful. Thus there then easily steps forth, in a
new shape, jealousy's heat, no less pernicious than in its previously
considered shapes. This is the worst to which it can come; and if any such
thing stirs up in the mind, then it is surely high time to consider that it
is better to flee than to suffer such ardour. Love shows here again its
double nature: how it is not merely of divine, but also of an altogether
earthly descent. The appetite for exclusive appropriation---to possess and
call one's own what is the object of pleasure---spoils here again the fair
possession which, when there was loving with more philosophical power and
depth, could be enjoyed and preserved. Therefore let the soul captive in
the hopeless love think already early of preserving for itself the fair,
full image it has taken up into itself. To bear a pure and glorious image
of a loved individuality's beauty and fullness and whole idiosyncrasy
within oneself, and precisely to bear it within oneself with a pure and
total devotion, free of desire, is something unspeakably refreshing,
enlivening and inwardly nourishing to the soul. Precisely of such a nature
shall also that rest be to which the happy lover is to come by a happy,
all-desire-stilling union; and for the hopeless lover there shall thus, out
of love's own fullness, the higher cognition and joy in the beautiful and
divine go forth, which can lead the soul---for a time ensnared and captive
in a narrow circle---back again to a freer existence, and become the
fruitful germ of a new life moving itself with new power in higher regions.
So that it must again be fulfilled, that to the one who knows how to let the
Divine come to dominion in his inmost, everything must come to good. For
which, moreover, a certain definite activity, penetrating right deep into
life---whether it be as study and development and contemplation of ideas, or
in a more practical manner reaching more immediately into the outer human
life---is a necessary condition. Especially able to free the mind and
further the new life's power is such an activity as stands in near
connection with the love for the Eternal, and offers manifold occasion for
high ideas and feelings for the great and beautiful to be set in motion in
a manner that tears the individual out of its individual life's circle of
thought, right into the great Nature and the universal, everywhere rich and
significant life.
\begin{center}\rule{0.28\linewidth}{0.4pt}\end{center}

But when we thus, by the hopeless love, point to love's own content and
fullness---independent of personal acquisition and possession---as to the
well from which a new, vigorous life is able to spring forth (surely not
without pain and resignation, but yet without stifling and losing the
working of the Divine, the θεῖον, that came near the moved mind and touched
it): how much more must we then point to the same, when we ask what the one
has left who, after having lived through love's fair spring in hearty joy
and devotion and sweet communication and confidingness, is now, by an
outward dispensation, altogether bereft---indeed not bereft of the returned
love, and therefore not either bereft of the beloved, but yet of this one's
presence and immediate possession; whereby we naturally especially think of
the dissolution of lovers' shared life that happens by death, though yet
still other nearly as decisive separations are conceivable.

Just as natural as it is vehement is the pain, when the loving heart is torn
out of the secure, calm possession, and now, feeling its whole existence
torn asunder, stares out into the dark, gloomy future. What gave life so
much significance and allurement, what constituted the centre about which
life in so many respects gathered itself, and whereby so many thoughts and
feelings, so much delight and activity, were set in motion, is gone, and
life seems---bereft of its best content---empty and desolate. Precisely in
that part and sphere of life which penetrated the mind with feelings of
eternity, the dispensation grasped annihilatingly in; and when now the
strained heart will, in thought, move itself over to the beloved, and asks:
where? whither?---and therefore also easily comes to ask: to what end, after
all, and why, it had to suffer the painful loss, and now seems to itself to
have no firm point to which it can fasten the thought---how near are then
despair's feelings, and easily ready to come up and reach about? But that
which is most easily able to loosen the temper---in the beginning often as
it were petrified---so that the pain makes itself air in tears and is
thereby softened, is: the living remembrance of the called-away beloved's
goodness and love, whereby its own deep love comes to life and movement
again, and thereby also all the fair emotions that lie in love or follow
with it---this is also that whereby a new centre for life can form itself
again, and existence can obtain new content again. It is not what one calls
distraction that is that which consoles or helps; at most it can yield an
aid toward procuring rest, when the pain's outbursts are all too violent
and assailing. It is not forgetting, or the mind's turning-away from the
beloved's image, but precisely an energetic holding-fast of this image---a
genuine and deep remembrance and recollection---that is able to uphold and
bring new life-power and life-delight back. The former also finds only
little entrance in the sorrowful heart, which rather thrusts it from itself.
Every true and good word over the snatched-away beloved, every living
painting-out of that one's lovableness, is on the contrary consoling and
comforting. What an inconsolable consolation is, after all, that too: that
with time the strong and inward feelings will surely lose their power, and a
thousand new impressions will surely crowd out the now so vehement one, in
that they will press the thought of the beloved more and more back in the
loving one's inmost. What an inconsolable consolation, that the pain will
surely lose its power, when to the loss we have suffered still a new loss is
to come: the loss of that which we still have left---the glorious, wholly
soul-filling and captivating memory, at which we, however much it brings the
pain to well forth anew, yet so gladly dwell, because by it, in the midst of
the pain, we nourish our love and move it in our inmost. How natural is it,
on the contrary, precisely when the beloved is snatched from us, to hold
love and its being fast, and with it at the same time to preserve the
infinite joy in the beloved, and as it were wholly to fill the soul with
this one's glorious image. When the soul with elastic power crowds out what
will overwhelm and darken it, then love's feelings themselves will go
strengthened forth out of the inner pain, as out of a cleansing fire, and
penetrate the soul again with new inwardness, and fill it again with pure,
fair feelings of eternity. The more essential the loss was, the more
beneficent the beloved's nearness, the more full the enjoyment of love's and
the shared life's whole beauty, the more will the lover keep back, and the
richer will, in the soul's inmost, the communion with the parted one still
be able to become. To have owned is already a glorious owning; and it is the
more glorious, the more we owned. As the beloved's image is kept and
preserved undarkened in a faithful imagination, fresh and blossoming, so
love too can keep itself fresh and blossoming in the faithful heart. Indeed,
it is now precisely nearest to being able to transfigure itself into the
highest purity, in that all the petty [things] that otherwise many a time
confuse fall away, and the desire that falls silent lets a pure
contemplation get room altogether. Love now melts most easily together with
devoutness's and God-devotion's pure feelings, with the thought of eternity
and with a perception of the Lord's Spirit, as the one in whom we all live
and are moved, and who therefore as it were constitutes the medium in which
all those who are loved are in and with each other. Separation shows here a
transfiguring and sanctifying power. It lifts the beloved figure away, out
of the moment's pressures and assaults, over into a region where everything
stands with pure, great features, and preserves itself in youthful
freshness. What is taken up into memory's sanctuary, that we keep
everlastingly undarkened. And now it is thereby brought---the purer and
clearer it stands for us, the more---back to its eternal source, so that this
its eternal source may rightly step forth therein. When the God-filled one
otherwise everywhere, in that which fills him with deep love, loves his
God's infinite nearness, then he will now with double power feel this in the
midst of his longings for the snatched-away one, and in his still-continuing
spiritual attachment and devotion to this one.

And when now the heart, still loving with the former inwardness and depth,
rightly resolves to hold its love fast and to keep on living for the
beloved, then let it live so completely for and with this one that it lets
new liveliness get room in it again, and opens itself anew for the joy in
life! Precisely in order completely to live to the beloved's pleasure---if
this one, though snatched away by death, can yet still (and this the loving
being so gladly imagines to itself) look on to it---let it let that
life-delight work again which in and for itself is so enlivening and glad
for the beholder, and without which nothing fair and good can live and grow
and blossom in the mind. Or shall, under the perhaps long waiting for
reunion, the faithful loving being be cowed and come to ruin precisely by
means of its faithful love, and just for its sake as it were not be able to
come to meet the loving [returning] one in complete fullness of life? And
if the hoping, God-devoted heart cannot do otherwise than think of the
state after death as that in which the good one beholds and lives in a
clarity, spirituality and blessedness to which our highest beholding and our
highest soul-joy relates only as a gleam to the full light, then let it also
give place to the glad in the thought that the beloved is gone hence to that
realm of spirituality and glory! If it is love's nature that it brings an
individual thus to give itself over to another and forget itself in and for
the other---that the beloved's weal lies in the highest degree on the loving
one's heart, and that the thought of the beloved's blessedness is highly
blessedness-making---then let love show here the same power, and, with the
continuing and now even more gathering and strengthening love, let the joy
too---that the departed one now fares well---win room and power in the loving
temper. Even if it is death that has parted the lovers from each other, the
words can yet gradually begin to find application which Goethe, in his
Tasso, lets the noble Princess Eleonore say: that with a distant friend,
when we know him happy, we still live together in a right particularly
kindly manner. It knits, on the whole, the feeling and hoping human being in
a particularly inward manner to the Yonder toward which we all go, to know
beloved beings there, in secure keeping, unassailed by earthly life's
upheavals and confusions. When the lovable figure really contained that
which was worthy to fill out a deeply feeling being's love and its longing,
and it now lives on thus in memory's heaven, then let the joy that such a
being has lived step to the [joy] that it is now in secure keeping and in
its home.

And thus, then, can also the pain over the beloved's loss show a sanctifying
power, and through it a fair consecration to a pure and God-devoted life be
prepared for the one who takes the dispensation, as every dispensation shall
and can be taken, in such a way that it can and must come to his good.

But in order to take it so, one must not see in it merely a fate, but also a
dispensation---namely a dispensation from the God whose ways are inscrutable,
but with whom there is prepared for the God-fearing a place where a glory
shall be revealed upon him against which this life's sufferings are nothing.
Regarded as mere and bare fate, suffering is cold, hard and icy; regarded as
a providence's dispensation, it draws the heart moved by sorrow over to an
infinitely rich heart, with whom all creatures---under the outer natural
necessity's iron-hard, unbending course---have a refuge, and with whom all
is in secure hands, but from whom also every fate is to be taken as a voice
and an admonition. Taken thus, it sets devotion's and humility's deep
feelings in motion; and when only the human being will not, in his despair,
go to law with the Lord---as if he had suffered wrong because he did not
enjoy the gift, which was always a pure grace, long enough---then even so
bitter a suffering, touching so strongly the most spiritual in life, can
become what suffering so often was in a life in which suffering still
belongs to the best we have: a means of awakening and a source for a new
life, which is lived, still with deeper feeling and more inward power than
before, as a struggling and self-cleansing life, as a life in God and a life
for eternity. To die---so taught already antiquity, and Socrates, in Plato,
in his last hours calls this thought forth for his friends' contemplation,
and again Christianity in many ways points to the same---to die, that is
something the wise man his whole life through strives after and exercises
himself in. Namely, to be rightly parted from the earthly and sensuous and
cleansed thereof, and to bring it so far that the spirit can rightly live
purely spiritually for itself and in itself; that is, to die away from the
earthly and from all desire. Out of such a death new life wells forth---indeed
in the sound human being new fresh joy in life and delight in grasping into
and taking part in its outer manifoldness, and in enjoying also the moment
and what things' perpetual alternation brings with it. This can precisely
the one enjoy purest and freest who, free of all passion and all burning
desire, has not fastened the heart to the alternating things, or made them
the ground and seat, or the centre, for his life's joy, but under all life's
alternation strives to live and be moved in the unchangeable. But to come so
far that all particular desire and yearning, with its heat and unrest, goes
up into an undivided joy in the Eternal and in a pure onward-streaming
activity, the human being is brought on many ways. Whither suffering and
pain are to lead the one [group], to the same shall precisely the full
satisfaction and fulfilment lead those to whom it was granted in its full
extent: for with them too, thereby, life's incitements are to be stilled,
and a calm, secure, passion-free life set in. And it can often be uncertain,
when one looks upon life as a whole, whose lot was most to be prized: that
of the one who enjoyed the shared life with the beloved unbrokenly his whole
life through, or that of the one who, because it was broken off early,
now---in a clear and pure memory---took it, in its full blossoming and its
purest beauty, without its being given up to be sacrificed to earthly life's
assaults, with him into his remaining life, and now therein owned something
out of which an unresting striving to contemplate the Eternal, and to dwell
in its contemplation, and to work toward its spreading, went forth in his
spirit. For it is not the mere possession that makes blessed, but the spirit
out of which life is set in motion, and the fullness of life coming from
within, which just as readily by possession as by loss can be called forth
to unfold itself. Only then can, on the whole, a genuine, life-animating and
forming principle go forth out of love: when it transfigures itself into a
spirituality and purity in which the joy in the beloved melts together with
the joy in the Eternal's glory and fullness present in everyone---with the
joy in the God in whom we all are, live and are moved.

\begin{center}\rule{0.28\linewidth}{0.4pt}\end{center}

\end{document}
